% ==========================================================
%   EST108 — Estadísticas y Data Science
%   Apuntes por Agustín Miranda
%   Universidad Adolfo Ibáñez — Primer semestre 2026
% ==========================================================

\documentclass[12pt, a4paper]{book}

% ── CODIFICACIÓN Y LENGUAJE ───────────────────────────────
\usepackage[utf8]{inputenc}
\usepackage[T1]{fontenc}
\usepackage[spanish, es-noshorthands]{babel}

% ── FUENTES ───────────────────────────────────────────────
\usepackage{lmodern}
\usepackage{microtype}

% ── MATEMÁTICAS ───────────────────────────────────────────
\usepackage{amsmath}
\usepackage{amssymb}
\usepackage{amsthm}
\usepackage{mathtools}

% ── DISEÑO DE PÁGINA ──────────────────────────────────────
\usepackage[top=3cm, bottom=3cm, left=3.2cm, right=3.2cm]{geometry}
\usepackage{fancyhdr}
\usepackage{titlesec}
\usepackage{emptypage}
\usepackage{float}

% ── COLORES Y CAJAS ───────────────────────────────────────
\usepackage[dvipsnames]{xcolor}
\usepackage[most, breakable]{tcolorbox}
\usepackage{colortbl}
\usepackage{multirow}

% ── GRÁFICOS ──────────────────────────────────────────────
\usepackage{tikz}
\usepackage{pgfplots}
\usetikzlibrary{patterns, positioning}
\pgfplotsset{compat=1.18}

% ── HIPERVÍNCULOS ─────────────────────────────────────────
\usepackage{hyperref}
\hypersetup{
    colorlinks = true,
    linkcolor  = NavyBlue,
    urlcolor   = NavyBlue,
    citecolor  = NavyBlue,
    pdftitle   = {EST108 --- Estadisticas y Data Science},
    pdfauthor  = {Agustin Miranda},
}

% ── MISCELÁNEOS ───────────────────────────────────────────
\usepackage{enumitem}
\usepackage{booktabs}
\usepackage{setspace}
\setstretch{1.15}

% ==========================================================
%                   PALETA DE COLORES
% ==========================================================
\definecolor{azul}      {RGB}{0,   51, 102}
\definecolor{azulclaro} {RGB}{230, 240, 255}
\definecolor{rojo}      {RGB}{139,  0,   0}
\definecolor{rojoclaro} {RGB}{255, 235, 235}
\definecolor{verde}     {RGB}{0,   90,  50}
\definecolor{verdeclaro}{RGB}{230, 255, 240}
\definecolor{gris}      {RGB}{245, 245, 245}
\definecolor{naranjo}      {RGB}{153,  76,   0}
\definecolor{naranjaclaro} {RGB}{255, 240, 220}

% ==========================================================
%               ENTORNOS MATEMÁTICOS
% Uso:
%   \begin{definicion}{Título}{etiqueta} ... \end{definicion}
%   \begin{axioma}{Título}{etiqueta}     ... \end{axioma}
%   \begin{teorema}{Título}{etiqueta}    ... \end{teorema}
%   \begin{proposicion}{Título}{etiqueta}... \end{proposicion}
%   \begin{observacion} ...              \end{observacion}
%   \begin{proof} ...                    \end{proof}
% ==========================================================

\newtcbtheorem[number within=chapter]{definicion}{Definici\'{o}n}{
    enhanced, unbreakable,
    colback=azulclaro, colframe=azul,
    fonttitle=\bfseries\color{white}, coltitle=white,
    attach boxed title to top left={yshift=-2mm, xshift=6mm},
    boxed title style={colback=azul, rounded corners},
    separator sign={\ ---\ },
}{def}

\newtcbtheorem[number within=chapter]{axioma}{Axioma}{
    enhanced, unbreakable,
    colback=rojoclaro, colframe=rojo,
    fonttitle=\bfseries\color{white}, coltitle=white,
    attach boxed title to top left={yshift=-2mm, xshift=6mm},
    boxed title style={colback=rojo, rounded corners},
    separator sign={\ ---\ },
}{ax}

\newtcbtheorem[number within=chapter]{teorema}{Teorema}{
    enhanced, unbreakable,
    colback=gris, colframe=azul!70!black,
    fonttitle=\bfseries\color{white}, coltitle=white,
    attach boxed title to top left={yshift=-2mm, xshift=6mm},
    boxed title style={colback=azul!70!black, rounded corners},
    separator sign={\ ---\ },
}{teo}

\newtcbtheorem[number within=chapter]{proposicion}{Proposici\'{o}n}{
    enhanced, unbreakable,
    colback=verdeclaro, colframe=verde,
    fonttitle=\bfseries\color{white}, coltitle=white,
    attach boxed title to top left={yshift=-2mm, xshift=6mm},
    boxed title style={colback=verde, rounded corners},
    separator sign={\ ---\ },
}{prop}

\newtcbtheorem[number within=chapter]{corolario}{Corolario}{
	enhanced, unbreakable,
	colback=green!10, colframe=green!50!black,
	fonttitle=\bfseries\color{white}, coltitle=white,
	attach boxed title to top left={yshift=-2mm, xshift=6mm},
	boxed title style={colback=green!50!black, rounded corners},
	separator sign={\ ---\ },
}{cor}

\newenvironment{observacion}{%
    \par\medskip\noindent\textbf{\color{azul}Observaci\'{o}n.}\enspace\itshape
}{\par\medskip}

\newtcbtheorem[number within=chapter]{ejemplo}{Ejemplo}{
	enhanced, breakable,
	colback=naranjaclaro, colframe=naranjo,
	fonttitle=\bfseries\color{white}, coltitle=white,
	attach boxed title to top left={yshift=-2mm, xshift=6mm},
	boxed title style={colback=naranjo, rounded corners},
	separator sign={\ ---\ },
}{ej}

\renewcommand{\qedsymbol}{$\blacksquare$}
\renewenvironment{proof}{%
    \par\medskip\noindent\textit{Demostraci\'{o}n.}\enspace
}{\qed\par\medskip}

% ==========================================================
%              ENCABEZADOS Y ESTILOS
% ==========================================================
\pagestyle{fancy}
\fancyhf{}
\fancyhead[LE]{\small\textit{EST108 --- Estad\'{\i}sticas y Data Science}}
\fancyhead[RO]{\small\textit{\leftmark}}
\fancyfoot[C]{\thepage}
\renewcommand{\headrulewidth}{0.4pt}

\titleformat{\chapter}[display]
  {\normalfont\huge\bfseries\color{azul}}{Tema \thechapter}{12pt}{\Huge}
\titlespacing*{\chapter}{0pt}{40pt}{30pt}

\titleformat{\section}
  {\normalfont\Large\bfseries\color{azul!80!black}}{\thesection}{1em}{}

\titleformat{\subsection}
  {\normalfont\large\bfseries\color{azul!60!black}}{\thesubsection}{1em}{}

\numberwithin{equation}{chapter}

% ==========================================================
\begin{document}

% ── PORTADA ───────────────────────────────────────────────
\begin{titlepage}
    \centering
    \vspace*{2.5cm}
    {\color{azul}\rule{\linewidth}{1.5pt}}\\[0.5cm]
    {\fontsize{32}{38}\selectfont\bfseries\color{azul}
        Estad\'{\i}sticas y Data Science}\\[0.4cm]
    {\Large EST108 \quad|\quad Ingenier\'{\i}a Comercial}\\[0.2cm]
    {\large Universidad Adolfo Ib\'a\~{n}ez}\\[0.1cm]
    {\color{azul}\rule{\linewidth}{1.5pt}}\\[2cm]
    {\Large Apuntes por \textbf{Agust\'{\i}n Miranda}}\\[0.5cm]
    {\large Primer semestre de 2026}
    \vfill
    {\small Basado en el material de la profesora Florencia Darrigrandi.}
\end{titlepage}

\newpage
\thispagestyle{empty}
\vspace*{\fill}
\begin{center}
	{\large\itshape «La matem\'{a}tica es aquello mediante lo cual
		los hombres gobiernan la naturaleza y a s\'{i} mismos.»}
	
	\bigskip
	{\normalsize --- Andr\'{e}i N. Kolmog\'{o}rov}\\[4pt]
	{\small Matem\'{a}tico sovi\'{e}tico (1903--1987)}
\end{center}
\vspace*{\fill}
\newpage

\tableofcontents
\clearpage

% ==========================================================
\chapter*{Introducci\'{o}n: \textquestiondown Qu\'{e} es la estad\'{\i}stica?}
\addcontentsline{toc}{chapter}{Introducci\'{o}n: \textquestiondown Qu\'{e} es la estad\'{\i}stica?}
\markboth{Introducci\'{o}n}{Introducci\'{o}n}

---El término «estadística» deriva del latín \textit{statisticum collegium} (consejo de Estado), y de la palabra italiana \textit{statista} («hombre de estado», «político» o «estadista»). Pero su uso contemporáneo fue acuñado por el historiador y jurista alemán Gottfried Achenwall en 1749 para referirse al análisis de datos acerca del Estado, significando «ciencia del Estado».

Durante el siglo XIX e inicios del XX, el término se fue transformando en una disciplina matemática propiamente rigurosa bajo el trabajo de Francis Galton, Karl Pearson, Ronald Fisher, entre otros que acuñaron los métodos estadísticos modernos que conocemos hoy en día.

De esta manera, podemos definir a la estadística como la ciencia de recolectar, analizar, presentar e interpretar datos.



% ==========================================================
\chapter{Probabilidad}

\section{\textquestiondown Qu\'{e} es la probabilidad?}

La probabilidad es el lenguaje matemático para cuantificar la incertidumbre.


\section{\textquestiondown Qu\'{e} es un experimento?}

Un experimento es el modelo matemático para cualquier procedimiento que puede ser repetido infinitamente y tiene un conjunto de posibles resultados, que conocemos como:

\section{Espacio muestral y eventos}

\begin{definicion}{Espacio muestral}{espacio_muestral}
	El espacio muestral (denotado usualmente $\Omega$ o $S$) es el conjunto de \textbf{todos los resultados posibles} de un experimento aleatorio.
	
	\medskip
	\noindent\textbf{Definición formal:}
	\[
	\Omega = \{ \omega : \omega \text{ es un resultado posible del experimento} \}
	\]
	
	\medskip
	Cada elemento $\omega \in \Omega$ se llama \textbf{punto muestral} o \textbf{resultado elemental}.
\end{definicion}

Los subconjuntos de $\Omega$ son llamados \textbf{eventos}. Dado que los eventos son subconjuntos de $\Omega$, las operaciones de la teoría de conjuntos ---unión, intersección y complemento--- se aplican directamente, adquiriendo aquí una interpretación probabilística:

\subsection{Operaciones entre eventos}

\begin{center}
	\begin{tabular}{@{}ll@{}}
		$\Omega$              & Espacio muestral \\[4pt]
		$\omega$              & Resultado (punto o elemento) \\[4pt]
		$A$                   & Evento (subconjunto de $\Omega$) \\[4pt]
		$A^c$                 & Complemento de $A$ (no $A$) \\[4pt]
		$A \cup B$            & Unión ($A$ o $B$) \\[4pt]
		$A \cap B$ o $AB$     & Intersección ($A$ y $B$) \\[4pt]
		$A - B$               & Diferencia ($\omega$ en $A$ pero no en $B$) \\[4pt]
		$A \subset B$         & Inclusión ($A$ implica $B$) \\[4pt]
		$\emptyset$           & Evento imposible (siempre falso) \\[4pt]
		$\Omega$              & Evento seguro (siempre verdadero) \\[4pt]
		$A \cap B = \emptyset$& $A$ y $B$ son mutuamente excluyentes \\
	\end{tabular}
\end{center}

\section{Interpretaciones de la probabilidad}

Antes de formalizar la probabilidad, es útil preguntarse qué \textit{significa} que un evento tenga probabilidad 
$p$. Existen tres respuestas clásicas a esta pregunta, cada una con 
supuestos y limitaciones distintas.

\subsection{Interpretación clásica o equiprobable}

\begin{definicion}{Espacio equiprobable}{equiprobable}
	Un espacio muestral $\Omega$ es \textbf{equiprobable} si todos sus 
	elementos tienen igual probabilidad de ocurrir. En este caso, la 
	probabilidad de un evento $A$ se define como:
	\[
	P(A) = \frac{\#A}{\#\Omega} = \frac{\text{casos favorables}}
	{\text{casos totales}}
	\]
	donde $\#A$ denota la cardinalidad del conjunto $A$.
\end{definicion}

Esta interpretación es la más antigua, se remonta a los trabajos del matemático francés 
Pierre-Simon Laplace (1749-1827) sobre juegos de azar, y es la que se aplica naturalmente cuando 
todos los resultados del experimento son igualmente posibles: lanzar una 
moneda justa, tirar un dado equilibrado, extraer una carta de una baraja 
bien mezclada.

\begin{observacion}
	La condición de equiprobabilidad debe \textit{justificarse}, no 
	asumirse. Un error común es definir el espacio muestral de forma 
	que los elementos no sean igualmente probables y luego aplicar la 
	fórmula de Laplace de todos modos. Por ejemplo, si una urna contiene 
	4 fichas con el número 1 y 2 fichas con el número 2, el espacio 
	muestral correcto es $\Omega = \{1_a, 1_b, 1_c, 1_d, 2_a, 2_b\}$ 
	con $\#\Omega = 6$, y no $\Omega = \{1, 2\}$ con $\#\Omega = 2$, 
	pues en este último los elementos \textit{no} son equiprobables.
\end{observacion}

\subsection{Interpretación frecuentista}

La interpretación frecuentista define la probabilidad de un evento $A$ 
como el límite de su frecuencia relativa cuando el experimento se repite 
un número muy grande de veces bajo condiciones idénticas:

\[
P(A) = \lim_{n \to \infty} \frac{n_A}{n}
\]

donde $n$ es el número total de repeticiones y $n_A$ es el número de 
veces que ocurrió $A$.

Esta es la interpretación que subyace a la estadística aplicada: si 
lanzamos una moneda 10.000 veces y obtenemos cara en 5.032 de ellas, 
estimamos $P(\text{cara}) \approx 0.5$. A mayor número de repeticiones, 
mejor la estimación.

\begin{observacion}
	La interpretación frecuentista tiene una limitación fundamental: 
	exige que el experimento sea \textit{repetible} en condiciones 
	idénticas. Esto la hace inaplicable a eventos únicos o irrepetibles 
	— como la probabilidad de que llueva mañana en Santiago, o de que 
	una empresa quiebre el próximo año — para los cuales no existe una 
	frecuencia relativa observable.
\end{observacion}

\subsection{Interpretación subjetiva (bayesiana)}

La interpretación subjetiva define la probabilidad como un \textbf{grado 
	de creencia racional} sobre la ocurrencia de un evento, basado en la 
información disponible. No requiere repetición del experimento ni 
equiprobabilidad de los resultados.

Bajo esta visión, dos agentes con información distinta pueden asignar 
probabilidades distintas al mismo evento sin que ninguno esté 
``equivocado''. Lo que se exige es que las creencias sean 
\textit{coherentes}, es decir, que satisfagan los axiomas de Kolmogórov que presentaremos en la siguiente sección.

\begin{observacion}
	Las tres interpretaciones no son mutuamente excluyentes en la 
	práctica. La interpretación clásica es un caso especial de la 
	frecuentista bajo simetría perfecta; y la frecuentista puede 
	entenderse como un caso especial de la subjetiva cuando la evidencia 
	empírica es abrumadora. Lo que unifica las tres es que, en todos los 
	casos, la función $P$ debe satisfacer los axiomas que veremos en la 
	sección \ref{ax:kolmogorov}.
\end{observacion}

\section{Los axiomas de Kolmogorov}

Los axiomas de probabilidad son las condiciones mínimas que deben verificarse para que una función definida sobre un conjunto de sucesos determine consistentemente sus probabilidades. Fueron formulados por el matemático soviético Andréi Kolmogórov en 1933. Y estos son:

\medskip
Sea $\Omega$ un espacio muestral y $\mathcal{A}$ una $\sigma$-álgebra de subconjuntos de $\Omega$.

\begin{definicion}{$\sigma$-álgebra}{sigma_algebra}
	Sea $\Omega$ un espacio muestral. Una \textbf{$\sigma$-álgebra} 
	$\mathcal{A}$ sobre $\Omega$ es una colección de subconjuntos de 
	$\Omega$ que satisface las siguientes tres propiedades:
	
	\medskip
	\noindent\textbf{S1)} $\Omega \in \mathcal{A}$.
	
	\medskip
	\noindent\textbf{S2)} Si $A \in \mathcal{A}$, entonces 
	$A^c \in \mathcal{A}$ \quad (cerrada bajo complementos).
	
	\medskip
	\noindent\textbf{S3)} Si $A_1, A_2, \ldots \in \mathcal{A}$, 
	entonces $\displaystyle\bigcup_{i=1}^{\infty} A_i \in \mathcal{A}$ 
	\quad (cerrada bajo uniones contables).
	
	\medskip
	A los elementos de $\mathcal{A}$ los llamamos \textbf{eventos}, y 
	a la terna $(\Omega, \mathcal{A})$ se la denomina 
	\textbf{espacio medible}.
\end{definicion}

En la práctica, para espacios muestrales finitos o contables, 
$\mathcal{A}$ corresponde simplemente al conjunto de todos los 
subconjuntos posibles de $\Omega$ — también llamado 
\textit{conjunto potencia} y denotado $2^\Omega$ — por lo que no es 
necesario preocuparse por esta distinción. La $\sigma$-álgebra cobra 
importancia técnica en espacios continuos, donde no todos los 
subconjuntos de $\Omega$ admiten una asignación de probabilidad 
consistente.

A los elementos de $\mathcal{A}$ los llamamos \textbf{eventos} o 
\textbf{sucesos}, y más precisamente \textbf{conjuntos 
	$\mathcal{A}$-medibles}. A la terna $(\Omega, \mathcal{A})$ se la 
denomina \textbf{espacio medible} o \textbf{espacio probabilizable}, 
según el contexto.

\medskip
\noindent\textbf{Corolario de S1 y S2:} $\emptyset \in \mathcal{A}$, 
pues $\emptyset = \Omega^c$ y $\Omega \in \mathcal{A}$ por S1, luego 
S2 garantiza que $\Omega^c = \emptyset \in \mathcal{A}$.

\begin{axioma}{Axiomas de Kolmog\'{o}rov}{kolmogorov}
	Una función $P : \mathcal{A} \longrightarrow \mathbb{R}$ es una \textbf{función de probabilidad} sobre $(\Omega, \mathcal{A})$ si se cumplen los siguientes axiomas:
	
	\medskip
	\noindent\textbf{A1)} La probabilidad de cualquier suceso $S$ es no negativa:
	\[
	P(S) \ge 0 \quad \forall S \in \mathcal{A}.
	\]
	
	\noindent\textbf{A2)} La probabilidad del suceso seguro es igual a uno:
	\[
	P(\Omega) = 1.
	\]
	
	\noindent\textbf{A3)} Si $\{S_n\}_{n \in \mathbb{N}} \subseteq \mathcal{A}$ son sucesos mutuamente excluyentes, es decir, $S_i \cap S_j = \emptyset \quad \forall i \neq j$, entonces:
	\[
	P\left(\bigcup_{n=1}^{\infty} S_n\right) = \sum_{n=1}^{\infty} P(S_n).
	\]
\end{axioma}

A la terna $(\Omega, \mathcal{A}, P)$ se la denomina \textbf{espacio de probabilidad}, esto es, un espacio de sucesos (espacio muestral) en el que se han definido los posibles sucesos a considerar (la $\sigma$-álgebra de sucesos) y la probabilidad de cada suceso (la función de probabilidad).

Una vez establecidos los tres axiomas, podemos nombrar formalmente el 
objeto que hemos construido.

\begin{definicion}{Espacio de probabilidad}{espacio_prob}
	Un \textbf{espacio de probabilidad} es una terna 
	$(\Omega, \mathcal{A}, P)$ donde:
	\begin{itemize}
		\item $\Omega$ es el espacio muestral,
		\item $\mathcal{A}$ es una $\sigma$-álgebra sobre $\Omega$ 
		(véase Definición~\ref{def:sigma_algebra}),
		\item $P : \mathcal{A} \to [0,1]$ es una función de 
		probabilidad que satisface los axiomas A1, A2 y A3.
	\end{itemize}
\end{definicion}

El espacio de probabilidad $(\Omega, \mathcal{A}, P)$ es el objeto 
matemático fundamental sobre el que descansa toda la teoría. Cualquier 
noción probabilística que desarrollemos — variables aleatorias, 
esperanza, independencia — vive dentro de este marco.

\subsection{Propiedades derivadas}

De los axiomas de Kolmogórov se deducen las siguientes propiedades. 
A diferencia de los axiomas, estas no se asumen sino que se 
\textit{demuestran}.

\begin{proposicion}{Probabilidad del complemento}{complemento}
	Para cualquier evento $A \in \mathcal{A}$:
	\[
	P(A^c) = 1 - P(A)
	\]
\end{proposicion}

\begin{proof}
	Como $A$ y $A^c$ son mutuamente excluyentes y $A \cup A^c = \Omega$, 
	por A3:
	\[
	P(A \cup A^c) = P(A) + P(A^c)
	\]
	Por A2, $P(\Omega) = 1$, luego $1 = P(A) + P(A^c)$, de donde:
	\[
	P(A^c) = 1 - P(A) \qedhere
	\]
\end{proof}

\begin{proposicion}{Probabilidad del evento imposible}{imposible}
	\[
	P(\emptyset) = 0
	\]
\end{proposicion}

\begin{proof}
	Basta notar que $\emptyset = \Omega^c$. Aplicando la proposición 
	anterior:
	\[
	P(\emptyset) = P(\Omega^c) = 1 - P(\Omega) = 1 - 1 = 0 \qedhere
	\]
\end{proof}

\begin{proposicion}{Monotonía y diferencia de eventos}{monotonia}
	Si $A \subset B$, entonces:
	\[
	P(A) \leq P(B) \qquad \text{y} \qquad P(B - A) = P(B) - P(A)
	\]
\end{proposicion}

\begin{proof}
	Si $A \subset B$, entonces $B$ puede descomponerse como unión 
	disjunta:
	\[
	B = A \cup (B - A), \qquad A \cap (B-A) = \emptyset
	\]
	Por A3:
	\[
	P(B) = P(A) + P(B - A)
	\]
	Como por A1 $P(B - A) \geq 0$, se sigue que $P(A) \leq P(B)$. 
	Despejando:
	\[
	P(B - A) = P(B) - P(A) \qedhere
	\]
\end{proof}

\begin{proposicion}{Aditividad general (regla de la suma)}{suma}
	Para cualquier par de eventos $A, B \in \mathcal{A}$:
	\[
	P(A \cup B) = P(A) + P(B) - P(A \cap B)
	\]
\end{proposicion}

\begin{proof}
	Descomponemos $A \cup B$ en tres partes mutuamente excluyentes:
	\[
	A \cup B = (A - B) \cup (A \cap B) \cup (B - A)
	\]
	Por A3:
	\[
	P(A \cup B) = P(A - B) + P(A \cap B) + P(B - A)
	\]
	Aplicando la proposición anterior a $A \cap B \subset A$ y 
	$A \cap B \subset B$:
	\[
	P(A - B) = P(A) - P(A \cap B), \qquad P(B - A) = P(B) - 
	P(A \cap B)
	\]
	Sustituyendo:
	\[
	P(A \cup B) = P(A) - P(A \cap B) + P(A \cap B) + P(B) - 
	P(A \cap B)
	\]
	\[
	P(A \cup B) = P(A) + P(B) - P(A \cap B) \qedhere
	\]
\end{proof}

\vspace{0.5em}
\begin{center}
	\begin{tabular}{@{}lll@{}}
		\toprule
		\textbf{Propiedad} & \textbf{Expresión} & \textbf{Condición} \\
		\midrule
		Complemento      & $P(A^c) = 1 - P(A)$                          & --- \\[6pt]
		Evento imposible & $P(\emptyset) = 0$                            & --- \\[6pt]
		Monotonía        & $P(A) \leq P(B)$                             & $A \subset B$ \\[6pt]
		Diferencia       & $P(B - A) = P(B) - P(A)$                    & $A \subset B$ \\[6pt]
		Regla de la suma & $P(A \cup B) = P(A) + P(B) - P(A \cap B)$   & --- \\
		\bottomrule
	\end{tabular}
\end{center}
\vspace{0.5em}

\section{Probabilidad condicional}

Hasta ahora hemos calculado probabilidades sobre el espacio muestral 
completo $\Omega$. Sin embargo, muchas veces disponemos de información 
parcial: sabemos que cierto evento $B$ \textit{ya ocurrió}, y queremos 
actualizar la probabilidad de otro evento $A$ a la luz de ese 
conocimiento.

\subsection*{Motivación: la reducción del espacio muestral}

La idea central es la siguiente. Si sabemos con certeza que $B$ ocurrió, 
entonces los resultados fuera de $B$ dejan de ser posibles. En efecto, 
nuestro universo de referencia se \textbf{reduce} de $\Omega$ a $B$: el 
evento $B$ pasa a funcionar como un nuevo espacio muestral.

Dentro de este universo reducido, la única parte de $A$ que sigue siendo 
relevante es la que cae dentro de $B$, es decir, $A \cap B$. La figura 
\ref{fig:condicional} ilustra esta reducción.

\begin{figure}[h]
	\centering
	\begin{tikzpicture}[scale=0.9]
		% --- Izquierda: sin condicionar ---
		\begin{scope}[shift={(0,0)}]
			\draw[thick] (-2.2,-1.6) rectangle (2.2,1.6);
			\node at (-1.9,1.3) {$\Omega$};
			\begin{scope}
				\clip (-0.6,0) circle (1.1);
				\fill[rojo!30] (0.6,0) circle (1.1);
			\end{scope}
			\draw[thick,fill=azul!12,fill opacity=0.5] (-0.6,0) circle (1.1);
			\draw[thick,fill=azul!12,fill opacity=0.5] (0.6,0) circle (1.1);
			\node at (-1.1,0) {$A$};
			\node at (1.1,0) {$B$};
			\node[font=\scriptsize] at (0,-0.05) {$A\cap B$};
			\node[font=\small] at (0,-2.1) {Sin condicionar};
		\end{scope}
		% --- Derecha: condicionado en B ---
		\begin{scope}[shift={(6,0)}]
			\draw[dashed,gray] (-2.2,-1.6) rectangle (2.2,1.6);
			\node[gray] at (-1.9,1.3) {$\Omega$};
			\draw[thick,fill=rojo!12] (0.6,0) circle (1.1);
			\begin{scope}
				\clip (0.6,0) circle (1.1);
				\fill[rojo!35] (-0.6,0) circle (1.1);
			\end{scope}
			\draw[thick,dotted] (-0.6,0) circle (1.1);
			\node at (1.3,0.7) {$B$};
			\node[font=\scriptsize] at (0.1,-0.05) {$A\cap B$};
			\node[font=\small] at (0,-2.1) {Condicionado en $B$};
		\end{scope}
		\draw[->,thick,azul] (2.6,0) -- (3.4,0);
	\end{tikzpicture}
	\caption{Al condicionar en $B$, el espacio muestral se reduce de 
		$\Omega$ a $B$. Lo único relevante de $A$ pasa a ser $A \cap B$.}
	\label{fig:condicional}
\end{figure}

Ahora bien, para que las probabilidades dentro de este nuevo universo 
$B$ vuelvan a sumar $1$, debemos \textit{re-escalar} dividiendo por 
$P(B)$. Esto motiva la siguiente definición.

\subsection{Definición}

\begin{definicion}{Probabilidad condicional}{condicional}
	Sean $A$ y $B$ dos eventos con $P(B) \neq 0$. La \textbf{probabilidad 
		condicional de $A$ dado $B$} se define como:
	\[
	P(A \mid B) = \frac{P(A \cap B)}{P(B)}
	\]
\end{definicion}

La condición $P(B) \neq 0$ es esencial: no tiene sentido condicionar en 
un evento que no puede ocurrir, pues el universo reducido sería vacío.

\begin{observacion}
	La probabilidad condicional $P(\,\cdot \mid B)$ es, ella misma, una 
	función de probabilidad legítima sobre el espacio reducido: satisface 
	los tres axiomas de Kolmogórov. En particular, $P(B \mid B) = 1$, 
	lo que confirma que $B$ actúa como el nuevo evento seguro.
\end{observacion}

\begin{proof}
	Verificamos los tres axiomas para $P(\,\cdot \mid B)$:
	
	\medskip
	\noindent\textbf{A1 — No negatividad:} Sea $A \in \mathcal{A}$. 
	Como $A \cap B \subseteq B$ y $P$ satisface A1, tenemos 
	$P(A \cap B) \geq 0$. Como además $P(B) > 0$ por hipótesis:
	\[
	P(A \mid B) = \frac{P(A \cap B)}{P(B)} \geq 0
	\]
	
	\medskip
	\noindent\textbf{A2 — Normalización:} 
	\[
	P(\Omega \mid B) = \frac{P(\Omega \cap B)}{P(B)} = 
	\frac{P(B)}{P(B)} = 1
	\]
	donde usamos que $\Omega \cap B = B$.
	
	\medskip
	\noindent\textbf{A3 — Aditividad contable:} Sean 
	$A_1, A_2, \ldots \in \mathcal{A}$ eventos mutuamente excluyentes. 
	Entonces los eventos $A_1 \cap B,\ A_2 \cap B, \ldots$ también son 
	mutuamente excluyentes, pues:
	\[
	(A_i \cap B) \cap (A_j \cap B) = (A_i \cap A_j) \cap B = 
	\emptyset \cap B = \emptyset \quad \forall\, i \neq j
	\]
	Aplicando A3 de Kolmogórov a la función $P$ original:
	\[
	P\!\left(\bigcup_{n=1}^{\infty} A_n \ \middle|\ B\right) 
	= \frac{P\!\left(\left(\bigcup_{n=1}^{\infty} A_n\right) 
		\cap B\right)}{P(B)}
	= \frac{P\!\left(\bigcup_{n=1}^{\infty} (A_n \cap B)\right)}
	{P(B)}
	= \frac{\sum_{n=1}^{\infty} P(A_n \cap B)}{P(B)}
	= \sum_{n=1}^{\infty} P(A_n \mid B)
	\]
	donde usamos la distributividad de la intersección sobre la unión. 
	$\qedhere$
\end{proof}

\subsection{Regla de la multiplicaci\'{o}n}

Despejando $P(A \cap B)$ de la definición anterior obtenemos una 
identidad muy útil para calcular probabilidades de intersecciones:

\begin{proposicion}{Regla de la multiplicación}{multiplicacion}
	Sean $A$ y $B$ dos eventos con $P(B) \neq 0$. Entonces:
	\[
	P(A \cap B) = P(A \mid B) \cdot P(B)
	\]
\end{proposicion}

\begin{proof}
	Despejando $P(A \cap B)$ directamente de la 
	Definición~\ref{def:condicional}:
	\[
	P(A \mid B) = \frac{P(A \cap B)}{P(B)} \implies 
	P(A \cap B) = P(A \mid B) \cdot P(B) \qedhere
	\]
\end{proof}

Esta regla es especialmente útil en experimentos secuenciales ---por 
ejemplo, extracciones sucesivas sin reposición--- donde la probabilidad 
de cada etapa depende del resultado de la anterior.

\section{Independencia de eventos}

Al estudiar la probabilidad condicional, vimos que $P(A \mid B)$ 
cuantifica cómo el conocimiento de $B$ modifica la probabilidad de $A$. 
Surge entonces una pregunta natural: ¿qué ocurre cuando ese conocimiento 
no modifica nada? Es decir, ¿cuándo saber que $B$ ocurrió es 
completamente irrelevante para $A$?

\begin{definicion}{Independencia de eventos}{independencia}
	Sean $A$ y $B$ dos eventos con $P(B) \neq 0$. Decimos que $A$ y $B$ 
	son \textbf{independientes} si:
	\[
	P(A \mid B) = P(A)
	\]
	Es decir, el conocimiento de que $B$ ocurrió no altera la 
	probabilidad de $A$.
\end{definicion}

\begin{observacion}
	Esta definición exige $P(B) \neq 0$ para que $P(A \mid B)$ esté 
	definida. La proposición siguiente elimina esa restricción y 
	constituye el criterio operacional que se usa en la práctica.
\end{observacion}

\begin{proposicion}{Criterio de independencia}{criterio_independencia}
	Sean $A$ y $B$ dos eventos. $A$ y $B$ son independientes si y solo 
	si:
	\[
	P(A \cap B) = P(A) \cdot P(B)
	\]
\end{proposicion}

\begin{proof}
	Partimos de la definición de independencia $P(A \mid B) = P(A)$ y 
	aplicamos la regla de la multiplicación 
	(Proposición~\ref{prop:multiplicacion}):
	\[
	P(A \cap B) = P(A \mid B) \cdot P(B) = P(A) \cdot P(B) \qedhere
	\]
\end{proof}

Esta caracterización es más general que la definición porque no requiere 
$P(B) \neq 0$, y es simétrica en $A$ y $B$ — lo que hace evidente que 
si $A$ es independiente de $B$, entonces $B$ es independiente de $A$.

\begin{observacion}
	\textbf{Independencia vs.\ exclusión mutua.} Estos conceptos son 
	frecuentemente confundidos, pero son casi opuestos. Si $A$ y $B$ 
	son mutuamente excluyentes ($A \cap B = \emptyset$) y ambos tienen 
	probabilidad positiva, entonces:
	\[
	P(A \cap B) = 0 \neq P(A) \cdot P(B)
	\]
	luego \textit{no} son independientes. De hecho, la exclusión mutua 
	implica la máxima dependencia posible: saber que $A$ ocurrió hace 
	imposible que ocurra $B$.
\end{observacion}

\begin{proposicion}{Independencia de complementos}{ind_complementos}
	Si $A$ y $B$ son independientes, entonces $A^c$ y $B^c$ también 
	son independientes.
\end{proposicion}

\begin{proof}
	Por la Proposición~\ref{prop:criterio_independencia}, basta mostrar 
	que $P(A^c \cap B^c) = P(A^c) \cdot P(B^c)$.
	
	\medskip
	Observamos primero que $A^c \cap B^c = (A \cup B)^c$ por las leyes 
	de De Morgan. Luego, usando la Proposición~\ref{prop:complemento} 
	y la regla de la suma (Proposición~\ref{prop:suma}):
	\[
	P(A^c \cap B^c) = P((A \cup B)^c) = 1 - P(A \cup B)
	\]
	\[
	= 1 - P(A) - P(B) + P(A \cap B)
	\]
	Como $A$ y $B$ son independientes, $P(A \cap B) = P(A) \cdot P(B)$:
	\[
	= 1 - P(A) - P(B) + P(A) \cdot P(B)
	\]
	Factorizando:
	\[
	= (1 - P(A))(1 - P(B)) = P(A^c) \cdot P(B^c) \qedhere
	\]
\end{proof}

\section{Probabilidad total y Teorema de Bayes}

\subsection{Partición del espacio muestral}

\begin{definicion}{Partición}{particion}
	Una colección de eventos $A_1, A_2, \ldots, A_k$ constituye una 
	\textbf{partición} del espacio muestral $\Omega$ si satisface:
	
	\medskip
	\noindent\textbf{P1)} Son \textbf{mutuamente excluyentes}: 
	$A_i \cap A_j = \emptyset$ para todo $i \neq j$.
	
	\medskip
	\noindent\textbf{P2)} Son \textbf{exhaustivos}: 
	$\displaystyle\bigcup_{i=1}^{k} A_i = \Omega$.
\end{definicion}

Intuitivamente, una partición divide a $\Omega$ en piezas que no se 
solapan y que, en conjunto, lo cubren por completo. Exactamente uno de 
los eventos $A_i$ ocurre siempre.

\subsection{Teorema de probabilidades totales}

A veces no es posible calcular $P(B)$ de manera directa, pero sí 
conocemos cómo se comporta $B$ dentro de cada pieza de una partición. 
El teorema de probabilidades totales permite recomponer $P(B)$ a partir 
de esa información parcial.

La idea geométrica se aprecia en la figura \ref{fig:totales}: la 
partición $\{A_1, A_2, A_3\}$ divide a $\Omega$ en franjas, y el evento 
$B$ queda automáticamente descompuesto en las piezas disjuntas 
$B \cap A_1$, $B \cap A_2$ y $B \cap A_3$.

\begin{figure}[h]
	\centering
	\begin{tikzpicture}[scale=1.1]
		% Omega
		\draw[thick] (0,0) rectangle (9,4.5);
		\node at (0.35,4.15) {$\Omega$};
		% Franjas de fondo (particion)
		\fill[azul!8]  (0,0) rectangle (3,4.5);
		\fill[rojo!8]  (3,0) rectangle (6,4.5);
		\fill[verde!8] (6,0) rectangle (9,4.5);
		% Intersecciones B cap Ai (elipse centrada en 4.5, radio 3.8 cruza las tres franjas)
		\begin{scope}
			\clip (0,0) rectangle (3,4.5);
			\fill[azul!35] (4.5,2.6) ellipse (3.8 and 1.2);
		\end{scope}
		\begin{scope}
			\clip (3,0) rectangle (6,4.5);
			\fill[rojo!35] (4.5,2.6) ellipse (3.8 and 1.2);
		\end{scope}
		\begin{scope}
			\clip (6,0) rectangle (9,4.5);
			\fill[verde!35] (4.5,2.6) ellipse (3.8 and 1.2);
		\end{scope}
		% Contorno de B
		\draw[thick] (4.5,2.6) ellipse (3.8 and 1.2);
		% Lineas divisorias
		\draw[thick] (3,0) -- (3,4.5);
		\draw[thick] (6,0) -- (6,4.5);
		% Borde exterior
		\draw[thick] (0,0) rectangle (9,4.5);
		% Etiquetas Ai
		\node at (1.5,0.4) {$A_1$};
		\node at (4.5,0.4) {$A_2$};
		\node at (7.5,0.4) {$A_3$};
		% Etiqueta B
		\node[font=\small] at (4.5,4.1) {$B$};
		% Etiquetas intersecciones
		\node[font=\small] at (1.5,2.6) {$B\cap A_1$};
		\node[font=\small] at (4.5,2.6) {$B\cap A_2$};
		\node[font=\small] at (7.5,2.6) {$B\cap A_3$};
	\end{tikzpicture}
	\caption{La partición $\{A_1, A_2, A_3\}$ descompone al evento $B$ 
		en las piezas disjuntas $B \cap A_i$.}
	\label{fig:totales}
\end{figure}

\begin{teorema}{Probabilidades totales}{totales}
	Sea $A_1, A_2, \ldots, A_k$ una partición de $\Omega$ con 
	$P(A_i) > 0$ para todo $i$. Entonces, para cualquier evento $B$:
	\[
	P(B) = \sum_{i=1}^{k} P(B \mid A_i)\, P(A_i)
	\]
\end{teorema}

\begin{proof}
	La demostración consta de tres pasos.
	
	\medskip
	\noindent\textbf{Paso 1 — Descomposición de $B$.} Como los $A_i$ 
	son exhaustivos, $B = B \cap \Omega = B \cap \left(\bigcup_{i=1}^k 
	A_i\right)$. Por distributividad:
	\[
	B = \bigcup_{i=1}^{k} (B \cap A_i)
	\]
	
	\medskip
	\noindent\textbf{Paso 2 — Aditividad.} Como los $A_i$ son 
	mutuamente excluyentes, también lo son los eventos $B \cap A_i$:
	\[
	(B \cap A_i) \cap (B \cap A_j) = B \cap (A_i \cap A_j) = 
	B \cap \emptyset = \emptyset \quad \forall\, i \neq j
	\]
	Aplicando el axioma A3:
	\[
	P(B) = \sum_{i=1}^{k} P(B \cap A_i)
	\]
	
	\medskip
	\noindent\textbf{Paso 3 — Regla de la multiplicación.} Aplicando la 
	Proposición~\ref{prop:multiplicacion} a cada término:
	\[
	P(B \cap A_i) = P(B \mid A_i)\, P(A_i)
	\]
	Sustituyendo:
	\[
	P(B) = \sum_{i=1}^{k} P(B \mid A_i)\, P(A_i) \qedhere
	\]
\end{proof}

\begin{observacion}
	El teorema puede leerse como un \textbf{promedio ponderado}: la 
	probabilidad total de $B$ es el promedio de las probabilidades 
	condicionales $P(B \mid A_i)$, donde cada una se pondera por el 
	``peso'' $P(A_i)$ de su pieza en la partición.
\end{observacion}

\subsection{Teorema de Bayes}

El teorema de probabilidades totales responde a la pregunta: ¿cuál es 
la probabilidad de $B$? El teorema de Bayes responde la pregunta 
\textit{inversa}: si observamos que $B$ ocurrió, ¿cuál es la 
probabilidad de que la pieza responsable haya sido $A$?

\subsubsection*{Versión simple: dos eventos}

Comencemos con el caso más sencillo, en que la partición consta de un 
evento $A$ y su complemento $A^c$.

\begin{teorema}{Regla de Bayes (forma simple)}{bayes_simple}
	Sean $A$ y $B$ dos eventos con $P(A) > 0$ y $P(B) > 0$. Entonces:
	\[
	P(A \mid B) = \frac{P(B \mid A)\, P(A)}{P(B)}
	\]
\end{teorema}

\begin{proof}
	Por la definición de probabilidad condicional aplicada a 
	$P(A \mid B)$, y luego la regla de la multiplicación en el numerador:
	\[
	P(A \mid B) = \frac{P(A \cap B)}{P(B)} = 
	\frac{P(B \mid A)\, P(A)}{P(B)} \qedhere
	\]
\end{proof}

Esta forma es la más fácil de memorizar y expone la idea esencial: el 
teorema \textbf{invierte} la condicional, permitiéndonos pasar de 
$P(B \mid A)$ a $P(A \mid B)$.

\subsubsection*{Versión general: una partición}

Cuando disponemos de una partición completa $A_1, \ldots, A_k$ y no 
conocemos $P(B)$ directamente, podemos expandir el denominador mediante 
el teorema de probabilidades totales. Esto da la forma general.

\begin{teorema}{Regla de Bayes (forma general)}{bayes}
	Sea $A_1, A_2, \ldots, A_k$ una partición de $\Omega$ con 
	$P(A_i) > 0$ para todo $i$, y sea $B$ un evento con $P(B) > 0$. 
	Entonces, para cada $j$:
	\[
	P(A_j \mid B) = \frac{P(B \mid A_j)\, P(A_j)}
	{\displaystyle\sum_{i=1}^{k} P(B \mid A_i)\, P(A_i)}
	\]
\end{teorema}

\begin{proof}
	Partimos de la forma simple aplicada a $A_j$:
	\[
	P(A_j \mid B) = \frac{P(B \mid A_j)\, P(A_j)}{P(B)}
	\]
	y reemplazamos el denominador $P(B)$ por su expresión según el 
	Teorema~\ref{teo:totales} de probabilidades totales:
	\[
	P(A_j \mid B) = \frac{P(B \mid A_j)\, P(A_j)}
	{\sum_{i=1}^{k} P(B \mid A_i)\, P(A_i)} \qedhere
	\]
\end{proof}

\begin{observacion}
	La forma simple es el caso particular de la general con $k = 2$, 
	tomando $A_1 = A$ y $A_2 = A^c$. En efecto, el denominador se 
	convierte en $P(B \mid A)\,P(A) + P(B \mid A^c)\,P(A^c) = P(B)$.
\end{observacion}

\subsubsection*{Interpretación geométrica y un ejemplo concreto}

La forma simple de Bayes tiene una lectura geométrica muy intuitiva, que 
además revela \textit{por qué} la fórmula es válida. La clave está en que 
la probabilidad de la intersección $P(A \cap B)$ puede calcularse de dos 
maneras distintas, según el orden en que descompongamos:
\[
P(A \cap B) = P(B \mid A)\,P(A) = P(A \mid B)\,P(B)
\]
Igualando las dos expresiones y despejando $P(A \mid B)$ se obtiene 
directamente el teorema de Bayes. Toda la fuerza del resultado descansa 
en esta simple identidad.

Consideremos un ejemplo concreto. Supongamos un grupo de 12 personajes, 
cada uno de los cuales puede ser un asesino ($A$) o no ($A^c$), y puede 
comportarse de manera sospechosa ($B$) o no ($B^c$). La siguiente tabla 
de contingencia resume cuántos hay de cada tipo:

\begin{figure}[h]
	\centering
	\begin{tikzpicture}[scale=1.0]
		\fill[verde!30] (0,2) rectangle (2,3);
		\fill[Goldenrod!35] (2,2) rectangle (4,3);
		\fill[azul!22] (0,1) rectangle (2,2);
		\fill[gray!18] (2,1) rectangle (4,2);
		\draw[thick] (0,0) rectangle (5,4);
		\draw[thick] (0,1) -- (5,1);
		\draw[thick] (0,3) -- (5,3);
		\draw[thick] (2,0) -- (2,4);
		\draw[thick] (4,0) -- (4,4);
		\node[font=\small] at (1,3.5) {$B$};
		\node[font=\small] at (3,3.5) {$B^c$};
		\node[font=\small] at (4.5,3.5) {suma};
		\node[font=\small] at (-0.6,2.5) {$A$};
		\node[font=\small] at (-0.6,1.5) {$A^c$};
		\node[font=\small] at (-0.6,0.5) {suma};
		\node at (1,2.5) {\textbf{3}};
		\node at (3,2.5) {\textbf{1}};
		\node at (1,1.5) {\textbf{2}};
		\node at (3,1.5) {\textbf{6}};
		\node at (4.5,2.5) {4};
		\node at (4.5,1.5) {8};
		\node at (1,0.5) {5};
		\node at (3,0.5) {7};
		\node at (4.5,0.5) {12};
	\end{tikzpicture}
	\caption{Tabla de contingencia: hay 3 asesinos sospechosos, 1 asesino no 
		sospechoso, 2 no asesinos sospechosos y 6 no asesinos no sospechosos.}
	\label{fig:contingencia}
\end{figure}

De la tabla leemos directamente:
\[
P(A) = \tfrac{4}{12}, \quad P(B) = \tfrac{5}{12}, \quad 
P(B \mid A) = \tfrac{3}{4}, \quad P(A \mid B) = \tfrac{3}{5}, \quad 
P(A \cap B) = \tfrac{3}{12}
\]

Ahora bien, la identidad geométrica se aprecia en la 
figura~\ref{fig:areas}: el área del rectángulo $A \cap B$ es la misma sin 
importar cómo la construyamos. A la izquierda la obtenemos tomando 
primero la fracción $P(A)$ y luego, dentro de ella, la fracción 
$P(B \mid A)$. A la derecha tomamos primero $P(B)$ y luego $P(A \mid B)$. 
Ambos caminos llegan al mismo rectángulo.

\begin{figure}[h]
	\centering
	\begin{tikzpicture}[scale=1.0]
		\begin{scope}[shift={(0,0)}]
			\draw[thick] (0,0) rectangle (4,4);
			\node[font=\small] at (2,4.4) {Primero $A$, luego $B\mid A$};
			\fill[Goldenrod!25] (0,0) rectangle (1.333,4);
			\fill[verde!35] (0,1) rectangle (1.333,4);
			\draw[thick] (1.333,0) -- (1.333,4);
			\draw[thick] (0,1) -- (1.333,1);
			\node[font=\scriptsize,rotate=90] at (-0.3,2) {$P(A)=\tfrac{4}{12}$};
			\node[font=\scriptsize] at (0.67,2.5) {$A\cap B$};
			\node[font=\scriptsize] at (2.6,3.5) {$P(B\mid A)=\tfrac{3}{4}$};
		\end{scope}
		\begin{scope}[shift={(6,0)}]
			\draw[thick] (0,0) rectangle (4,4);
			\node[font=\small] at (2,4.4) {Primero $B$, luego $A\mid B$};
			\fill[azul!18] (0,0) rectangle (1.667,4);
			\fill[verde!35] (0,1.6) rectangle (1.667,4);
			\draw[thick] (1.667,0) -- (1.667,4);
			\draw[thick] (0,1.6) -- (1.667,1.6);
			\node[font=\scriptsize,rotate=90] at (-0.3,2) {$P(B)=\tfrac{5}{12}$};
			\node[font=\scriptsize] at (0.83,2.8) {$A\cap B$};
			\node[font=\scriptsize] at (2.9,3.5) {$P(A\mid B)=\tfrac{3}{5}$};
		\end{scope}
		\node at (5,2) {$=$};
	\end{tikzpicture}
	\caption{El área de $A \cap B$ se construye de dos formas: 
		$P(A)\,P(B\mid A)$ (izquierda) y $P(B)\,P(A\mid B)$ (derecha). Como ambas 
		dan la misma área, igualarlas produce el teorema de Bayes.}
	\label{fig:areas}
\end{figure}

Verifiquémoslo numéricamente. Por la izquierda:
\[
P(B \mid A)\,P(A) = \tfrac{3}{4} \cdot \tfrac{4}{12} = \tfrac{3}{12}
\]
Por la derecha:
\[
P(A \mid B)\,P(B) = \tfrac{3}{5} \cdot \tfrac{5}{12} = \tfrac{3}{12}
\]
Ambas coinciden, como debía ser. Y despejando $P(A \mid B)$ de la 
igualdad recuperamos Bayes:
\[
P(A \mid B) = \frac{P(B \mid A)\,P(A)}{P(B)} = 
\frac{\tfrac{3}{4} \cdot \tfrac{4}{12}}{\tfrac{5}{12}} = 
\frac{3}{5}
\]
que es exactamente el valor que habíamos leído directamente de la tabla.

\begin{observacion}
	El teorema de Bayes es la herramienta formal que permite 
	\textbf{invertir} condicionales: pasar de $P(B \mid A_j)$ —lo que 
	típicamente conocemos— a $P(A_j \mid B)$ —lo que deseamos inferir. 
	En el lenguaje bayesiano, $P(A_j)$ es la \textbf{probabilidad a 
		priori} (lo que sabemos antes de observar) y $P(A_j \mid B)$ es la 
	\textbf{probabilidad a posteriori} (lo que sabemos después de 
	observar la evidencia $B$).
\end{observacion}

El teorema de Bayes nació de una pregunta filosófica, no técnica. En su 
\textit{Investigación sobre el entendimiento humano}, Hume había 
planteado que la inducción carece de justificación lógica (que consiste en observar patrones o casos específicos en el mundo real para llegar a una conclusión general): por más 
veces que hayamos observado que el sol sale cada mañana, ninguna 
cantidad de observaciones pasadas garantiza lógicamente que saldrá 
mañana. Los fundamentos de esto se deben a la visión empirista radical de este filósofo, pues para Hume la única fuente del conocimiento debe ser la experiencia sensible. Si esto es así, el conocimiento empírico —incluida toda la 
ciencia newtoniana— descansaría sobre una base irracional.

Thomas Bayes, ministro presbiteriano inglés y contemporáneo de Hume, 
buscaba una respuesta rigurosa a ese problema. La encontró escondida 
dentro de la definición de probabilidad condicional que ya existía: su 
fórmula revela que conocer $P(B \mid A)$ es suficiente para inferir 
$P(A \mid B)$ —pasar del efecto a la causa, de la evidencia observable 
a la hipótesis no observable. No se necesita certeza lógica para 
aprender racionalmente; basta un mecanismo coherente de actualización 
de creencias, y ese mecanismo es la fórmula de Bayes.

A más evidencia acumulada, más convergen las probabilidades a 
posteriori hacia la verdad, independientemente de cuál fuera la 
creencia inicial. Es una respuesta pragmática y matemáticamente precisa 
al escepticismo de Hume. Por eso el teorema —publicado póstumamente en 
1763 por su amigo Richard Price— aparece hoy en dominios tan distintos 
como el diagnóstico médico, los filtros de spam, la navegación autónoma 
y el aprendizaje automático: en todos ellos subyace el mismo problema 
que preocupaba a Hume y a Bayes, esto es, dado lo que observo, ¿qué 
puedo inferir sobre lo que no veo?

% ==========================================================
\chapter{Variable aleatoria discreta}

\section{¿Qué es una variable aleatoria?}

Hasta ahora hemos trabajado con eventos — subconjuntos abstractos del 
espacio muestral $\Omega$. Sin embargo, en la práctica queremos asociar 
\textit{números} a los resultados de un experimento aleatorio: el número 
de clientes que llegan en una hora, el monto de una transacción, la 
cantidad de fallas en un proceso. Para eso necesitamos una función que 
traduzca los elementos de $\Omega$ a la recta numérica. Ese objeto es 
la \textbf{variable aleatoria}.

\begin{definicion}{Variable aleatoria}{va}
	Sea $(\Omega, \mathcal{A}, P)$ un espacio de probabilidad. Una 
	\textbf{variable aleatoria} es una función:
	\[
	X : \Omega \longrightarrow \mathbb{R}
	\]
	que satisface la condición de \textbf{medibilidad}: para todo 
	$x \in \mathbb{R}$,
	\[
	\{\omega \in \Omega : X(\omega) \leq x\} \in \mathcal{A}
	\]
	Es decir, el conjunto de resultados para los cuales $X$ no supera 
	el valor $x$ siempre es un evento — pertenece a la $\sigma$-álgebra 
	$\mathcal{A}$.
\end{definicion}

\begin{observacion}
	La condición de medibilidad exige que 
	$\{\omega \in \Omega : X(\omega) \leq x\} \in \mathcal{A}$ para 
	todo $x \in \mathbb{R}$. Esto garantiza que siempre tenga sentido 
	calcular $P(X \leq x)$ — si ese conjunto no perteneciera a 
	$\mathcal{A}$, la función $P$ no estaría definida sobre él. En 
	espacios finitos o contables, toda función $X: \Omega \to \mathbb{R}$ 
	es automáticamente medible, por lo que la condición solo cobra 
	relevancia técnica en espacios continuos.
\end{observacion}

\begin{observacion}
	Usaremos letras mayúsculas ($X$, $Y$, $Z$) para denotar variables 
	aleatorias y letras minúsculas ($x$, $y$, $z$) para los valores 
	concretos que estas pueden tomar. La notación $P(X = x)$ significa 
	la probabilidad del evento $\{\omega \in \Omega : X(\omega) = x\}$.
\end{observacion}

\subsection*{Ejemplo: suma de dos dados}

Lanzamos dos dados equilibrados. El espacio muestral es:
\[
\Omega = \{(i,j) : i,j \in \{1,2,3,4,5,6\}\}, \quad \#\Omega = 36
\]
Definimos la variable aleatoria $X$ como la \textit{suma} de los 
valores obtenidos:
\[
X(i,j) = i + j
\]
Por ejemplo, $X(1,2) = 3$, $X(3,3) = 6$ y $X(6,6) = 12$. La figura 
\ref{fig:va_funcion} ilustra cómo $X$ mapea elementos de $\Omega$ a 
valores en $\mathbb{R}$.

\begin{figure}[h]
	\centering
	\begin{tikzpicture}[scale=1.1]
		\draw[thick,rounded corners=8pt] (-2.5,-1.8) rectangle (2.5,1.8);
		\node at (-2.1,1.5) {$\Omega$};
		\foreach \x/\y/\lab in {
			-1.5/1.1/{(1,1)},
			0.2/1.2/{(1,2)},
			1.7/0.9/{(2,1)},
			-1.8/0.1/{(2,3)},
			0.0/0.2/{(3,3)},
			1.6/0.1/{(4,2)},
			-1.5/-0.9/{(5,6)},
			0.1/-0.8/{(6,5)},
			1.7/-1.0/{(6,6)}
		}{
			\node[font=\scriptsize] at (\x,\y) {\lab};
		}
		\draw[->,thick,azul] (2.7,0) -- (4.3,0);
		\node[font=\small,azul] at (3.5,0.35) {$X$};
		\draw[thick] (4.5,-1.8) -- (4.5,1.8);
		\node at (4.5,2.1) {$\mathbb{R}$};
		\foreach \y/\lab in {
			1.4/2, 0.9/3, 0.2/6, -0.4/9, -1.0/11, -1.6/12
		}{
			\fill[rojo!70] (4.5,\y) circle (2pt);
			\node[font=\scriptsize,right] at (4.6,\y) {\lab};
		}
		\draw[->,gray,thin] (0.6,1.2)  to[out=10,in=150] (4.4,0.9);
		\draw[->,gray,thin] (2.0,0.9)  to[out=10,in=150] (4.4,0.9);
		\draw[->,gray,thin] (0.3,0.2)  to[out=10,in=170] (4.4,0.2);
		\draw[->,gray,thin] (2.0,0.1)  to[out=10,in=170] (4.4,0.2);
		\draw[->,gray,thin] (2.0,-1.0) to[out=10,in=190] (4.4,-1.6);
	\end{tikzpicture}
	\caption{La variable aleatoria $X$ como función $\Omega \to \mathbb{R}$.
		Distintos elementos de $\Omega$ pueden mapearse al mismo valor real
		(por ejemplo, $(1,2)$ y $(2,1)$ ambos mapean a $X = 3$).}
	\label{fig:va_funcion}
\end{figure}

Nótese un hecho importante: varios elementos de $\Omega$ mapean al 
mismo valor de $X$. Los pares $(1,2)$ y $(2,1)$ producen ambos $X=3$; 
los pares $(1,5)$, $(2,4)$, $(3,3)$, $(4,2)$ y $(5,1)$ producen todos 
$X=6$. Esto es lo esperable — la variable aleatoria \textit{resume} 
información del resultado, colapsando múltiples resultados en un mismo 
número.

\section{Clasificación de las variables aleatorias}

\begin{definicion}{Variable aleatoria discreta}{va_discreta}
	Una variable aleatoria $X$ es \textbf{discreta} si su recorrido 
	$\text{Rec}(X) = \{x \in \mathbb{R} : P(X = x) > 0\}$ es finito o 
	contable. Es decir, sus valores posibles pueden listarse como una 
	sucesión $x_1, x_2, x_3, \ldots$
\end{definicion}

\begin{definicion}{Variable aleatoria continua}{va_continua}
	Una variable aleatoria $X$ es \textbf{continua} si su recorrido 
	abarca un intervalo (o unión de intervalos) de la recta numérica. 
	En este caso, $P(X = x) = 0$ para todo valor particular $x$, y la 
	probabilidad solo tiene sentido sobre intervalos.
\end{definicion}

\begin{observacion}
	La distinción entre discreta y continua no es solo técnica — 
	implica herramientas matemáticas distintas. Para variables discretas 
	usaremos sumas; para variables continuas, integrales. En este 
	capítulo nos concentramos en el caso discreto. El caso continuo se 
	trata en el Tema~3.
\end{observacion}

En el ejemplo de los dos dados, $X = i+j$ toma valores en 
$\{2, 3, 4, \ldots, 12\}$ — un conjunto finito — por lo que es una 
variable aleatoria \textbf{discreta}. En cambio, el tiempo de espera 
hasta que llega el próximo cliente puede tomar cualquier valor en 
$[0, +\infty)$, por lo que sería \textbf{continua}.


\section{Funci\'{o}n de masa de probabilidad}

Una vez definida una variable aleatoria discreta $X$, la pregunta 
natural es: ¿cómo se distribuye la probabilidad sobre sus valores 
posibles? La función de masa de probabilidad responde 
exactamente eso — asigna a cada valor posible de $X$ la probabilidad 
de que $X$ tome ese valor.

\begin{definicion}{Función de masa de probabilidad}{fmp}
	Sea $X$ una variable aleatoria discreta con recorrido 
	$\text{Rec}(X)$. La \textbf{función de masa de probabilidad} 
	 de $X$ es 
	la función $p_X : \mathbb{R} \to [0,1]$ definida por:
	\[
	p_X(x) = P(X = x) = P(\{\omega \in \Omega : X(\omega) = x\})
	\]
	para todo $x \in \mathbb{R}$, donde $p_X(x) = 0$ para todo 
	$x \notin \text{Rec}(X)$.
\end{definicion}

El nombre \textbf{función de masa} no es arbitrario — proviene de la 
mecánica clásica. Si imaginamos los valores del recorrido $x \in 
\text{Rec}(X)$ como posiciones sobre una recta, y las probabilidades 
$p_X(x)$ como \textbf{masas puntuales} concentradas en esas posiciones, 
entonces la distribución de probabilidad discreta es formalmente 
idéntica a un sistema físico de partículas. Esta analogía no es solo 
pedagógica — es una identidad matemática exacta que veremos reflejada 
en las definiciones de valor esperado y varianza.

\begin{observacion}
	\textbf{Sobre la terminología.} En el curso, y en muchos textos 
	de estadística aplicada, esta función se denomina 
	\textit{función de distribución de probabilidad}. Esta terminología 
	es técnicamente imprecisa: el término \textit{distribución} está 
	reservado más apropiadamente para la función de distribución 
	\textit{acumulada} $F_X(x) = P(X \leq x)$, que veremos en la 
	sección siguiente. La denominación correcta para $p_X(x) = P(X=x)$ 
	es \textbf{función de masa de probabilidad}, que es la que usaremos 
	en este apunte. La distinción se vuelve especialmente importante en 
	el Tema~3, cuando introduzcamos la \textbf{función de densidad de 
		probabilidad} para variables continuas — un objeto distinto que 
	también suele llamarse coloquialmente "función de distribución", 
	generando confusión.
\end{observacion}

\begin{observacion}
	\textbf{Sobre la notación.} En este apunte usamos $p_X(x)$ para 
	la función de masa de probabilidad, distinguiéndola de $P$ — la 
	función de probabilidad sobre eventos definida en los axiomas de 
	Kolmogórov (Definición~\ref{def:espacio_prob}). La distinción es:
	
	\medskip
	\begin{center}
		\begin{tabular}{@{}lll@{}}
			$P(A)$ & probabilidad de un evento $A \in \mathcal{A}$ & 
			(Kolmogórov) \\[4pt]
			$p_X(x)$ & probabilidad de que $X$ tome el valor $x$ & 
			(función de masa) \\
		\end{tabular}
	\end{center}
	
	\medskip
	La relación entre ambas es: $p_X(x) = P(\{\omega : X(\omega) = x\})$
	— la función de masa es simplemente $P$ aplicada al evento 
	$\{X = x\}$.
	
	\medskip
	En el curso y en muchos textos aplicados esta función también 
	aparece denotada $P(x)$, $p(x)$ o $f(x)$. Todas son la misma 
	cosa. Usaremos $p_X(x)$ por claridad, pero reconoceremos cualquiera 
	de estas notaciones en evaluaciones.
\end{observacion}

\subsection{Propiedades}

La función de masa de probabilidad satisface dos propiedades que se 
deducen directamente de los axiomas de Kolmogórov.

\begin{proposicion}{Propiedades de la función de masa}{props_fmp}
	Sea $p_X$ la función de masa de probabilidad de una variable 
	aleatoria discreta $X$. Entonces:
	
	\medskip
	\noindent\textbf{P1)} $0 \leq p_X(x) \leq 1$ para todo 
	$x \in \mathbb{R}$.
	
	\medskip
	\noindent\textbf{P2)} $\displaystyle\sum_{x \in \text{Rec}(X)} 
	p_X(x) = 1$.
\end{proposicion}

\begin{proof}
	\textbf{P1)} Por definición $p_X(x) = P(\{X = x\})$. Como $P$ 
	satisface A1:
	\[
	P(\{X = x\}) \geq 0
	\]
	Como $\{X = x\} \subseteq \Omega$ y $P(\Omega) = 1$, por la 
	Proposición~\ref{prop:monotonia}:
	\[
	P(\{X = x\}) \leq P(\Omega) = 1
	\]
	Luego $0 \leq p_X(x) \leq 1$. \hfill$\square$
	
	\medskip
	\textbf{P2)} Los eventos $\{X = x\}$ para $x \in \text{Rec}(X)$ 
	son mutuamente excluyentes y exhaustivos, luego forman una 
	partición de $\Omega$. Aplicando A3:
	\[
	\sum_{x \in \text{Rec}(X)} p_X(x) 
	= \sum_{x \in \text{Rec}(X)} P(X = x)
	= P\!\left(\bigcup_{x \in \text{Rec}(X)} \{X = x\}\right)
	= P(\Omega) = 1 \qedhere
	\]
\end{proof}

\begin{observacion}
	Estas dos propiedades también caracterizan completamente a las 
	funciones de masa: cualquier función $p: \mathbb{R} \to [0,1]$ que 
	satisfaga P1 y P2 es la función de masa de alguna variable 
	aleatoria discreta. Son el análogo discreto de los axiomas de 
	Kolmogórov para la función de probabilidad sobre eventos.
\end{observacion}

\section{Funci\'{o}n de distribuci\'{o}n acumulada}


La función de masa $p_X(x)$ responde preguntas del tipo $P(X = x)$ — 
la probabilidad de un valor exacto. Sin embargo, en la práctica casi 
siempre nos interesan preguntas sobre intervalos: ¿cuál es la 
probabilidad de que $X$ no supere cierto umbral? ¿De que caiga entre 
dos valores? La función de distribución acumulada es la herramienta 
que responde esas preguntas de manera eficiente.

\begin{definicion}{Función de distribución acumulada}{fda}
	Sea $X$ una variable aleatoria discreta con función de masa 
	$p_X$. La \textbf{función de distribución acumulada} (FDA) de 
	$X$ es la función $F_X : \mathbb{R} \to [0,1]$ definida por:
	\[
	F_X(x) = P(X \leq x) = \sum_{t \leq x} p_X(t) 
	\qquad \text{para todo } x \in \mathbb{R}
	\]
	donde la suma recorre todos los valores $t \in \text{Rec}(X)$ 
	con $t \leq x$.
\end{definicion}

\subsection{Propiedades}

\begin{proposicion}{Propiedades de la FDA}{props_fda}
	Sea $F_X$ la función de distribución acumulada de una variable 
	aleatoria discreta $X$. Entonces:
	
	\medskip
	\noindent\textbf{P1)} $F_X$ es \textbf{no decreciente}: si 
	$a \leq b$ entonces $F_X(a) \leq F_X(b)$.
	
	\medskip
	\noindent\textbf{P2)} $F_X$ tiene \textbf{límites en el infinito}:
	\[
	\lim_{x \to -\infty} F_X(x) = 0 \qquad \text{y} \qquad 
	\lim_{x \to +\infty} F_X(x) = 1
	\]
	
	\medskip
	\noindent\textbf{P3)} $F_X$ es \textbf{continua por la derecha}: 
	para todo $x \in \mathbb{R}$,
	\[
	\lim_{t \to x^+} F_X(t) = F_X(x)
	\]
\end{proposicion}

\begin{proof}
	\textbf{P1)} Si $a \leq b$, entonces el conjunto 
	$\{t \in \text{Rec}(X) : t \leq a\} \subseteq 
	\{t \in \text{Rec}(X) : t \leq b\}$. Como $p_X(t) \geq 0$ 
	para todo $t$, sumar sobre un conjunto más grande no puede 
	disminuir el resultado:
	\[
	F_X(a) = \sum_{t \leq a} p_X(t) \leq \sum_{t \leq b} 
	p_X(t) = F_X(b) \qquad \square
	\]
	
	\medskip
	\textbf{P2)} Cuando $x \to -\infty$, no hay valores de 
	$\text{Rec}(X)$ que satisfagan $t \leq x$, luego la suma es 
	vacía y vale $0$. Cuando $x \to +\infty$, todos los valores de 
	$\text{Rec}(X)$ satisfacen $t \leq x$, luego por P2 de la 
	Proposición~\ref{prop:props_fmp}:
	\[
	\lim_{x \to +\infty} F_X(x) = \sum_{t \in \text{Rec}(X)} 
	p_X(t) = 1 \qquad \square
	\]
	
	\medskip
	\textbf{P3)} Para $t$ suficientemente cerca de $x$ por la 
	derecha, no hay valores de $\text{Rec}(X)$ en el intervalo 
	$(x, t]$ — el recorrido es discreto, luego existe un salto 
	mínimo entre valores consecutivos. Por tanto $F_X(t) = F_X(x)$ 
	para $t$ suficientemente próximo a $x$ por la derecha, lo que 
	prueba la continuidad por la derecha. $\qedhere$
\end{proof}

\begin{observacion}
	La propiedad P3 explica la convención de los puntos 
	\textbf{cerrados} y \textbf{abiertos} en la gráfica de la FDA: 
	en cada salto, el punto cerrado (valor que toma $F_X$) está 
	siempre a la derecha, reflejando que $F_X$ incluye el valor 
	$x$ mismo en la suma — es decir, $P(X \leq x)$ y no 
	$P(X < x)$.
\end{observacion}

\subsection{Cálculo de probabilidades con la FDA}

La utilidad principal de la FDA es que permite calcular 
probabilidades sobre intervalos con una simple resta, sin necesidad 
de sumar término a término.

\begin{proposicion}{Probabilidades de intervalos}{intervalos_fda}
	Sea $F_X$ la FDA de una variable aleatoria discreta $X$. 
	Para cualquier $a < b$:
	\[
	P(a < X \leq b) = F_X(b) - F_X(a)
	\]
\end{proposicion}

\begin{proof}
	Los eventos $\{X \leq a\}$ y $\{a < X \leq b\}$ son mutuamente 
	excluyentes y su unión es $\{X \leq b\}$. Por A3:
	\[
	P(X \leq b) = P(X \leq a) + P(a < X \leq b)
	\]
	Despejando:
	\[
	P(a < X \leq b) = P(X \leq b) - P(X \leq a) = 
	F_X(b) - F_X(a) \qedhere
	\]
\end{proof}

\begin{observacion}
	Nótese que el intervalo es abierto en $a$ y cerrado en $b$: 
	$P(a < X \leq b)$ excluye $a$ e incluye $b$. Esto es 
	consecuencia directa de la continuidad por la derecha de $F_X$. 
	Para otros tipos de intervalos:
	\medskip
	\begin{center}
		\begin{tabular}{@{}ll@{}}
			$P(X > a)$       & $= 1 - F_X(a)$ \\[4pt]
			$P(X \geq a)$    & $= 1 - F_X(a) + p_X(a) = 
			1 - F_X(a^-)$ \\[4pt]
			$P(a \leq X \leq b)$ & $= F_X(b) - F_X(a) + p_X(a)$ \\
		\end{tabular}
	\end{center}
\end{observacion}

\subsection*{Ejemplo: demanda diaria de un tractor}

Una empresa arrienda su tractor para trabajos de remoción. La demanda 
diaria $X$ — el número de veces que se usa el tractor en un día — 
tiene la siguiente distribución de probabilidad, donde $P(5)$ se 
deduce de la propiedad P2 de la función de masa: la suma de todas 
las probabilidades debe ser $1$, luego 
$p_X(5) = 1 - 0.10 - 0.20 - 0.30 - 0.25 - 0.10 = 0.05$.

\medskip
\begin{center}
	\begin{tabular}{@{}ccc@{}}
		\toprule
		$x$ & $p_X(x)$ & $F_X(x)$ \\
		\midrule
		0 & 0.10 & 0.10 \\
		1 & 0.20 & 0.30 \\
		2 & 0.30 & 0.60 \\
		3 & 0.25 & 0.85 \\
		4 & 0.10 & 0.95 \\
		5 & 0.05 & 1.00 \\
		\bottomrule
	\end{tabular}
\end{center}
\medskip

La figura \ref{fig:fda} muestra la FDA como función escalera. Cada 
salto ocurre en un valor de $\text{Rec}(X)$ y su altura es 
exactamente $p_X(x)$.

\begin{figure}[H]
	\centering
	\begin{tikzpicture}
		\begin{axis}[
			width=11cm, height=7cm,
			xmin=-0.5, xmax=6.5,
			ymin=-0.05, ymax=1.15,
			xlabel={$x$},
			ylabel={$F_X(x)$},
			axis lines=left,
			xtick={0,1,2,3,4,5},
			ytick={0,0.10,0.30,0.60,0.85,0.95,1.00},
			yticklabels={$0$,$0.10$,$0.30$,$0.60$,$0.85$,$0.95$,$1.00$},
			grid=both,
			grid style={line width=0.2pt, draw=gray!30},
			tick label style={font=\small},
			label style={font=\small},
			]
			\addplot[thick,azul] coordinates {(-0.5,0) (0,0)};
			\addplot[thick,azul] coordinates {(0,0.10) (1,0.10)};
			\addplot[thick,azul] coordinates {(1,0.30) (2,0.30)};
			\addplot[thick,azul] coordinates {(2,0.60) (3,0.60)};
			\addplot[thick,azul] coordinates {(3,0.85) (4,0.85)};
			\addplot[thick,azul] coordinates {(4,0.95) (5,0.95)};
			\addplot[thick,azul] coordinates {(5,1.00) (6.5,1.00)};
			\addplot[only marks, mark=*, mark size=2.5pt, azul]
			coordinates {(0,0.10)(1,0.30)(2,0.60)(3,0.85)(4,0.95)(5,1.00)};
			\addplot[only marks, mark=o, mark size=2.5pt, azul,
			every mark/.append style={fill=white}]
			coordinates {(0,0)(1,0.10)(2,0.30)(3,0.60)(4,0.85)(5,0.95)};
			\node[font=\scriptsize,rojo] at (axis cs:0.5,0.05) {$+0.10$};
			\node[font=\scriptsize,rojo] at (axis cs:1.5,0.20) {$+0.20$};
			\node[font=\scriptsize,rojo] at (axis cs:2.5,0.45) {$+0.30$};
			\node[font=\scriptsize,rojo] at (axis cs:3.5,0.73) {$+0.25$};
			\node[font=\scriptsize,rojo] at (axis cs:4.5,0.90) {$+0.10$};
		\end{axis}
	\end{tikzpicture}
	\caption{Función de distribución acumulada $F_X$ para el ejemplo del 
		tractor. Los puntos cerrados indican continuidad por la derecha. 
		Los saltos en rojo corresponden a $p_X(x)$.}
	\label{fig:fda}
\end{figure}

Usando la proposición de intervalos:
\[
P(X > 3) = 1 - F_X(3) = 1 - 0.85 = 0.15
\]
\[
P(1 < X \leq 4) = F_X(4) - F_X(1) = 0.95 - 0.30 = 0.65
\]

\section{Valor esperado y varianza}

\subsection{Valor esperado}

Sabemos que una variable aleatoria discreta $X$ puede tomar distintos 
valores con distintas probabilidades. Una pregunta natural es: ¿cuál 
es el valor \textit{típico} o \textit{central} de $X$? La respuesta 
es el \textbf{valor esperado} — un promedio ponderado de los valores 
posibles de $X$, donde cada valor se pondera por su probabilidad.

\begin{definicion}{Valor esperado}{ev}
	Sea $X$ una variable aleatoria discreta con función de masa 
	$p_X$ y recorrido $\text{Rec}(X)$. El \textbf{valor esperado}, 
	\textbf{valor medio} o \textbf{esperanza} de $X$, denotado 
	$E(X)$, $\mu_X$ o $\mu$, es:
	\[
	E(X) = \sum_{x \in \text{Rec}(X)} x \cdot p_X(x)
	\]
\end{definicion}

La definición del valor esperado no es arbitraria — tiene un origen 
físico preciso. En mecánica clásica, el \textbf{centro de masa} de un 
sistema de partículas con masas $m_i$ en posiciones $x_i$ es:
\[
\bar{x} = \frac{\sum_i m_i x_i}{\sum_i m_i}
\]
Si normalizamos las masas para que sumen $1$ — definiendo $p_i = 
m_i / \sum_j m_j$ — obtenemos exactamente:
\[
\bar{x} = \sum_i x_i\, p_i = E(X)
\]
El valor esperado es el punto donde la distribución estaría en 
equilibrio si la colocáramos sobre un fulcro. La figura 
\ref{fig:fisica_ev} muestra este paralelismo explícitamente.

\begin{figure}[H]
	\centering
	\begin{tikzpicture}[scale=1.0]
		\begin{scope}[shift={(0,0)}]
			\node[font=\small\bfseries] at (3,4.2) {Sistema físico};
			\draw[thick] (0,0) -- (6,0);
			\fill[azul!60] (1,0) circle (8pt);
			\node[font=\scriptsize,white] at (1,0) {$m_1$};
			\node[font=\scriptsize,below] at (1,-0.4) {$x_1$};
			\fill[azul!60] (2.4,0) circle (11pt);
			\node[font=\scriptsize,white] at (2.4,0) {$m_2$};
			\node[font=\scriptsize,below] at (2.4,-0.4) {$x_2$};
			\fill[azul!60] (3.8,0) circle (13pt);
			\node[font=\scriptsize,white] at (3.8,0) {$m_3$};
			\node[font=\scriptsize,below] at (3.8,-0.4) {$x_3$};
			\fill[azul!60] (5.2,0) circle (6pt);
			\node[font=\scriptsize,white] at (5.2,0) {$m_4$};
			\node[font=\scriptsize,below] at (5.2,-0.4) {$x_4$};
			\draw[rojo,very thick,dashed] (3.1,0) -- (3.1,1.2);
			\node[rojo,font=\small] at (3.1,1.5) 
			{$\bar{x} = \dfrac{\sum m_i x_i}{\sum m_i}$};
			\fill[rojo!70] (3.1,-0.15) -- (2.7,-0.7) -- (3.5,-0.7) -- cycle;
			\draw[thick] (2.5,-0.7) -- (3.7,-0.7);
			\node[font=\scriptsize,rojo] at (3.1,-1.0) {centro de masa};
		\end{scope}
		\node[font=\Large] at (7,0) {$\equiv$};
		\begin{scope}[shift={(7.5,0)}]
			\node[font=\small\bfseries] at (3,4.2) {Modelo probabilístico};
			\draw[thick] (0,0) -- (6,0);
			\foreach \x/\h/\lab in {1/1.2/0.20, 2.4/1.8/0.30, 3.8/2.1/0.35, 
				5.2/0.9/0.15}{
				\fill[azul!30] (\x-0.25,0) rectangle (\x+0.25,\h);
				\draw[azul,thick] (\x-0.25,0) rectangle (\x+0.25,\h);
				\node[font=\scriptsize,azul] at (\x,\h+0.2) {$\lab$};
			}
			\foreach \x/\lab in {1/$x_1$, 2.4/$x_2$, 3.8/$x_3$, 5.2/$x_4$}{
				\node[font=\scriptsize,below] at (\x,-0.4) {\lab};
			}
			\draw[rojo,very thick,dashed] (3.1,0) -- (3.1,3.0);
			\node[rojo,font=\small] at (3.1,3.4) {$E(X) = \sum x\, p_X(x)$};
			\fill[rojo!70] (3.1,-0.15) -- (2.7,-0.7) -- (3.5,-0.7) -- cycle;
			\draw[thick] (2.5,-0.7) -- (3.7,-0.7);
			\node[font=\scriptsize,rojo] at (3.1,-1.0) {valor esperado};
		\end{scope}
	\end{tikzpicture}
	\caption{El valor esperado como centro de masa. Las masas físicas 
		$m_i$ normalizadas corresponden a las probabilidades $p_X(x_i)$; 
		el punto de equilibrio físico corresponde a $E(X)$.}
	\label{fig:fisica_ev}
\end{figure}

\begin{proposicion}{Valor esperado de una función de $X$}{ev_hx}
	Sea $h : \mathbb{R} \to \mathbb{R}$ una función. El valor 
	esperado de $h(X)$ es:
	\[
	E(h(X)) = \sum_{x \in \text{Rec}(X)} h(x) \cdot p_X(x)
	\]
\end{proposicion}

\begin{proposicion}{Linealidad del valor esperado}{ev_lineal}
	Si $a$ y $b$ son constantes, entonces:
	\[
	E(aX + b) = aE(X) + b
	\]
\end{proposicion}

\begin{proof}
	Aplicamos la Proposición~\ref{prop:ev_hx} con $h(x) = ax + b$:
	\[
	E(aX + b) 
	= \sum_{x \in \text{Rec}(X)} (ax + b)\, p_X(x)
	= a \sum_{x \in \text{Rec}(X)} x\, p_X(x) 
	+ b \sum_{x \in \text{Rec}(X)} p_X(x)
	= aE(X) + b \qedhere
	\]
	donde en el último paso usamos la definición de $E(X)$ y la 
	propiedad P2 de la función de masa 
	($\sum p_X(x) = 1$).
\end{proof}

\subsection*{Ejemplo: demanda mensual}

Una empresa registra la demanda mensual $X$ de un producto durante 
dos años y obtiene la siguiente distribución:

\medskip
\begin{center}
	\begin{tabular}{@{}ccccc@{}}
		\toprule
		$x$ & 300 & 400 & 500 & 600 \\
		\midrule
		$p_X(x)$ & 0.20 & 0.30 & 0.35 & 0.15 \\
		\bottomrule
	\end{tabular}
\end{center}
\medskip

El valor esperado es:
\[
E(X) = 300(0.20) + 400(0.30) + 500(0.35) + 600(0.15)
= 60 + 120 + 175 + 90 = 445
\]

La figura~\ref{fig:ev} muestra la función de masa con el valor 
esperado marcado como centro de la distribución.

\begin{figure}[H]
	\centering
	\begin{tikzpicture}
		\begin{axis}[
			width=11cm, height=6cm,
			xmin=250, xmax=650,
			ymin=0, ymax=0.60,
			xlabel={$x$},
			ylabel={$p_X(x)$},
			axis lines=left,
			xtick={300,400,445,500,600},
			xticklabels={$300$,$400$,$\mu$,$500$,$600$},
			ytick={0.10,0.20,0.30,0.35},
			yticklabels={$0.10$,$0.20$,$0.30$,$0.35$},
			tick label style={font=\small},
			label style={font=\small},
			bar width=18pt,
			]
			\addplot[ybar, fill=azul!30, draw=azul]
			coordinates {(300,0.20)(400,0.30)(500,0.35)(600,0.15)};
			\addplot[thick, rojo, dashed] coordinates {(445,0) (445,0.48)};
			\node[font=\small, rojo, above right] at (axis cs:445,0.48) {$E(X)=445$};
		\end{axis}
	\end{tikzpicture}
	\caption{Función de masa $p_X$ para la demanda mensual. La línea 
		roja marca el valor esperado $E(X) = 445$ — el centro de masa de 
		la distribución.}
	\label{fig:ev}
\end{figure}

\subsection{Varianza}

El valor esperado nos dice dónde está el centro de la distribución, 
pero no dice nada sobre qué tan dispersos están los valores de $X$ 
alrededor de ese centro. Dos distribuciones pueden tener el mismo 
$E(X)$ y ser completamente distintas en su dispersión. La 
\textbf{varianza} cuantifica esa dispersión.

\begin{definicion}{Varianza}{var}
	Sea $X$ una variable aleatoria discreta con $E(X) = \mu$. La 
	\textbf{varianza} de $X$, denotada $\text{Var}(X)$, $V(X)$, 
	$\sigma^2_X$ o $\sigma^2$, es:
	\[
	\text{Var}(X) = E\!\left[(X - \mu)^2\right] 
	= \sum_{x \in \text{Rec}(X)} (x - \mu)^2\, p_X(x)
	\]
	La \textbf{desviación estándar} de $X$ es $\sigma_X = 
	\sqrt{\text{Var}(X)}$.
\end{definicion}

La razón de elevar al cuadrado en la definición de la varianza es 
puramente \textbf{matemática}, y se entiende mejor a través del 
\textbf{método de mínimos cuadrados}. El problema fundamental es 
el siguiente: dado un conjunto de observaciones $x_1, x_2, \ldots, 
x_n$, ¿qué valor $\mu$ minimiza la dispersión total? Es decir, 
¿cuál es el mejor representante central de los datos?

La respuesta depende de cómo midamos la dispersión. Si definimos 
la dispersión total como la suma de las desviaciones al cuadrado:
\[
S(\mu) = \sum_{i=1}^n (x_i - \mu)^2
\]
entonces minimizar $S(\mu)$ respecto a $\mu$ es trivial — derivamos 
e igualamos a cero:
\[
\frac{dS}{d\mu} = -2\sum_{i=1}^n (x_i - \mu) = 0 
\implies \mu = \frac{1}{n}\sum_{i=1}^n x_i
\]
El minimizador es exactamente la \textbf{media aritmética}. Esto no 
es una coincidencia — es la razón por la que la media y la varianza 
están tan profundamente ligadas: la varianza mide la dispersión 
respecto al único punto que la minimiza.

Gauss y Legendre eligieron el cuadrado sobre cualquier alternativa 
porque produce estas propiedades matemáticas superiores:

\begin{itemize}
	\item Es \textbf{convexa} — tiene un único mínimo global, 
	garantizando que la solución es única.
	\item Es \textbf{diferenciable en todas partes} — permite 
	minimizar analíticamente sin casos especiales.
	\item Tiene \textbf{solución analítica cerrada} — la media, 
	obtenida directamente sin métodos numéricos.
	\item Conecta naturalmente con la \textbf{distribución normal}, 
	cuya densidad es proporcional a $e^{-x^2/2}$.
\end{itemize}

La alternativa más intuitiva habría sido minimizar la suma de valores 
absolutos $\sum_i |x_i - \mu|$ — la llamada \textbf{desviación media 
	absoluta} — que tiene interpretación más directa pero no es 
diferenciable en $x = 0$ y no tiene solución analítica cerrada. 
El cuadrado ganó por sus propiedades matemáticas, no por intuición.

El costo de elevar al cuadrado es \textbf{dimensional}: si $X$ mide 
demanda en unidades de producto, $\text{Var}(X)$ está en 
unidades$^2$, lo que hace imposible comparar la dispersión 
directamente con el valor esperado. La \textbf{desviación estándar} 
$\sigma_X = \sqrt{\text{Var}(X)}$ existe precisamente para resolver 
eso — devuelve la dispersión a las mismas unidades que $X$, 
permitiendo afirmaciones del tipo \textit{«el valor típico de $X$ 
	se aleja de $\mu$ aproximadamente $\sigma_X$ unidades»}. En la 
práctica, la varianza se usa en demostraciones y cálculos por su 
mejor comportamiento algebraico, mientras que la desviación estándar 
se usa para interpretar y comunicar resultados.

El paralelismo con la física se extiende a la varianza. En mecánica 
clásica, el \textbf{momento de inercia} de un sistema de partículas 
respecto a un eje que pasa por el centro de masa es:
\[
I = \sum_i m_i (x_i - \bar{x})^2
\]
Bajo normalización de las masas, esto es exactamente:
\[
I = \sum_i (x_i - \bar{x})^2\, p_i = \text{Var}(X)
\]
La varianza mide lo mismo que el momento de inercia: qué tan 
\textit{lejos} están las masas — o las probabilidades — del centro. 
Un sistema con masas muy dispersas tiene momento de inercia grande; 
una distribución con valores muy dispersos alrededor de $\mu$ tiene 
varianza grande. La figura~\ref{fig:fisica_var} ilustra esta 
equivalencia.

\begin{figure}[H]
	\centering
	\begin{tikzpicture}[scale=1.0]
		% ============ PANEL IZQUIERDO: física ============
		\begin{scope}[shift={(0,0)}]
			\node[font=\small\bfseries] at (3,4.0) {Sistema físico};
			\draw[thick] (0,0) -- (6.5,0);
			\draw[rojo, very thick, dashed] (3.1,-1.2) -- (3.1,3.2);
			\node[rojo, font=\scriptsize] at (3.1,3.45) {$\bar{x}$};
			\foreach \x/\r in {1.0/8pt, 2.4/11pt, 3.8/13pt, 5.2/6pt}{
				\draw[gray, thin] (\x, 0) -- (\x, -0.5);
				\fill[azul!60] (\x, 0) circle (\r);
			}
			\node[font=\scriptsize, white] at (1.0,0) {$m_1$};
			\node[font=\scriptsize, white] at (2.4,0) {$m_2$};
			\node[font=\scriptsize, white] at (3.8,0) {$m_3$};
			\node[font=\scriptsize, white] at (5.2,0) {$m_4$};
			\draw[verde, very thick] (1.0,-0.5) -- (3.1,-0.5);
			\draw[verde, thick] (1.0,-0.55) -- (1.0,-0.45);
			\draw[verde, thick] (3.1,-0.55) -- (3.1,-0.45);
			\node[verde, font=\scriptsize] at (2.05,-0.82) {$|x_1-\bar{x}|$};
			\draw[azul!70, very thick] (2.4,-0.5) -- (3.1,-0.5);
			\draw[Goldenrod!80!black, very thick] (3.1,-0.5) -- (3.8,-0.5);
			\draw[rojo!70, very thick] (3.1,-0.5) -- (5.2,-0.5);
			\draw[rojo!70, thick] (5.2,-0.55) -- (5.2,-0.45);
			\node[rojo!70, font=\scriptsize] at (4.15,-0.82) {$|x_4-\bar{x}|$};
			\node[font=\scriptsize] at (3.1,-1.5) 
			{$I = \displaystyle\sum_i m_i(x_i-\bar{x})^2$};
			\node[font=\scriptsize] at (3.1,-2.0) {momento de inercia};
		\end{scope}
	% Signo de igualdad + etiqueta de mapeo
	\node[font=\Large] at (7.2, 0.5) {$\equiv$};
	\node[font=\scriptsize, gray, align=center] at (7.2, -0.75) 
	{masa (área)\\$=$\\prob. (altura)};
		% ============ PANEL DERECHO: probabilidad ============
		\begin{scope}[shift={(7.8,0)}]
			\node[font=\small\bfseries] at (3,4.0) {Modelo probabilístico};
			\draw[thick] (0,0) -- (6.5,0);
			\draw[rojo, very thick, dashed] (3.1,-1.2) -- (3.1,3.2);
			\node[rojo, font=\scriptsize] at (3.1,3.45) {$\mu$};
			\foreach \x/\h in {1.0/1.2, 2.4/1.8, 3.8/2.1, 5.2/0.9}{
				\fill[azul!30] (\x-0.22,0) rectangle (\x+0.22,\h);
				\draw[azul, thick] (\x-0.22,0) rectangle (\x+0.22,\h);
				\draw[gray, thin] (\x,0) -- (\x,-0.5);
			}
			\draw[verde, very thick] (1.0,-0.5) -- (3.1,-0.5);
			\draw[verde, thick] (1.0,-0.55) -- (1.0,-0.45);
			\draw[verde, thick] (3.1,-0.55) -- (3.1,-0.45);
			\node[verde, font=\scriptsize] at (2.05,-0.82) {$|x_1-\mu|$};
			\draw[azul!70, very thick] (2.4,-0.5) -- (3.1,-0.5);
			\draw[Goldenrod!80!black, very thick] (3.1,-0.5) -- (3.8,-0.5);
			\draw[rojo!70, very thick] (3.1,-0.5) -- (5.2,-0.5);
			\draw[rojo!70, thick] (5.2,-0.55) -- (5.2,-0.45);
			\node[rojo!70, font=\scriptsize] at (4.15,-0.82) {$|x_4-\mu|$};
			\node[font=\scriptsize] at (3.1,-1.5)
			{$\text{Var}(X) = \displaystyle\sum_x (x-\mu)^2 p_X(x)$};
			\node[font=\scriptsize] at (3.1,-2.0) {varianza};
		\end{scope}
	\end{tikzpicture}
	\caption{La varianza como momento de inercia. Los segmentos de color 
		bajo el eje muestran la distancia $|x_i - \bar{x}|$ de cada valor 
		al centro. La varianza pondera el \textit{cuadrado} de esas distancias 
		por la masa (o probabilidad) correspondiente.}
	\label{fig:fisica_var}
\end{figure}

\begin{observacion}
	La varianza promedia el \textit{cuadrado} de las desviaciones 
	respecto a $\mu$, no las desviaciones directas. Esto tiene dos 
	razones: primero, las desviaciones positivas y negativas no se 
	cancelan; segundo, las desviaciones grandes pesan 
	desproporcionadamente más, lo que hace a la varianza sensible 
	a valores extremos. La desviación estándar $\sigma_X$ devuelve 
	la dispersión a las unidades originales de $X$.
\end{observacion}

\begin{figure}[H]
	\centering
	\begin{tikzpicture}
		\begin{axis}[
			width=11cm, height=5.5cm,
			xmin=250, xmax=650,
			ymin=-0.08, ymax=0.50,
			xlabel={$x$},
			ylabel={$p_X(x)$},
			axis lines=left,
			xtick={300,400,500,600},
			ytick={0.10,0.20,0.30,0.35},
			yticklabels={$0.10$,$0.20$,$0.30$,$0.35$},
			tick label style={font=\small},
			label style={font=\small},
			]
			\addplot[ybar, bar width=30pt, fill=azul!30, draw=azul]
			coordinates {(300,0.20)(400,0.30)(500,0.35)(600,0.15)};
			\addplot[thick, rojo, dashed] coordinates {(445,0) (445,0.42)};
			\node[font=\small, rojo] at (axis cs:445,0.44) {$\mu=445$};
		\draw[<->,verde,thick] (axis cs:300,-0.04) -- (axis cs:445,-0.04);
		\node[font=\scriptsize,verde] at (axis cs:372,-0.062) 
		{$|300-\mu|=145$};
		\draw[<->,verde,thick] (axis cs:445,-0.04) -- (axis cs:600,-0.04);
		\node[font=\scriptsize,verde] at (axis cs:522,-0.062) 
		{$|600-\mu|=155$};
		\end{axis}
	\end{tikzpicture}
	\caption{La varianza mide la dispersión de $X$ alrededor de 
		$\mu = E(X) = 445$. Las flechas ilustran las desviaciones 
		$|x - \mu|$ para los valores extremos del recorrido.}
	\label{fig:var}
\end{figure}

\begin{proposicion}{Fórmula alternativa de la varianza}{var_alt}
	\[
	\text{Var}(X) = E(X^2) - (E(X))^2
	\]
\end{proposicion}

\begin{proof}
	Sea $\mu = E(X)$. Expandimos el cuadrado en la definición:
	\[
	\text{Var}(X) 
	= \sum_{x \in \text{Rec}(X)} (x - \mu)^2\, p_X(x)
	= \sum_{x} (x^2 - 2\mu x + \mu^2)\, p_X(x)
	\]
	Distribuimos la suma:
	\[
	= \sum_{x} x^2\, p_X(x) 
	- 2\mu \sum_{x} x\, p_X(x) 
	+ \mu^2 \sum_{x} p_X(x)
	\]
	Reconocemos cada término: el primero es $E(X^2)$ por la 
	Proposición~\ref{prop:ev_hx} con $h(x) = x^2$; el segundo 
	usa $\sum x\, p_X(x) = \mu$; el tercero usa $\sum p_X(x) = 1$:
	\[
	= E(X^2) - 2\mu \cdot \mu + \mu^2 \cdot 1
	= E(X^2) - 2\mu^2 + \mu^2
	= E(X^2) - \mu^2
	= E(X^2) - (E(X))^2 \qedhere
	\]
\end{proof}

\begin{observacion}
	La fórmula alternativa es la que se usa en la práctica: calcular 
	$E(X^2) - (E(X))^2$ requiere solo dos sumas, mientras que la 
	definición directa requiere calcular $(x - \mu)^2$ para cada 
	valor del recorrido — más laborioso cuando el recorrido es grande.
\end{observacion}

\begin{proposicion}{Propiedades de la varianza}{props_var}
	Sean $a$ y $b$ constantes. Entonces:
	
	\medskip
	\noindent\textbf{i)} $\text{Var}(X) \geq 0$
	
	\medskip
	\noindent\textbf{ii)} $\text{Var}(a) = 0$
	
	\medskip
	\noindent\textbf{iii)} $\text{Var}(aX + b) = a^2\,\text{Var}(X)$
	\quad y \quad $\sigma_{aX+b} = |a|\,\sigma_X$
\end{proposicion}

\begin{proof}
	\textbf{i)} Cada término $(x-\mu)^2 p_X(x) \geq 0$, luego 
	la suma también. $\square$
	
	\medskip
	\textbf{ii)} Si $a$ es constante, $E(a) = a$ y 
	$(a - E(a))^2 = 0$ para todo resultado, luego 
	$\text{Var}(a) = 0$. $\square$
	
	\medskip
	\textbf{iii)} Por la linealidad del valor esperado 
	(Proposición~\ref{prop:ev_lineal}), $E(aX+b) = aE(X)+b = 
	a\mu + b$. Aplicando la fórmula alternativa:
	\[
	\text{Var}(aX+b) 
	= E\!\left[(aX+b)^2\right] - (aE(X)+b)^2
	\]
	\[
	= E\!\left[a^2X^2 + 2abX + b^2\right] - 
	(a^2\mu^2 + 2ab\mu + b^2)
	\]
	\[
	= a^2E(X^2) + 2ab\mu + b^2 - a^2\mu^2 - 2ab\mu - b^2
	= a^2\!\left(E(X^2) - \mu^2\right)
	= a^2\,\text{Var}(X)
	\]
	La desviación estándar sigue como 
	$\sigma_{aX+b} = \sqrt{a^2\,\text{Var}(X)} = 
	|a|\,\sigma_X$. $\qedhere$
\end{proof}

\begin{observacion}
	La propiedad iii revela algo importante: sumar una constante 
	$b$ a $X$ \textbf{no cambia la varianza} — desplaza la 
	distribución pero no altera su dispersión. Solo la 
	multiplicación por $a$ escala la varianza, y lo hace de manera 
	cuadrática.
\end{observacion}

\section{Distribuciones importantes}

Hasta ahora hemos definido la variable aleatoria discreta en su forma más
general: una función $X:\Omega\to\mathbb{R}$ equipada con una función de masa
$p_X$ arbitraria. Sin embargo, en la práctica aparecen ciertas
\textbf{estructuras de experimento} que se repiten una y otra vez ---y cada
estructura genera siempre la misma forma de distribución.

\subsection{Distribuci\'{o}n de Bernoulli}

\begin{definicion}{Experimento de Bernoulli}{exp_bernoulli}
	Sea $(\Omega, \mathcal{A}, P)$ un espacio de probabilidad. Un
	\textbf{experimento de Bernoulli} es aquel en que el resultado de
	interés queda determinado por un evento $A \in \mathcal{A}$,
	llamado \textbf{éxito}, con $P(A) = p \in (0,1)$. Su complemento
	$A^c$ se denomina \textbf{fracaso} y tiene probabilidad $1-p$.
\end{definicion}

El espacio muestral $\Omega$ puede tener cualquier cardinalidad ---lo que
define al experimento como de Bernoulli es la \textbf{partición binaria}
$\{A,\, A^c\}$ del espacio, no la estructura de $\Omega$. Por ejemplo,
al lanzar un dado de seis caras y definir $A = \{\text{sale par}\}$, el
espacio muestral tiene seis elementos, pero el experimento es de Bernoulli
con $p = 1/2$.

\begin{definicion}{Variable aleatoria Bernoulli}{va_bernoulli}
	Sea $(\Omega, \mathcal{A}, P)$ el espacio de probabilidad de un
	experimento de Bernoulli con evento de éxito $A$ y probabilidad
	$p = P(A)$. La \textbf{variable aleatoria Bernoulli} asociada a $A$
	es la \textbf{función indicadora} de $A$:
	\[
	X(\omega) = \mathbf{1}_A(\omega) \;=\;
	\begin{cases}
		1 & \text{si } \omega \in A, \\
		0 & \text{si } \omega \notin A.
	\end{cases}
	\]
	Se escribe $X \sim \mathrm{Bernoulli}(p)$.
\end{definicion}

\begin{proposicion}{Funci\'{o}n de masa Bernoulli}{pmf_bern}
	Si $X \sim \mathrm{Bernoulli}(p)$, entonces
	\[
	p_X(x) = p^x(1-p)^{1-x}, \qquad x \in \{0,1\}.
	\]
\end{proposicion}

\begin{proof}
	Directamente de la definición de $X$ y de $P(A) = p$:
	\[
	p_X(1) = P(X=1) = P(A) = p,
	\qquad
	p_X(0) = P(X=0) = P(A^c) = 1-p,
	\]
	que se unifica en $p^x(1-p)^{1-x}$ para $x \in \{0,1\}$. \qedhere
\end{proof}

\begin{observacion}
	La \textbf{función indicadora} $\mathbf{1}_A : \Omega \to \{0,1\}$
	es un objeto general de la matemática: dado cualquier conjunto
	$A \subseteq \Omega$, vale $1$ exactamente donde ocurre el evento y
	$0$ en el resto. La variable aleatoria Bernoulli no se \emph{parece}
	a una función indicadora --- es literalmente la misma cosa, vista
	ahora sobre un espacio de probabilidad. En análisis aparece con el
	nombre de función característica de un conjunto y se denota
	$\chi_A$; en probabilidad el nombre estándar es el que usamos aquí.
\end{observacion}

\begin{observacion}
	La expresión $p^x(1-p)^{1-x}$ es una fórmula compacta que unifica
	ambos casos: si $x=1$, entrega $p^1(1-p)^0 = p$; si $x=0$, entrega
	$p^0(1-p)^1 = 1-p$. Es habitual denotar $q \coloneqq 1-p$.
\end{observacion}

\begin{proposicion}{Momentos de la distribuci\'{o}n Bernoulli}
	{momentos_bern}
	Si $X \sim \mathrm{Bernoulli}(p)$ con $p \in (0,1)$, entonces
	\[
	E(X) = p \qquad \text{y} \qquad \mathrm{Var}(X) = p(1-p).
	\]
\end{proposicion}

\begin{proof}
	\textbf{Valor esperado.} Aplicando la definición directamente:
	\[
	E(X) = 0\cdot p_X(0) + 1\cdot p_X(1) = 0(1-p) + 1\cdot p = p.
	\]
	
	\textbf{Varianza.} Usamos la fórmula alternativa
	$\mathrm{Var}(X) = E(X^2) - (E(X))^2$. Como $X$ solo toma valores
	en $\{0,1\}$, se tiene $x^2 = x$ para todo $x$ en el recorrido de
	$X$, luego $E(X^2) = E(X) = p$. Por lo tanto:
	\[
	\mathrm{Var}(X) = p - p^2 = p(1-p). \qedhere
	\]
\end{proof}

\begin{figure}[H]
	\centering
	\begin{tikzpicture}
		% Ejes
		\draw[thick, ->] (-0.5,0) -- (4.0,0) node[right] {$x$};
		\draw[thick, ->] (0,-0.1) -- (0,3.3)
		node[above] {$p_X(x)$};
		
		% Barra x=0: altura = (1-p)*4 = 0.65*4 = 2.6  (p=0.35)
		\fill[azul!20] (0.75,0) rectangle (1.25,2.6);
		\draw[azul,thick] (0.75,0) rectangle (1.25,2.6);
		\node[below] at (1.0,0) {$0$};
		\draw[azul,dashed,thin] (0,2.6) -- (0.75,2.6);
		\node[azul,left,font=\small] at (0,2.6) {$1{-}p$};
		
		% Barra x=1: altura = p*4 = 0.35*4 = 1.4
		\fill[rojo!20] (2.75,0) rectangle (3.25,1.4);
		\draw[rojo,thick] (2.75,0) rectangle (3.25,1.4);
		\node[below] at (3.0,0) {$1$};
		\draw[rojo,dashed,thin] (0,1.4) -- (2.75,1.4);
		\node[rojo,left,font=\small] at (0,1.4) {$p$};
		
		% Marcas en eje y
		\foreach \y in {0.5,1.0,1.5,2.0,2.5}{
			\draw[gray,very thin] (0.03,\y)--(-0.03,\y);
		}
		
		% Nota
		\node[gray,font=\small] at (2.5,3.1) {Ejemplo: $p = 0.35$};
	\end{tikzpicture}
	\caption{Función de masa de $X \sim \mathrm{Bernoulli}(0.35)$.
		La probabilidad de fracaso es $1-p = 0.65$ y la de éxito es $p = 0.35$.}
	\label{fig:pmf_bernoulli}
\end{figure}

La distribución Bernoulli modela un único ensayo. La pregunta natural que
sigue es: ¿qué ocurre si repetimos el experimento $n$ veces? Para que esa
repetición tenga sentido probabilístico se necesita precisar las condiciones
bajo las cuales opera.

\begin{definicion}{Variables aleatorias i.i.d.}{iid}
	Una colección de variables aleatorias $X_1, X_2, \ldots, X_n$ se
	dice \textbf{independientes e idénticamente distribuidas} (i.i.d.)
	si satisface dos condiciones:
	\begin{enumerate}[label=\textbf{\roman*)}]
		\item \textbf{Independencia:} para todo $(x_1, x_2, \ldots, x_n)$,
		\[
		P(X_1 = x_1,\, X_2 = x_2,\, \ldots,\, X_n = x_n)
		= \prod_{i=1}^{n} P(X_i = x_i).
		\]
		\item \textbf{Identidad de distribución:} todas las $X_i$ tienen
		la misma función de masa, es decir, $p_{X_i} = p_{X_j}$ para
		todo $i \neq j$.
	\end{enumerate}
\end{definicion}

\begin{observacion}
	El concepto i.i.d.\ es uno de los pilares de la estadística: la
	mayor parte de los resultados clásicos --- incluyendo el Teorema del
	Límite Central --- asumen que las observaciones de una muestra son
	i.i.d. La condición de independencia formaliza la intuición de que
	el resultado de un ensayo no afecta a los demás; la de identidad de
	distribución formaliza que todos los ensayos operan bajo las mismas
	condiciones.
\end{observacion}

\begin{definicion}{Proceso de Bernoulli}{proceso_bernoulli}
	Un \textbf{proceso de Bernoulli} con $n$ ensayos y probabilidad de
	éxito $p \in (0,1)$ es una colección de variables aleatorias
	i.i.d.\ $X_1, X_2, \ldots, X_n$ con
	\[
	X_i \sim \mathrm{Bernoulli}(p) \qquad \text{para todo }
	i = 1, 2, \ldots, n.
	\]
\end{definicion}

\begin{observacion}
	Las tres condiciones que se enuncian habitualmente --- ensayos
	binarios, independencia y probabilidad constante --- están todas
	contenidas en esta definición: la distribución Bernoulli garantiza
	la binaridad, la i.i.d.\ garantiza simultáneamente la independencia
	y la constancia de $p$.
\end{observacion}

\subsection{Distribuci\'{o}n Binomial}

\begin{definicion}{Distribuci\'{o}n Binomial}{dist_binomial}
	Sea $X_1, X_2, \ldots, X_n$ un proceso de Bernoulli con probabilidad
	de éxito $p \in (0,1)$ (Definición~\ref{def:proceso_bernoulli}).
	La \textbf{variable aleatoria Binomial} con parámetros $n$ y $p$ es
	la suma
	\[
	X = \sum_{i=1}^{n} X_i,
	\]
	que cuenta el número de éxitos en los $n$ ensayos. Se escribe
	$X \sim \mathrm{Binomial}(n,p)$.
\end{definicion}

\begin{proposicion}{Funci\'{o}n de masa Binomial}{pmf_binom}
	Si $X \sim \mathrm{Binomial}(n,p)$, entonces
	\[
	p_X(x) = \binom{n}{x}\, p^{\,x}(1-p)^{n-x},
	\qquad x = 0, 1, \ldots, n.
	\]
\end{proposicion}

\begin{proof}
	Sea $x \in \{0,1,\ldots,n\}$. El evento $\{X = x\}$ ocurre cuando
	exactamente $x$ de los $n$ ensayos resultan en éxito. Consideremos
	primero una secuencia específica con éxitos en las posiciones
	$i_1,\ldots,i_x$ y fracasos en las restantes $j_1,\ldots,j_{n-x}$.
	Por la independencia de las $X_i$:
	\[
	P\!\left(X_{i_1}=1,\ldots,X_{i_x}=1,\;
	X_{j_1}=0,\ldots,X_{j_{n-x}}=0\right)
	= p^{\,x}(1-p)^{n-x}.
	\]
	Esta probabilidad es la misma para \emph{toda} secuencia con
	exactamente $x$ éxitos, independientemente de en qué posiciones
	ocurran. El número de secuencias distintas con $x$ éxitos en $n$
	posiciones es $\binom{n}{x}$, y los eventos correspondientes son
	mutuamente excluyentes. Por lo tanto:
	\[
	p_X(x) = P(X=x) = \binom{n}{x}\,p^{\,x}(1-p)^{n-x}. \qedhere
	\]
\end{proof}

El nombre \emph{distribución binomial} no es arbitrario. Para que
$p_X$ sea una función de masa válida debe cumplirse que
$\sum_{x=0}^{n} p_X(x) = 1$, y la verificación revela la conexión:
\[
\sum_{x=0}^{n} \binom{n}{x}\, p^{\,x}(1-p)^{n-x}
= \bigl(p + (1-p)\bigr)^n = 1.
\]
El paso central es el \textbf{teorema del binomio}, que establece que
para cualesquiera $a, b \in \mathbb{R}$ y $n \in \mathbb{N}$,
\[
(a+b)^n = \sum_{x=0}^{n} \binom{n}{x}\, a^{\,x}\, b^{n-x}.
\]
Es decir, los coeficientes $\binom{n}{x}$ que aparecen en la función
de masa son exactamente los coeficientes que surgen al expandir un
binomio elevado a la $n$ --- de ahí el nombre. La distribución binomial
no hace más que repartir la unidad $(p + (1-p))^n = 1$ entre todos los
posibles conteos de éxitos.

\begin{proposicion}{Momentos de la distribuci\'{o}n Binomial}
	{momentos_binom}
	Si $X \sim \mathrm{Binomial}(n,p)$ con $p \in (0,1)$, entonces
	\[
	E(X) = np \qquad \text{y} \qquad \mathrm{Var}(X) = np(1-p).
	\]
\end{proposicion}

\begin{proof}
	Escribimos $X = X_1 + \cdots + X_n$ con $X_i \sim \mathrm{Bernoulli}(p)$
	i.i.d.\ (Definición~\ref{def:dist_binomial}).
	
	\textbf{Valor esperado.} Por linealidad del valor esperado
	(Proposición~\ref{prop:ev_lineal}):
	\[
	E(X) = \sum_{i=1}^{n} E(X_i) = \sum_{i=1}^{n} p = np.
	\]
	
	\textbf{Varianza.} La independencia de las $X_i$ implica que no
	hay términos cruzados en la varianza de la suma, de modo que
	\[
	\mathrm{Var}(X) = \sum_{i=1}^{n} \mathrm{Var}(X_i)
	= \sum_{i=1}^{n} p(1-p) = np(1-p). \qedhere
	\]
\end{proof}

\begin{observacion}
	La aditividad de la varianza bajo independencia --- usada en la
	demostración anterior --- es un resultado que se justificará
	formalmente en el Tema~3 al introducir el concepto de covarianza.
	Por ahora basta saber que cuando las variables son independientes,
	los términos cruzados son cero y la varianza de la suma se reduce
	a la suma de las varianzas.
\end{observacion}

\begin{figure}[H]
	\centering
	\begin{tikzpicture}[scale=0.90]
		% Binomial(n=10, p=0.30) escalado por 10
		% k=0:  0.0282  -> 0.282
		% k=1:  0.1211  -> 1.211
		% k=2:  0.2335  -> 2.335
		% k=3:  0.2668  -> 2.668  (moda)
		% k=4:  0.2001  -> 2.001
		% k=5:  0.1029  -> 1.029
		% k=6:  0.0368  -> 0.368
		% k=7:  0.0090  -> 0.090
		% k=8:  0.0014  -> 0.014
		
		\def\bw{0.30}
		
		\draw[thick,->] (-0.5,0) -- (11.0,0) node[right] {$x$};
		\draw[thick,->] (0,-0.1) -- (0,3.3) node[above] {$p_X(x)$};
		
		\draw[gray,thin,dashed] (0,2.668) -- (10.5,2.668);
		\node[gray,left,font=\scriptsize] at (0,2.668) {$0.267$};
		
		\foreach \k/\h in {%
			0/0.282, 1/1.211, 2/2.335, 3/2.668, 4/2.001,%
			5/1.029, 6/0.368, 7/0.090, 8/0.014}{
			\pgfmathsetmacro{\xp}{\k+1}
			\fill[azul!25] (\xp-\bw,0) rectangle (\xp+\bw,\h);
			\draw[azul,thick] (\xp-\bw,0) rectangle (\xp+\bw,\h);
			\node[below,font=\scriptsize] at (\xp,0) {$\k$};
		}
		
		\draw[rojo,thick,dashed] (4,0) -- (4,2.9);
		\node[rojo,above,font=\small] at (4,2.9) {$\mu = np = 3$};
		
		\node[font=\small,gray] at (7.5,3.15)
		{$X\sim\mathrm{Binomial}(10,\;0.30)$};
	\end{tikzpicture}
	\caption{Función de masa de $X\sim\mathrm{Binomial}(10,0.30)$.
		La distribución es asimétrica hacia la derecha para $p<0.5$,
		con moda en $x=3$ y media $E(X)=np=3$.}
	\label{fig:pmf_binomial}
\end{figure}

Vale la pena contrastar ambas distribuciones: Bernoulli responde
\emph{¿ocurrió el evento?} en un único ensayo --- la variable solo
puede valer $0$ o $1$. Binomial responde \emph{¿cuántas veces ocurrió?}
en $n$ ensayos --- la variable toma valores en $\{0,1,\ldots,n\}$.
La relación entre ambas es exacta: $\mathrm{Bernoulli}(p) =
\mathrm{Binomial}(1,p)$.

\subsection{Distribuci\'{o}n de Poisson}

La distribución Binomial requiere que el número de ensayos $n$ sea
finito y fijo. Sin embargo, hay fenómenos donde los eventos pueden
ocurrir en cualquier punto de un intervalo continuo --- el tiempo, el
espacio, una superficie --- y no tiene sentido hablar de un $n$ fijo.
¿Cuántos autos pasan por una intersección en una hora? ¿Cuántos
mensajes llegan a un servidor en un minuto? ¿Cuántas imperfecciones
hay en un metro de cable?

La idea para modelar este tipo de fenómenos es la siguiente: dividamos
el intervalo en $n$ subintervalos tan pequeños que en cada uno pueda
ocurrir a lo más un evento, con probabilidad $p = \lambda/n$, donde
$\lambda > 0$ es la tasa promedio de ocurrencia. Esto define un proceso
de Bernoulli con $n$ ensayos y probabilidad $\lambda/n$, de modo que
el número de eventos $X_n \sim \mathrm{Binomial}(n, \lambda/n)$. La
distribución de Poisson emerge al tomar el límite $n \to \infty$ ---
subintervalos infinitesimales.

\begin{proposicion}{L\'{\i}mite de la Binomial}{limite_binomial}
	Sea $k \in \mathbb{N}_0$ fijo y $\lambda > 0$. Entonces
	\[
	\lim_{n\to\infty}
	\binom{n}{k}
	\left(\frac{\lambda}{n}\right)^k
	\left(1-\frac{\lambda}{n}\right)^{n-k}
	= \frac{\lambda^k e^{-\lambda}}{k!}.
	\]
\end{proposicion}

\begin{proof}
	Separamos el término en tres factores:
	\[
	\binom{n}{k}
	\left(\frac{\lambda}{n}\right)^k
	\left(1-\frac{\lambda}{n}\right)^{n-k}
	= \frac{\lambda^k}{k!}
	\cdot
	\underbrace{\frac{n(n-1)\cdots(n-k+1)}{n^k}}_{(A)}
	\cdot
	\underbrace{\left(1-\frac{\lambda}{n}\right)^{n}}_{(B)}
	\cdot
	\underbrace{\left(1-\frac{\lambda}{n}\right)^{-k}}_{(C)}.
	\]
	Analizamos el límite de cada uno.
	
	\textbf{Factor (A).} Reescribiendo:
	\[
	\frac{n(n-1)\cdots(n-k+1)}{n^k}
	= \prod_{j=0}^{k-1}\!\left(1-\frac{j}{n}\right) \;\to\; 1
	\quad\text{cuando } n\to\infty.
	\]
	
	\textbf{Factor (C).} Como $\lambda/n \to 0$, se tiene
	$\left(1-\lambda/n\right)^{-k} \to 1$.
	
	\textbf{Factor (B).} Tomamos logaritmo y reescribimos la
	indeterminación $\infty\cdot 0$ como $0/0$:
	\[
	\ln\!\left(1-\frac{\lambda}{n}\right)^{\!n}
	= n\ln\!\left(1-\frac{\lambda}{n}\right)
	= \frac{\ln\!\left(1-\dfrac{\lambda}{n}\right)}{1/n}.
	\]
	Aplicando la regla de L'H\^{o}pital (tratando $n$ como variable
	continua):
	\[
	\lim_{n\to\infty}
	\frac{\ln(1-\lambda/n)}{1/n}
	= \lim_{n\to\infty}
	\frac{\dfrac{\lambda/n^2}{1-\lambda/n}}{-1/n^2}
	= \lim_{n\to\infty}
	\frac{-\lambda}{1-\lambda/n} = -\lambda.
	\]
	Por continuidad de la exponencial:
	$\left(1-\lambda/n\right)^n \to e^{-\lambda}$.
	
	Combinando los tres factores:
	\[
	\lim_{n\to\infty}
	\binom{n}{k}
	\left(\frac{\lambda}{n}\right)^k
	\left(1-\frac{\lambda}{n}\right)^{n-k}
	= \frac{\lambda^k}{k!}\cdot 1 \cdot e^{-\lambda}\cdot 1
	= \frac{\lambda^k e^{-\lambda}}{k!}. \qedhere
	\]
\end{proof}

La proposición anterior muestra que al tomar subintervalos
infinitesimales, la probabilidad de observar exactamente $k$ eventos
converge a $\lambda^k e^{-\lambda}/k!$. Esto motiva la siguiente
definición.

\begin{definicion}{Distribuci\'{o}n de Poisson}{dist_poisson}
	Una variable aleatoria discreta $X$ tiene \textbf{distribución de
		Poisson} con parámetro $\lambda > 0$ si su función de masa es
	\[
	p_X(k) = \frac{\lambda^k e^{-\lambda}}{k!},
	\qquad k = 0, 1, 2, \ldots
	\]
	Se escribe $X \sim \mathrm{Poisson}(\lambda)$.
\end{definicion}

\begin{observacion}
	La verificación de que $p_X$ es una función de masa válida es
	inmediata a partir de la serie de Taylor de $e^x$:
	\[
	\sum_{k=0}^{\infty} \frac{\lambda^k e^{-\lambda}}{k!}
	= e^{-\lambda} \sum_{k=0}^{\infty} \frac{\lambda^k}{k!}
	= e^{-\lambda} \cdot e^{\lambda} = 1.
	\]
	La serie de Taylor $e^x = \sum_{k=0}^{\infty} x^k/k!$ es el
	análogo continuo del teorema del binomio: así como la Binomial
	se apoya en $(p+q)^n = 1$, la Poisson se apoya en
	$e^{-\lambda}\cdot e^{\lambda} = 1$.
\end{observacion}

\begin{proposicion}{Momentos de la distribuci\'{o}n de Poisson}
	{momentos_poisson}
	Si $X \sim \mathrm{Poisson}(\lambda)$ con $\lambda > 0$, entonces
	\[
	E(X) = \lambda \qquad \text{y} \qquad \mathrm{Var}(X) = \lambda.
	\]
\end{proposicion}

\begin{proof}
	\textbf{Valor esperado.} El término $k=0$ no contribuye, luego:
	\[
	E(X)
	= \sum_{k=1}^{\infty} k\,\frac{e^{-\lambda}\lambda^k}{k!}
	= e^{-\lambda}\sum_{k=1}^{\infty} \frac{\lambda^k}{(k-1)!}.
	\]
	Sustituyendo $j = k-1$:
	\[
	= e^{-\lambda}\cdot\lambda \sum_{j=0}^{\infty}
	\frac{\lambda^j}{j!}
	= e^{-\lambda}\cdot\lambda\cdot e^{\lambda} = \lambda.
	\]
	
	\textbf{Varianza.} Usamos $\mathrm{Var}(X) = E(X^2)-(E(X))^2$
	y escribimos $E(X^2) = E(X(X-1)) + E(X)$. Para el primer término,
	los casos $k=0$ y $k=1$ no contribuyen:
	\[
	E(X(X-1))
	= \sum_{k=2}^{\infty} k(k-1)\,\frac{e^{-\lambda}\lambda^k}{k!}
	= e^{-\lambda}\sum_{k=2}^{\infty}\frac{\lambda^k}{(k-2)!}.
	\]
	Sustituyendo $j = k-2$:
	\[
	= e^{-\lambda}\cdot\lambda^2\sum_{j=0}^{\infty}
	\frac{\lambda^j}{j!}
	= \lambda^2.
	\]
	Por lo tanto $E(X^2) = \lambda^2 + \lambda$, y:
	\[
	\mathrm{Var}(X) = (\lambda^2+\lambda) - \lambda^2
	= \lambda. \qedhere
	\]
\end{proof}

\begin{observacion}
	La igualdad $E(X) = \mathrm{Var}(X) = \lambda$ es una propiedad
	característica de la Poisson: si en datos de conteo se observa que
	la media y la varianza muestral son aproximadamente iguales, la
	Poisson es un modelo candidato razonable. El parámetro $\lambda$
	tiene además una interpretación directa: es la tasa promedio de
	ocurrencia en el intervalo de referencia. Si los eventos ocurren
	a tasa $\lambda$ por hora, en un intervalo de $t$ horas se modela
	con $\mathrm{Poisson}(\lambda t)$.
\end{observacion}

Vale la pena contrastar las tres distribuciones del tema. Bernoulli
responde \emph{¿ocurrió el evento?} en un único ensayo. Binomial
responde \emph{¿cuántas veces ocurrió?} en $n$ ensayos fijos. Poisson
responde \emph{¿cuántos eventos ocurrieron?} en un intervalo continuo,
sin $n$ fijo, solo con una tasa $\lambda$ estable. La regla práctica:
si hay un $n$ contable y un $p$ identificable, Binomial; si los
eventos se distribuyen continuamente en el tiempo o el espacio con
tasa conocida, Poisson.

\begin{figure}[H]
	\centering
	\begin{tikzpicture}[scale=0.90]
		% Poisson(lambda=3) escalado por 10
		% k=0:  0.0498  -> 0.498
		% k=1:  0.1494  -> 1.494
		% k=2:  0.2240  -> 2.240
		% k=3:  0.2240  -> 2.240
		% k=4:  0.1680  -> 1.680
		% k=5:  0.1008  -> 1.008
		% k=6:  0.0504  -> 0.504
		% k=7:  0.0216  -> 0.216
		% k=8:  0.0081  -> 0.081
		% k=9:  0.0027  -> 0.027
		
		\def\bw{0.30}
		
		\draw[thick,->] (-0.5,0) -- (11.0,0) node[right] {$k$};
		\draw[thick,->] (0,-0.1) -- (0,3.0) node[above] {$p_X(k)$};
		
		\draw[gray,thin,dashed] (0,2.240) -- (10.5,2.240);
		\node[gray,left,font=\scriptsize] at (0,2.240) {$0.224$};
		
		\foreach \k/\h in {%
			0/0.498, 1/1.494, 2/2.240, 3/2.240, 4/1.680,%
			5/1.008, 6/0.504, 7/0.216, 8/0.081, 9/0.027}{
			\pgfmathsetmacro{\xp}{\k+1}
			\fill[verde!25] (\xp-\bw,0) rectangle (\xp+\bw,\h);
			\draw[verde!80!black,thick] (\xp-\bw,0) rectangle
			(\xp+\bw,\h);
			\node[below,font=\scriptsize] at (\xp,0) {$\k$};
		}
		
		\draw[rojo,thick,dashed] (4,0) -- (4,2.6);
		\node[rojo,above,font=\small] at (4,2.6) {$\lambda=3$};
		
		\node[font=\small,gray] at (7.0,2.85)
		{$X\sim\mathrm{Poisson}(3)$};
	\end{tikzpicture}
	\caption{Función de masa de $X\sim\mathrm{Poisson}(3)$. Media y
		varianza son ambas iguales a $\lambda=3$. La moda se ubica en
		$k=2$ y $k=3$ simultáneamente, ya que $\lambda$ no es entero más uno.}
	\label{fig:pmf_poisson}
\end{figure}

\begin{observacion}
	La distribución de Poisson es \textbf{discreta}: los valores que
	toma $X$ son $0, 1, 2, \ldots$, que son contables. Lo que es
	continuo no es la variable, sino el \emph{soporte} donde ocurren
	los eventos --- el tiempo, el espacio, una superficie. La variable
	cuenta cuántos eventos ocurrieron ahí, y ese conteo siempre es un
	número entero. En el ejemplo de mensajes a un servidor: el número
	de mensajes recibidos en una hora puede ser $0, 1, 2, 3, \ldots$
	--- nunca $2.7$. Lo continuo es el intervalo de tiempo en que cada
	mensaje puede llegar, no la cantidad de mensajes.
\end{observacion}

% ==========================================================
\chapter{Variable aleatoria continua}

\section{De lo discreto a lo continuo}


En el Tema~2 modelamos fenómenos cuyos resultados son contables:
número de éxitos, número de mensajes, número de fallas. La herramienta
central fue la función de masa $p_X$, que asigna probabilidad a cada
valor individual del recorrido. Pero hay fenómenos que no encajan en
ese marco: el tiempo de espera en una cola, la temperatura de una
ciudad, la altura de una persona. En estos casos el recorrido de $X$
no es un conjunto contable --- es un intervalo continuo de la recta
real, y la probabilidad de observar cualquier valor \emph{exacto} es
cero.

Para entender por qué, y cómo modelar estos fenómenos, conviene
partir de algo familiar: el histograma.

\subsection{Del histograma a la densidad}

Supongamos que medimos una variable continua --- digamos, el tiempo
de atención en ventanilla en segundos --- sobre una muestra de $n$
observaciones. Dividimos el rango de los datos en $k$ intervalos
$[a_0, a_1), [a_1, a_2), \ldots, [a_{k-1}, a_k)$ de ancho $\Delta x
= a_{i+1} - a_i$, y construimos un histograma. La altura de cada
barra puede definirse de dos maneras:

\begin{itemize}[leftmargin=2em]
	\item \textbf{Histograma de frecuencias:} la altura es el conteo
	de observaciones en el intervalo. El área no tiene interpretación
	probabilística directa.
	\item \textbf{Histograma de densidad:} la altura $d_i$ se define
	como la frecuencia relativa $fr_i$ dividida por el ancho del
	intervalo:
	\[
	d_i = \frac{fr_i}{\Delta x}.
	\]
	De este modo, el \textbf{área} de cada barra es igual a la
	frecuencia relativa:
	\[
	\text{área}_i = d_i \cdot \Delta x = fr_i,
	\]
	y la suma de todas las áreas es $1$. La probabilidad de que $X$
	caiga en el intervalo $[a_{i-1}, a_i)$ se aproxima por el área
	de la barra correspondiente.
\end{itemize}

El histograma de densidad es el punto de partida. Ahora preguntamos:
¿qué ocurre si refinamos la partición indefinidamente, haciendo
$\Delta x \to 0$?

\subsection{El límite: de la suma a la integral}

Con una partición de $k$ intervalos de ancho $\Delta x$, la
probabilidad de que $X$ caiga en $[a, b]$ se aproxima por la suma
de las áreas de las barras contenidas en ese intervalo:
\[
P(a \leq X \leq b) \approx \sum_{i} d_i \cdot \Delta x.
\]
Esta es una suma de Riemann. Al hacer $\Delta x \to 0$ --- es decir,
al refinar la partición hasta tener intervalos infinitesimales ---
las alturas $d_i$ convergen a una función continua $f(x)$ y la suma
converge a una integral:
\[
P(a \leq X \leq b) = \lim_{\Delta x \to 0}
\sum_{i} d_i \cdot \Delta x = \int_a^b f(x)\, dx.
\]
La función $f(x)$ que emerge en este límite es la
\textbf{función de densidad de probabilidad} de $X$. Ya no
asigna probabilidad a valores individuales --- asigna
\emph{densidad de probabilidad} a cada punto, de modo que la
probabilidad de un intervalo es el área bajo la curva $f$ en ese
intervalo.

\begin{observacion}
	Una consecuencia inmediata e importante: para una variable
	aleatoria continua,
	\[
	P(X = x) = \int_x^x f(t)\, dt = 0
	\]
	para todo $x \in \mathbb{R}$. La probabilidad de observar
	cualquier valor \emph{exacto} es cero. Esto no significa que
	el evento sea imposible --- significa que en un espacio continuo,
	la probabilidad se concentra en intervalos, no en puntos. En
	consecuencia, para una v.a.\ continua es irrelevante si los
	extremos del intervalo están incluidos o no:
	\[
	P(a \leq X \leq b) = P(a < X \leq b) =
	P(a \leq X < b) = P(a < X < b).
	\]
	Este resultado no tiene análogo en el mundo discreto y suele
	sorprender la primera vez que se encuentra.
\end{observacion}

\begin{figure}[H]
	\centering
	\begin{tikzpicture}[scale=0.95]
		
		% ── Panel izquierdo: histograma grueso ──────────────────
		\begin{scope}[shift={(0,0)}]
			\node[font=\small\bfseries, azul] at (3.2, 4.6)
			{Partici\'{o}n gruesa ($\Delta x$ grande)};
			
			\draw[thick,->] (-0.2,0) -- (5.5,0) node[right] {$x$};
			\draw[thick,->] (0,-0.1) -- (0,4.2) node[above] {$d$};
			
			\foreach \x/\h in {0/0.8, 1/2.2, 2/3.6, 3/2.4, 4/0.9}{
				\fill[azul!20] (\x,0) rectangle (\x+1,\h);
				\draw[azul,thick] (\x,0) rectangle (\x+1,\h);
			}
			\node[below, font=\scriptsize] at (0,0) {$a$};
			\node[below, font=\scriptsize] at (5,0) {$b$};
		\end{scope}
		
		% Flecha central
		\draw[->, very thick, azul!60] (5.9, 2.0) -- (7.1, 2.0)
		node[midway, above, font=\scriptsize] {$\Delta x \to 0$};
		
		% ── Panel derecho: curva límite ─────────────────────────
		\begin{scope}[shift={(7.5,0)}]
			\node[font=\small\bfseries, azul] at (3.0, 4.3)
			{L\'{i}mite ($\Delta x \to 0$)};
			
			\draw[thick,->] (-0.2,0) -- (5.5,0) node[right] {$x$};
			\draw[thick,->] (0,-0.1) -- (0,4.2) node[above] {$f(x)$};
			
			\fill[azul!15] (0,0)
			.. controls (0.5,0.3) and (1.0,1.8) .. (2.0,3.6)
			.. controls (3.0,1.8) and (4.0,0.4) .. (5.0,0.05)
			-- (5,0) -- (0,0);
			\draw[azul, very thick] (0,0)
			.. controls (0.5,0.3) and (1.0,1.8) .. (2.0,3.6)
			.. controls (3.0,1.8) and (4.0,0.4) .. (5.0,0.05);
			
			% Etiqueta del área — posicionada debajo de la curva
			\node[azul!70!black, font=\small] at (2.5, 0.55)
			{$\displaystyle\int_a^b f(x)\,dx = 1$};
			
			\node[below, font=\scriptsize] at (0,0) {$a$};
			\node[below, font=\scriptsize] at (5,0) {$b$};
		\end{scope}
		
	\end{tikzpicture}
	\caption{Al refinar la partici\'{o}n del histograma de densidad
		($\Delta x \to 0$), las alturas de las barras convergen a una
		funci\'{o}n continua $f(x)$. La probabilidad de un intervalo pasa
		de ser una suma de \'{a}reas de barras a ser el \'{a}rea bajo la curva.}
	\label{fig:histograma_densidad}
\end{figure}

\begin{definicion}{Variable aleatoria continua}{va_continua}
	Una variable aleatoria $X$ es \textbf{continua} si existe una
	función $f: \mathbb{R} \to [0, \infty)$ tal que para todo
	intervalo $[a, b]$:
	\[
	P(a \leq X \leq b) = \int_a^b f(x)\, dx.
	\]
\end{definicion}

\section{Funci\'{o}n de densidad de probabilidad}

El argumento del histograma nos dejó con una función $f(x)$ tal que
la probabilidad de un intervalo es el área bajo $f$ en ese intervalo.
Ahora formalizamos qué condiciones debe satisfacer $f$ para ser una
función de densidad válida --- el análogo continuo de los axiomas de
la función de masa.

\begin{definicion}{Funci\'{o}n de densidad de probabilidad}{fdp}
	Sea $X$ una variable aleatoria continua. Una función
	$f: \mathbb{R} \to \mathbb{R}$ es una \textbf{función de densidad
		de probabilidad} (fdp) de $X$ si satisface:
	\begin{enumerate}[label=\textbf{\roman*)}]
		\item $f(x) \geq 0$ para todo $x \in \mathbb{R}$.
		\item $\displaystyle\int_{-\infty}^{\infty} f(x)\,dx = 1$.
		\item Para todo $a \leq b$:
		\[
		P(a \leq X \leq b) = \int_a^b f(x)\,dx.
		\]
	\end{enumerate}
\end{definicion}

\begin{observacion}
	El paralelo con la función de masa discreta es exacto: la condición
	\textbf{i} es la no negatividad ($p_X(x) \geq 0$), la condición
	\textbf{ii} es la normalización ($\sum_x p_X(x) = 1$), y la
	condición \textbf{iii} es el cálculo de probabilidades (suma
	$\to$ integral). La única diferencia estructural es que la suma
	se reemplaza por la integral en todo lugar.
	
	Nótese además que $f(x)$ \textbf{no es una probabilidad} ---
	es una \emph{densidad}. Puede tomar valores mayores que $1$; lo
	que debe ser menor o igual a $1$ es el área bajo $f$ en cualquier
	intervalo, no el valor de $f$ en un punto.
\end{observacion}

\section{Funci\'{o}n de distribuci\'{o}n acumulada}

\begin{definicion}{Funci\'{o}n de distribuci\'{o}n acumulada
		(continua)}{fda_continua}
	Sea $X$ una variable aleatoria continua con fdp $f$. La
	\textbf{función de distribución acumulada} de $X$ es
	\[
	F(x) = P(X \leq x) = \int_{-\infty}^{x} f(t)\,dt,
	\qquad x \in \mathbb{R}.
	\]
\end{definicion}

\begin{proposicion}{Propiedades de la FDA continua}{props_fda_continua}
	Sea $X$ una variable aleatoria continua con fdp $f$ y FDA $F$.
	Entonces:
	\begin{enumerate}[label=\textbf{\roman*)}]
		\item $F$ es no decreciente: si $x_1 < x_2$ entonces
		$F(x_1) \leq F(x_2)$.
		\item $\lim_{x \to -\infty} F(x) = 0$.
		\item $\lim_{x \to \infty} F(x) = 1$.
		\item Para todo $a \leq b$:
		$P(a \leq X \leq b) = F(b) - F(a)$.
	\end{enumerate}
\end{proposicion}

\begin{proof}
	\textbf{i)} Si $x_1 < x_2$, entonces
	\[
	F(x_2) - F(x_1)
	= \int_{-\infty}^{x_2} f(t)\,dt -
	\int_{-\infty}^{x_1} f(t)\,dt
	= \int_{x_1}^{x_2} f(t)\,dt \geq 0,
	\]
	donde la última desigualdad se sigue de que $f(t) \geq 0$
	para todo $t$ (condición \textbf{i} de la
	Definición~\ref{def:fdp}).
	
	\medskip
	\textbf{ii)} Por la condición \textbf{ii} de la
	Definición~\ref{def:fdp}, $\int_{-\infty}^{\infty} f(t)\,dt = 1$,
	lo que implica en particular que $f$ es integrable en
	$\mathbb{R}$. Para una función integrable, la integral sobre
	un intervalo $(-\infty, x]$ tiende a $0$ cuando $x \to -\infty$:
	\[
	\lim_{x \to -\infty} F(x)
	= \lim_{x \to -\infty} \int_{-\infty}^{x} f(t)\,dt = 0.
	\]
	
	\medskip
	\textbf{iii)} Directamente de la condición \textbf{ii} de la
	Definición~\ref{def:fdp}:
	\[
	\lim_{x \to \infty} F(x)
	= \lim_{x \to \infty} \int_{-\infty}^{x} f(t)\,dt
	= \int_{-\infty}^{\infty} f(t)\,dt = 1.
	\]
	
	\medskip
	\textbf{iv)} Por aditividad de la integral:
	\[
	\int_{-\infty}^{b} f(t)\,dt
	= \int_{-\infty}^{a} f(t)\,dt + \int_{a}^{b} f(t)\,dt,
	\]
	de modo que
	\[
	P(a \leq X \leq b)
	= \int_a^b f(t)\,dt
	= F(b) - F(a). \qedhere
	\]
\end{proof}

La diferencia central con el caso discreto es que ahora $F$ es
continua y diferenciable, lo que permite recuperar la densidad a
partir de la FDA vía el Teorema Fundamental del Cálculo.

\begin{proposicion}{Relaci\'{o}n entre $F$ y $f$}{fda_fdp}
	Sea $X$ una variable aleatoria continua con fdp $f$ y FDA $F$.
	En todo punto donde $f$ es continua,
	\[
	F'(x) = f(x).
	\]
\end{proposicion}

\begin{proposicion}{Relaci\'{o}n entre $F$ y $f$}{fda_fdp}
	Sea $X$ una variable aleatoria continua con fdp $f$ y FDA $F$.
	En todo punto donde $f$ es continua,
	\[
	F'(x) = f(x).
	\]
\end{proposicion}

\begin{proof}
	Recordamos primero el \textbf{Teorema Fundamental del C\'{a}lculo}
	(TFC): si $g$ es una función continua en $[a,b]$ y se define
	\[
	G(x) = \int_a^x g(t)\,dt, \qquad x \in [a,b],
	\]
	entonces $G$ es diferenciable en $(a,b)$ y $G'(x) = g(x)$.
	Es decir, la derivada de una integral con límite superior
	variable recupera el integrando evaluado en ese límite.
	
	Aplicamos este resultado a nuestra situación. Por definición,
	\[
	F(x) = \int_{-\infty}^{x} f(t)\,dt.
	\]
	En todo punto donde $f$ es continua, el TFC garantiza que $F$
	es diferenciable y que
	\[
	F'(x) = f(x). \qedhere
	\]
\end{proof}

\begin{observacion}
	Esta proposición tiene una consecuencia práctica importante:
	si conocemos la FDA $F$, podemos recuperar la fdp simplemente
	derivando. Y recíprocamente, si conocemos la fdp $f$, podemos
	obtener la FDA integrando. Las dos funciones contienen la misma
	información sobre $X$ --- solo difieren en la forma de
	presentarla.
\end{observacion}


Esta proposición es la razón por la que $f$ se llama
\emph{densidad}: así como la densidad física de un objeto es la
derivada de su masa acumulada respecto a la posición, la fdp es
la derivada de la probabilidad acumulada respecto a $x$. La FDA
acumula; la fdp describe cómo se distribuye esa acumulación.


\section{Valor esperado y varianza}

Los conceptos de valor esperado y varianza se trasladan al caso
continuo reemplazando la suma por la integral --- la misma sustitución
que hicimos en la fdp.

En el caso discreto, el valor esperado era una suma ponderada
$\sum_x x\, p_X(x)$: cada valor posible se pesa por su
probabilidad. En el caso continuo, la probabilidad de que $X$
tome valores en un entorno infinitesimal de $x$ es $f(x)\,dx$,
de modo que la suma se convierte naturalmente en una integral.
La definición es la única extensión consistente del caso discreto.

\begin{definicion}{Valor esperado continuo}{ev_continuo}
	Sea $X$ una variable aleatoria continua con fdp $f$. El
	\textbf{valor esperado} de $X$ es
	\[
	E(X) = \int_{-\infty}^{\infty} x\, f(x)\,dx,
	\]
	siempre que la integral converja absolutamente.
\end{definicion}

De manera análoga, la varianza se extiende al caso continuo
reemplazando la suma por la integral. La interpretación es
idéntica a la del caso discreto: mide la dispersión promedio
de $X$ alrededor de su media, ponderada ahora por la densidad
en lugar de por la función de masa.

\begin{definicion}{Varianza continua}{var_continua}
	Sea $X$ una variable aleatoria continua con fdp $f$ y
	$\mu = E(X)$. La \textbf{varianza} de $X$ es
	\[
	\mathrm{Var}(X) = E\!\left[(X-\mu)^2\right]
	= \int_{-\infty}^{\infty} (x-\mu)^2\, f(x)\,dx,
	\]
	y la \textbf{desviación estándar} es
	$\sigma_X = \sqrt{\mathrm{Var}(X)}$.
\end{definicion}

\begin{observacion}
	La fórmula alternativa $\mathrm{Var}(X) = E(X^2) - (E(X))^2$
	se mantiene válida en el caso continuo, con la misma demostración
	algebraica que en el caso discreto
	(Proposición~\ref{prop:var_alt}). Las propiedades de
	linealidad del valor esperado y las propiedades de la varianza
	también se transfieren sin cambio.
\end{observacion}

\section{Distribuci\'{o}n Uniforme continua}

La distribución uniforme continua modela una variable aleatoria
que puede tomar cualquier valor en un intervalo $[a,b]$ con la
misma probabilidad relativa. Esto impone una exigencia concreta
sobre la fdp: debe ser \textbf{constante} en $[a,b]$, es decir
$f(x) = c$ para alguna constante $c > 0$, y cero fuera del
intervalo. La condición de normalización determina $c$
unívocamente:
\[
\int_{-\infty}^{\infty} f(x)\,dx
= \int_a^b c\,dx = c(b-a) = 1
\implies c = \frac{1}{b-a}.
\]
La fdp queda así completamente determinada por la sola exigencia
de equiprobabilidad.

\begin{definicion}{Distribuci\'{o}n Uniforme continua}{dist_uniforme}
	Una variable aleatoria $X$ tiene \textbf{distribución uniforme
		continua} en el intervalo $[a,b]$, con $a < b$, si su función
	de densidad de probabilidad es
	\[
	f(x) = \begin{cases}
		\dfrac{1}{b-a} & \text{si } a \leq x \leq b, \\[6pt]
		0              & \text{en otro caso.}
	\end{cases}
	\]
	Se escribe $X \sim U(a,b)$.
\end{definicion}

\begin{proposicion}{FDA, valor esperado y varianza de la Uniforme}
	{momentos_uniforme}
	Si $X \sim U(a,b)$, entonces:
	
	\medskip
	\noindent\textbf{i)} La función de distribución acumulada es
	\[
	F(x) = \begin{cases}
		0                      & \text{si } x < a,        \\[4pt]
		\dfrac{x-a}{b-a}       & \text{si } a \leq x < b, \\[6pt]
		1                      & \text{si } x \geq b.
	\end{cases}
	\]
	
	\medskip
	\noindent\textbf{ii)} $E(X) = \dfrac{a+b}{2}$.
	
	\medskip
	\noindent\textbf{iii)} $\mathrm{Var}(X) = \dfrac{(b-a)^2}{12}$.
\end{proposicion}

\begin{proof}
	\textbf{i)} Para $x < a$: $F(x) = \int_{-\infty}^x 0\,dt = 0$.
	
	Para $a \leq x < b$:
	\[
	F(x) = \int_{-\infty}^{x} f(t)\,dt
	= \int_a^x \frac{1}{b-a}\,dt
	= \frac{x-a}{b-a}.
	\]
	
	Para $x \geq b$: $F(x) = \int_{-\infty}^{\infty} f(t)\,dt = 1$.
	
	\medskip
	\textbf{ii)} Aplicando la definición de valor esperado continuo:
	\[
	E(X) = \int_{-\infty}^{\infty} x\,f(x)\,dx
	= \int_a^b \frac{x}{b-a}\,dx
	= \frac{1}{b-a} \cdot \frac{x^2}{2}\Bigg|_a^b
	= \frac{b^2 - a^2}{2(b-a)}
	= \frac{(b+a)(b-a)}{2(b-a)}
	= \frac{a+b}{2}.
	\]
	
	\medskip
	\textbf{iii)} Usamos la fórmula alternativa
	$\mathrm{Var}(X) = E(X^2) - (E(X))^2$. Calculamos $E(X^2)$:
	\[
	E(X^2) = \int_a^b \frac{x^2}{b-a}\,dx
	= \frac{1}{b-a} \cdot \frac{x^3}{3}\Bigg|_a^b
	= \frac{b^3 - a^3}{3(b-a)}
	= \frac{(b^2+ab+a^2)(b-a)}{3(b-a)}
	= \frac{a^2+ab+b^2}{3}.
	\]
	Por lo tanto:
	\[
	\mathrm{Var}(X)
	= \frac{a^2+ab+b^2}{3} - \left(\frac{a+b}{2}\right)^2
	= \frac{a^2+ab+b^2}{3} - \frac{a^2+2ab+b^2}{4}.
	\]
Llevando a denominador com\'{u}n $12$:
\begin{align*}
	\mathrm{Var}(X)
	&= \frac{4(a^2+ab+b^2) - 3(a^2+2ab+b^2)}{12} \\[6pt]
	&= \frac{4a^2+4ab+4b^2-3a^2-6ab-3b^2}{12} \\[6pt]
	&= \frac{a^2-2ab+b^2}{12} \\[6pt]
	&= \frac{(b-a)^2}{12}. \qedhere
\end{align*}
\end{proof}

\begin{observacion}
	El valor esperado $E(X) = (a+b)/2$ es el punto medio del
	intervalo --- exactamente lo que la intuición sugiere para una
	distribución completamente simétrica. La varianza
	$(b-a)^2/12$ crece con el cuadrado del largo del intervalo:
	mientras más amplio es $[a,b]$, más dispersa es $X$.
\end{observacion}

\begin{figure}[H]
	\centering
	\begin{tikzpicture}[scale=0.95]
		
		% ── Panel izquierdo: fdp ────────────────────────────────
		\begin{scope}[shift={(0,0)}]
			\node[font=\small\bfseries, azul] at (3.0, 3.0)
			{Funci\'{o}n de densidad $f(x)$};
			
			\draw[thick,->] (-0.5,0) -- (6.0,0) node[right] {$x$};
			\draw[thick,->] (0,-0.2) -- (0,2.8) node[above] {};
			
			% Rectángulo de la fdp
			\fill[azul!20] (1.5,0) rectangle (4.5,2.0);
			\draw[azul,thick] (1.5,0) -- (1.5,2.0) -- (4.5,2.0)
			-- (4.5,0);
			\draw[azul,thick] (-0.5,0) -- (1.5,0);
			\draw[azul,thick] (4.5,0) -- (6.0,0);
			
			% Etiqueta altura
			\draw[gray,dashed,thin] (0,2.0) -- (1.5,2.0);
			\node[left, font=\small] at (0,2.0)
			{$\frac{1}{b-a}$};
			
			% Etiquetas eje x
			\node[below, font=\small] at (1.5,0) {$a$};
			\node[below, font=\small] at (4.5,0) {$b$};
			
			% Media
			\draw[rojo,dashed,thick] (3.0,0) -- (3.0,2.0);
			\node[rojo,below,font=\small] at (3.0,0)
			{$\frac{a+b}{2}$};
		\end{scope}
		
		% ── Panel derecho: FDA ──────────────────────────────────
		\begin{scope}[shift={(7.5,0)}]
			\node[font=\small\bfseries, azul] at (3.0, 3.0)
			{FDA $F(x)$};
			
			\draw[thick,->] (-0.5,0) -- (6.0,0) node[right] {$x$};
			\draw[thick,->] (0,-0.2) -- (0,2.8) node[above] {};
			
			% F(x) = 0 para x < a
			\draw[azul,thick] (-0.5,0) -- (1.5,0);
			
			% F(x) = (x-a)/(b-a) rampa lineal
			\fill[azul!20] (1.5,0) -- (4.5,2.0) -- (4.5,0) -- cycle;
			\draw[azul,thick] (1.5,0) -- (4.5,2.0);
			
			% F(x) = 1 para x >= b
			\draw[azul,thick] (4.5,2.0) -- (6.0,2.0);
			
			% Marca en y=1
			\draw[gray,dashed,thin] (0,2.0) -- (4.5,2.0);
			\node[left,font=\small] at (0,2.0) {$1$};
			
			% Etiquetas eje x
			\node[below,font=\small] at (1.5,0) {$a$};
			\node[below,font=\small] at (4.5,0) {$b$};
		\end{scope}
		
	\end{tikzpicture}
	\caption{Función de densidad y FDA de $X \sim U(a,b)$. La fdp es
		un rectángulo de altura $1/(b-a)$ y la FDA es una rampa lineal
		entre $a$ y $b$. La línea roja marca la media $E(X) = (a+b)/2$.}
	\label{fig:uniforme_continua}
\end{figure}

\section{Distribuci\'{o}n Exponencial}

En la distribución de Poisson modelamos el número de eventos
que ocurren en un intervalo de tiempo fijo, asumiendo una tasa
promedio $\lambda$ constante. Pero hay una pregunta dual que
Poisson no responde directamente: ¿cuánto tiempo transcurre
entre un evento y el siguiente? La distribución exponencial
es la respuesta --- y no por coincidencia usa el mismo parámetro
$\lambda$.

\subsubsection{Derivaci\'{o}n de la fdp}

Sea $X$ el tiempo de espera hasta el próximo evento en un proceso
de Poisson con tasa $\lambda$. Queremos encontrar la distribución
de $X$. El punto de partida es calcular $P(X > x)$ --- la
probabilidad de que el próximo evento ocurra \emph{después} de
$x$ unidades de tiempo.

El evento $\{X > x\}$ es equivalente a que no ocurra ningún
evento en el intervalo $[0, x]$. Si $Y$ denota el número de
eventos en $[0,x]$, entonces $Y \sim \mathrm{Poisson}(\lambda x)$
--- la tasa es $\lambda$ por unidad de tiempo, luego en $x$
unidades la tasa es $\lambda x$. Por lo tanto:
\[
P(X > x) = P(Y = 0)
= \frac{(\lambda x)^0 e^{-\lambda x}}{0!}
= e^{-\lambda x}.
\]
De aquí obtenemos la FDA de $X$ directamente:
\[
F(x) = P(X \leq x) = 1 - P(X > x) = 1 - e^{-\lambda x},
\qquad x \geq 0.
\]
Aplicando la Proposición~\ref{prop:fda_fdp}, la fdp se obtiene
derivando la FDA:
\[
f(x) = F'(x) = \lambda e^{-\lambda x}, \qquad x > 0.
\]
La fdp no emerge de una postulación arbitraria --- es la
consecuencia directa de la estructura del proceso de Poisson
subyacente.

\begin{definicion}{Distribuci\'{o}n Exponencial}{dist_exponencial}
	Una variable aleatoria continua $X$ tiene
	\textbf{distribución exponencial} con parámetro $\lambda > 0$
	si su función de densidad de probabilidad es
	\[
	f(x) = \begin{cases}
		\lambda e^{-\lambda x} & \text{si } x > 0, \\
		0                      & \text{en otro caso.}
	\end{cases}
	\]
	Se escribe $X \sim \mathrm{Exp}(\lambda)$.
\end{definicion}

\begin{observacion}
	La verificación de que $f$ es una fdp válida es inmediata:
	$f(x) \geq 0$ para todo $x$, y
	\[
	\int_{-\infty}^{\infty} f(x)\,dx
	= \int_0^{\infty} \lambda e^{-\lambda x}\,dx
	= \lambda \cdot \left[-\frac{1}{\lambda}e^{-\lambda x}
	\right]_0^{\infty}
	= \left[-e^{-\lambda x}\right]_0^{\infty}
	= 0 - (-1) = 1.
	\]
\end{observacion}

\begin{proposicion}{FDA, valor esperado y varianza de la
		Exponencial}{momentos_exp}
	Si $X \sim \mathrm{Exp}(\lambda)$ con $\lambda > 0$, entonces:
	
	\medskip
	\noindent\textbf{i)} La función de distribución acumulada es
	\[
	F(x) = \begin{cases}
		0                    & \text{si } x < 0,    \\[4pt]
		1 - e^{-\lambda x}   & \text{si } x \geq 0.
	\end{cases}
	\]
	
	\medskip
	\noindent\textbf{ii)} $E(X) = \dfrac{1}{\lambda}$.
	
	\medskip
	\noindent\textbf{iii)} $\mathrm{Var}(X) = \dfrac{1}{\lambda^2}$.
\end{proposicion}

\begin{proof}
	\textbf{i)} Ya fue derivada en la sección anterior.
	
	\medskip
	\textbf{ii)} Aplicamos la definición de valor esperado:
	\[
	E(X) = \int_0^{\infty} x \cdot \lambda e^{-\lambda x}\,dx
	= \lambda \int_0^{\infty} x e^{-\lambda x}\,dx.
	\]
	Integramos por partes con $u = x$, $dv = e^{-\lambda x}\,dx$,
	luego $du = dx$, $v = -\frac{1}{\lambda}e^{-\lambda x}$:
	\[
	\int_0^{\infty} x e^{-\lambda x}\,dx
	= \left[-\frac{x}{\lambda}e^{-\lambda x}\right]_0^{\infty}
	+ \int_0^{\infty} \frac{1}{\lambda}e^{-\lambda x}\,dx.
	\]
	El primer término se anula en ambos extremos: en $x = 0$ da
	$0$, y cuando $x \to \infty$, $xe^{-\lambda x} \to 0$ porque
	la exponencial domina al polinomio. El segundo término:
	\[
	\int_0^{\infty} \frac{1}{\lambda}e^{-\lambda x}\,dx
	= \frac{1}{\lambda}\left[-\frac{1}{\lambda}
	e^{-\lambda x}\right]_0^{\infty}
	= \frac{1}{\lambda^2}.
	\]
	Por lo tanto:
	\[
	E(X) = \lambda \cdot \frac{1}{\lambda^2} = \frac{1}{\lambda}.
	\]
	
	\medskip
	\textbf{iii)} Usamos la fórmula alternativa
	$\mathrm{Var}(X) = E(X^2) - (E(X))^2$. Calculamos $E(X^2)$
	integrando por partes dos veces con $u = x^2$,
	$dv = e^{-\lambda x}\,dx$:
	\[
	E(X^2) = \lambda\int_0^{\infty} x^2 e^{-\lambda x}\,dx
	= \lambda\left(
	\left[-\frac{x^2}{\lambda}e^{-\lambda x}
	\right]_0^{\infty}
	+ \frac{2}{\lambda}\int_0^{\infty}
	x e^{-\lambda x}\,dx\right).
	\]
	El término evaluado se anula por el mismo argumento que antes.
	El integral restante ya lo calculamos en \textbf{ii)}:
	$\int_0^{\infty} x e^{-\lambda x}\,dx = 1/\lambda^2$. Entonces:
	\[
	E(X^2) = \lambda \cdot \frac{2}{\lambda} \cdot
	\frac{1}{\lambda^2} = \frac{2}{\lambda^2}.
	\]
	Por lo tanto:
	\[
	\mathrm{Var}(X) = \frac{2}{\lambda^2} -
	\left(\frac{1}{\lambda}\right)^2
	= \frac{2}{\lambda^2} - \frac{1}{\lambda^2}
	= \frac{1}{\lambda^2}. \qedhere
	\]
\end{proof}

\begin{observacion}
	El valor esperado $E(X) = 1/\lambda$ tiene una interpretación
	directa: si los eventos ocurren a tasa $\lambda$ por unidad de
	tiempo, el tiempo promedio de espera entre eventos es
	$1/\lambda$. Esto cierra la conexión con Poisson: si
	$\lambda = 5$ eventos por hora, el tiempo promedio entre
	eventos es $1/5$ de hora, es decir, 12 minutos.
\end{observacion}

La distribución exponencial tiene una propiedad que la hace única
entre las distribuciones continuas y que resulta fundamental en
aplicaciones de confiabilidad y teoría de colas.

\begin{proposicion}{Propiedad de falta de memoria}
	{falta_memoria}
	Si $X \sim \mathrm{Exp}(\lambda)$, entonces para todo
	$s, t > 0$:
	\[
	P(X > s + t \mid X > s) = P(X > t).
	\]
\end{proposicion}

\begin{proof}
	Aplicando la definición de probabilidad condicional y usando
	que $P(X > x) = e^{-\lambda x}$:
	\[
	P(X > s+t \mid X > s)
	= \frac{P(X > s+t,\; X > s)}{P(X > s)}
	= \frac{P(X > s+t)}{P(X > s)}
	= \frac{e^{-\lambda(s+t)}}{e^{-\lambda s}}
	= e^{-\lambda t}
	= P(X > t). \qedhere
	\]
\end{proof}

\begin{observacion}
	La propiedad de falta de memoria dice que el tiempo de espera
	restante no depende del tiempo ya esperado. Si $X$ es la vida
	útil de un componente y el componente ha sobrevivido $s$ horas,
	la probabilidad de que sobreviva $t$ horas más es exactamente
	la misma que si fuera nuevo. El componente no \emph{recuerda}
	cuánto tiempo lleva funcionando --- no se desgasta en sentido
	probabilístico. Esta propiedad caracteriza a la distribución
	exponencial: es la \textbf{única} distribución continua con
	falta de memoria.
\end{observacion}

\begin{figure}[H]
	\centering
	\begin{tikzpicture}[scale=0.95]
		
		% ── Panel izquierdo: fdp ────────────────────────────────
		\begin{scope}[shift={(0,0)}]
			\node[font=\small\bfseries, azul] at (3.0, 3.8)
			{Funci\'{o}n de densidad $f(x)$};
			
			\draw[thick,->] (-0.3,0) -- (6.2,0) node[right] {$x$};
			\draw[thick,->] (0,-0.2) -- (0,3.6) node[above] {};
			
		% lambda = 0.5
		\draw[azul!50, thick, domain=0:6, samples=100]
		plot (\x, {1.5*0.5*exp(-0.5*\x)});
		\node[azul!60, font=\scriptsize] at (4.5, 0.6)
		{$\lambda = 0.5$};
		
		% lambda = 1
		\draw[azul, thick, domain=0:6, samples=100]
		plot (\x, {1.5*1*exp(-1*\x)});
		\node[azul, font=\scriptsize] at (2.5, 0.9)
		{$\lambda = 1$};
		
		% lambda = 2
		\draw[azul!80!black, very thick, domain=0:6, samples=100]
		plot (\x, {1.5*2*exp(-2*\x)});
		\node[azul!80!black, font=\scriptsize] at (1.2, 1.8)
		{$\lambda = 2$};
		
		% Puntos en x=0
		\foreach \lam/\col in {0.5/azul!50, 1/azul,
			2/azul!80!black}{
			\fill[\col] (0, {1.5*\lam}) circle (2pt);
		}
		\end{scope}
		
		% ── Panel derecho: FDA ──────────────────────────────────
		\begin{scope}[shift={(7.5,0)}]
			\node[font=\small\bfseries, azul] at (3.0, 3.8)
			{FDA $F(x)$};
			
			\draw[thick,->] (-0.3,0) -- (6.2,0) node[right] {$x$};
			\draw[thick,->] (0,-0.2) -- (0,3.6) node[above] {};
			
			% Línea y=3 (representa F=1)
			\draw[gray, thin, dashed] (0,3.0) -- (6.0,3.0);
			\node[gray, left, font=\scriptsize] at (0,3.0) {$1$};
			
			% lambda = 0.5
			\draw[azul!50, thick, domain=0:6, samples=100]
			plot (\x, {3*(1-exp(-0.5*\x))});
			\node[azul!60, font=\scriptsize] at (5.0, 1.8)
			{$\lambda = 0.5$};
			
			% lambda = 1
			\draw[azul, thick, domain=0:6, samples=100]
			plot (\x, {3*(1-exp(-1*\x))});
			\node[azul, font=\scriptsize] at (3.5, 2.4)
			{$\lambda = 1$};
			
			% lambda = 2
			\draw[azul!80!black, very thick, domain=0:6, samples=100]
			plot (\x, {3*(1-exp(-2*\x))});
			\node[azul!80!black, font=\scriptsize] at (2.0, 2.8)
			{$\lambda = 2$};
			
			% F(x) = 0 para x < 0
			\draw[azul, thick] (-0.3,0) -- (0,0);
		\end{scope}
		
	\end{tikzpicture}
	\caption{Función de densidad y FDA de $X \sim \mathrm{Exp}(\lambda)$
		para distintos valores de $\lambda$. A mayor $\lambda$, la densidad
		cae más rápidamente y la FDA satura antes --- el tiempo de espera
		promedio $1/\lambda$ es menor.}
	\label{fig:exponencial}
\end{figure}

\section{Distribuci\'{o}n Normal}

La distribución Normal es la más importante de toda la
estadística. Aparece en fenómenos tan distintos como las alturas
de una población, los errores de medición en un experimento
físico, o las fluctuaciones de precios en un mercado financiero.
Su ubicuidad no es accidental --- tiene una justificación
matemática profunda que exploraremos en la subsección siguiente.
Por ahora la introducimos formalmente.

\begin{definicion}{Distribuci\'{o}n Normal}{dist_normal}
	Una variable aleatoria continua $X$ tiene
	\textbf{distribución Normal} con parámetros $\mu \in \mathbb{R}$
	y $\sigma^2 > 0$ si su función de densidad de probabilidad es
	\[
	f(x) = \frac{1}{\sqrt{2\pi}\,\sigma}
	\exp\!\left(-\frac{(x-\mu)^2}{2\sigma^2}\right),
	\qquad x \in \mathbb{R}.
	\]
Se escribe $X \sim N(\mu, \sigma^2)$.
\end{definicion}

\begin{figure}[H]
	\centering
	\begin{tikzpicture}[scale=1.0]
		
		\draw[thick,->] (-4.5,0) -- (4.5,0) node[right] {$x$};
		\draw[thick,->] (0,-0.2) -- (0,4.2) node[above] {$f(x)$};
		
		% Curva Normal(0,1) escalada -- 1/sqrt(2pi) = 0.3989
		\fill[azul!15, domain=-4:4, samples=200]
		plot (\x, {3.5*0.3989*exp(-\x*\x/2)})
		-- (4,0) -- (-4,0) -- cycle;
		\draw[azul, very thick, domain=-4:4, samples=200]
		plot (\x, {3.5*0.3989*exp(-\x*\x/2)});
		
	    % Línea en la media
	\draw[rojo, dashed, thick] (0,0) -- (0, {3.5*0.3989});
	\node[rojo, right, font=\small] at (0.2, {3.5*0.3989+0.4})
	{$\dfrac{1}{\sqrt{2\pi}\,\sigma}$};
	\draw[rojo, ->] (0.2, {3.5*0.3989+0.35})
	-- (0.05, {3.5*0.3989+0.05});
		
		% Marcas en mu +/- sigma
		\draw[verde!70!black, dashed, thick]
		(1.0,0) -- (1.0, {3.5*0.3989*exp(-0.5)});
		\draw[verde!70!black, dashed, thick]
		(-1.0,0) -- (-1.0, {3.5*0.3989*exp(-0.5)});
		\node[verde!70!black, below, font=\small] at (1.0,0)
		{$\mu+\sigma$};
		\node[verde!70!black, below, font=\small] at (-1.0,0)
		{$\mu-\sigma$};
		
	    % Simetría — nota lateral en lugar de flecha
	    \node[gray, font=\small, align=center] at (3.5, 3.2)
	    {sim\'{e}trica\\respecto a $\mu$};
	    \draw[->, gray, thin] (3.5, 2.9) -- (0.3, 2.1);
		
	\end{tikzpicture}
	\caption{Funci\'{o}n de densidad de $X \sim N(\mu, \sigma^2)$. La
		curva es sim\'{e}trica respecto a $\mu$, alcanza su m\'{a}ximo en
		$x = \mu$ con altura $1/(\sqrt{2\pi}\,\sigma)$, y tiene puntos
		de inflexi\'{o}n en $\mu \pm \sigma$.}
	\label{fig:normal_intro}
\end{figure}

\begin{observacion}
	La fdp Normal tiene tres propiedades geométricas que se
	desprenden directamente de su fórmula:
	\begin{enumerate}[label=\textbf{\roman*)}]
		\item Es \textbf{simétrica} respecto a $x = \mu$ --- el
		exponente $(x-\mu)^2$ es el mismo para $\mu + t$ y
		$\mu - t$.
		\item Alcanza su \textbf{máximo} en $x = \mu$, donde
		$f(\mu) = 1/(\sqrt{2\pi}\,\sigma)$.
		\item Las \textbf{colas} decaen a cero cuando
		$x \to \pm\infty$, más rápidamente cuanto menor es
		$\sigma$.
	\end{enumerate}
	El parámetro $\mu$ controla la \emph{ubicación} de la curva
	y $\sigma$ controla su \emph{dispersión}: una $\sigma$ pequeña
	produce una campana alta y estrecha; una $\sigma$ grande
	produce una campana baja y ancha.
\end{observacion}

\subsection{\textquestiondown Por qu\'{e} $\sqrt{2\pi}$?}

La fórmula de la fdp Normal no es arbitraria. Cada uno de sus
componentes --- la exponencial cuadrática, la constante
$\sqrt{2\pi}$ --- tiene una justificación matemática precisa.
Para entenderla es útil remontarse a su origen histórico.

A mediados del siglo XIX, tres matemáticos llegaron a la misma
función por caminos completamente distintos: \textbf{Carl
	Friedrich Gauss} desde la teoría de errores de medición
astronómica, \textbf{Pierre-Simon Laplace} desde la combinatoria
y los límites de distribuciones binomiales, y \textbf{James
	Clerk Maxwell} --- el mismo de las ecuaciones del
electromagnetismo --- desde la física cinética de los gases.
Que tres rutas tan distintas converjan a la misma función es la
primera señal de que algo profundo está ocurriendo.

\subsubsection{El argumento de Maxwell}

En 1860, Maxwell quería describir cómo se distribuyen las
velocidades de las moléculas de un gas. Postuló dos condiciones
que parecen muy razonables:

\begin{enumerate}[label=\textbf{\roman*)}]
	\item \textbf{Isotropía:} la distribución de velocidades en
	cada dirección ($x$, $y$, $z$) es la misma función $f$.
	No hay ninguna dirección privilegiada.
	\item \textbf{Independencia direccional:} las velocidades en
	distintas direcciones son independientes entre sí.
\end{enumerate}

Si $v^2 = v_x^2 + v_y^2 + v_z^2$ es la rapidez total, la
probabilidad conjunta debe depender solo de $v^2$ --- no de
la dirección. Esto impone la condición funcional:
\[
f(v_x)\cdot f(v_y)\cdot f(v_z) = g(v_x^2 + v_y^2 + v_z^2).
\]
Tomando logaritmo:
\[
\ln f(v_x) + \ln f(v_y) + \ln f(v_z)
= \ln g(v_x^2 + v_y^2 + v_z^2).
\]
La única familia de funciones que satisface esta ecuación
funcional es
\[
f(t) = c\, e^{-\alpha t^2}, \qquad \alpha > 0,
\]
para alguna constante $c > 0$. La demostración es la siguiente:
diferenciando la ecuación funcional respecto a $v_x$ con $v_y$
y $v_z$ fijos, se obtiene que $(\ln f)'(v_x)$ debe ser
proporcional a $v_x$, lo que implica $\ln f(t) = -\alpha t^2 +
\text{cte}$, y por lo tanto $f(t) = ce^{-\alpha t^2}$.

\begin{figure}[H]
	\centering
	\begin{tikzpicture}[scale=1.0]
		
		% ── Panel izquierdo: isotropía ──────────────────────────
		\begin{scope}[shift={(0,0)}]
			\node[font=\small\bfseries, azul] at (0, 3.2)
			{Isotrop\'{\i}a};
			\node[font=\scriptsize, gray] at (0, 2.75)
			{ninguna direcci\'{o}n privilegiada};
			
			% Círculo de referencia
			\draw[gray, dashed, thin, opacity=0.5] (0,0) circle (2);
			
			% Vectores en 8 direcciones
			\foreach \ang in {0,45,...,315}{
				\draw[azul, thick, ->] (0,0) --
				(\ang:2);
			}
			
			% Molécula central
			\fill[azul!80!black] (0,0) circle (3pt);
		\end{scope}
		
		% Separador vertical
		\draw[gray, thin, opacity=0.3] (3.5,-2.8) -- (3.5,3.0);
		
% ── Panel derecho: descomposición ───────────────────────
\begin{scope}[shift={(7.0,0)}]
	\node[font=\small\bfseries, azul] at (1.0, 3.2)
	{Descomposici\'{o}n};
	\node[font=\scriptsize, gray] at (1.0, 2.78)
	{depende solo de la rapidez};
	
	% Ejes — el eje y sube solo hasta 2.0 para no chocar
	\draw[thick, ->] (-0.3,-2) -- (3.0,-2)
	node[right] {$v_x$};
	\draw[thick, ->] (-0.3,-2) -- (-0.3,2.0)
	node[above] {$v_y$};
	
	% Arco de rapidez constante — centrado en el origen,
	% pasando por la punta del vector
	\draw[rojo, dashed, thin, opacity=0.6]
	(-0.3,-2) ++(0:2.95) arc (0:50:2.95);
	
	% Vector velocidad (de origen a la punta)
	\draw[rojo, very thick, ->] (-0.3,-2) --
	(1.6, 0.3) node[midway, above left, font=\small]
	{$\vec{v}$};
	
	% Componentes punteadas
	\draw[azul, thick, dashed] (-0.3,-2) -- (1.6,-2);
	\draw[azul, thick, dashed] (1.6,-2) -- (1.6,0.3);
	
	% Etiquetas componentes
	\node[azul, below, font=\small] at (0.65,-2) {$v_x$};
	\node[azul, right, font=\small] at (1.6,-0.85) {$v_y$};
	
	% Origen
	\fill[rojo] (-0.3,-2) circle (2pt);
\end{scope}
		
	\end{tikzpicture}
	\caption{El argumento de Maxwell. \textbf{Izquierda:} la isotropía
		exige que la distribución de velocidades sea la misma en todas las
		direcciones --- no hay ninguna dirección privilegiada.
		\textbf{Derecha:} al descomponer la velocidad en componentes
		$v_x$ y $v_y$, la condición de que la distribución conjunta dependa
		solo de la rapidez $v = \sqrt{v_x^2 + v_y^2}$ --- es decir, que sea
		constante sobre cada circunferencia --- es la que fuerza la forma
		gaussiana.}
	\label{fig:maxwell}
\end{figure}

\begin{observacion}
	Maxwell no buscaba justificar la distribución Normal ---
	buscaba entender el movimiento molecular. Pero las dos
	condiciones que postuló (isotropía e independencia
	direccional) son tan naturales que la gaussiana emerge
	como la \textbf{única solución posible}. La campana no se
	elige: se impone.
\end{observacion}

\subsubsection{La integral gaussiana y el origen de $\sqrt{2\pi}$}

El argumento de Maxwell justifica que la fdp tiene la forma
$ce^{-\alpha x^2}$. Para que sea una fdp válida debe integrar
a $1$, lo que requiere calcular:
\[
I = \int_{-\infty}^{\infty} e^{-x^2/2}\,dx.
\]
Esta integral no tiene antiderivada elemental --- no se puede
resolver con las técnicas estándar de cálculo. El truco
clásico es elevar $I$ al cuadrado y convertir la integral
doble resultante a coordenadas polares.

\begin{proposicion}{Integral gaussiana}{integral_gaussiana}
	\[
	\int_{-\infty}^{\infty} e^{-x^2/2}\,dx = \sqrt{2\pi}.
	\]
\end{proposicion}

\begin{proof}
	Llamemos $I = \int_{-\infty}^{\infty} e^{-x^2/2}\,dx$.
	Entonces:
	\[
	I^2 = \left(\int_{-\infty}^{\infty} e^{-x^2/2}\,dx\right)
	\left(\int_{-\infty}^{\infty} e^{-y^2/2}\,dy\right)
	= \int_{-\infty}^{\infty}\int_{-\infty}^{\infty}
	e^{-(x^2+y^2)/2}\,dx\,dy.
	\]
	Pasamos a coordenadas polares: $x = r\cos\theta$,
	$y = r\sin\theta$, con $r \geq 0$ y $\theta \in [0, 2\pi)$.
	El jacobiano de la transformación es $r$, y
	$x^2 + y^2 = r^2$:
	\[
	I^2 = \int_0^{2\pi}\int_0^{\infty}
	e^{-r^2/2}\,r\,dr\,d\theta
	= 2\pi \int_0^{\infty} r\,e^{-r^2/2}\,dr.
	\]
	La integral restante se resuelve con la sustitución
	$u = r^2/2$, $du = r\,dr$:
	\[
	\int_0^{\infty} r\,e^{-r^2/2}\,dr
	= \int_0^{\infty} e^{-u}\,du
	= \left[-e^{-u}\right]_0^{\infty} = 1.
	\]
	Por lo tanto $I^2 = 2\pi$, y como $I > 0$:
	\[
	I = \int_{-\infty}^{\infty} e^{-x^2/2}\,dx
	= \sqrt{2\pi}. \qedhere
	\]
\end{proof}

La constante $1/\sqrt{2\pi}$ en la fdp Normal es exactamente
la constante de normalización que garantiza que $f$ integre a
$1$ --- ni más ni menos. El $\sqrt{2\pi}$ no es un número
misterioso: es el área bajo la curva $e^{-x^2/2}$, calculada
por el argumento de coordenadas polares.

\begin{proposicion}{Momentos de la distribuci\'{o}n Normal}
	{momentos_normal}
	Si $X \sim N(\mu, \sigma^2)$, entonces
	\[
	E(X) = \mu \qquad \text{y} \qquad \mathrm{Var}(X) = \sigma^2.
	\]
\end{proposicion}

\begin{proof}
	Hacemos el cambio de variable $z = (x-\mu)/\sigma$, de modo
	que $x = \mu + \sigma z$ y $dx = \sigma\,dz$.
	
	\textbf{Valor esperado.}
	\[
	E(X) = \int_{-\infty}^{\infty} x\,\frac{1}{\sqrt{2\pi}\,\sigma}
	e^{-\frac{(x-\mu)^2}{2\sigma^2}}\,dx
	= \int_{-\infty}^{\infty} (\mu + \sigma z)\,
	\frac{1}{\sqrt{2\pi}}\,e^{-z^2/2}\,dz.
	\]
	Separamos en dos integrales:
	\[
	= \mu \int_{-\infty}^{\infty} \frac{1}{\sqrt{2\pi}}\,
	e^{-z^2/2}\,dz
	+ \sigma \int_{-\infty}^{\infty} z\,\frac{1}{\sqrt{2\pi}}\,
	e^{-z^2/2}\,dz.
	\]
	La primera integral es $1$ por la
	Proposición~\ref{prop:integral_gaussiana}. La segunda es $0$
	porque el integrando $z\,e^{-z^2/2}$ es una función
	\textbf{impar} --- y la integral de una función impar sobre
	$(-\infty, \infty)$ es cero. Por lo tanto:
	\[
	E(X) = \mu \cdot 1 + \sigma \cdot 0 = \mu.
	\]
	
	\textbf{Varianza.} Usamos la fórmula alternativa
	$\mathrm{Var}(X) = E\!\left[(X-\mu)^2\right]$. Con el mismo
	cambio de variable:
	\[
	\mathrm{Var}(X)
	= \int_{-\infty}^{\infty} (x-\mu)^2\,
	\frac{1}{\sqrt{2\pi}\,\sigma}\,
	e^{-\frac{(x-\mu)^2}{2\sigma^2}}\,dx
	= \sigma^2 \int_{-\infty}^{\infty} z^2\,
	\frac{1}{\sqrt{2\pi}}\,e^{-z^2/2}\,dz.
	\]
	Integramos por partes con $u = z$ y
	$dv = z\,e^{-z^2/2}\,dz$, luego $du = dz$ y
	$v = -e^{-z^2/2}$:
	\[
	\int_{-\infty}^{\infty} z^2\,e^{-z^2/2}\,dz
	= \left[-z\,e^{-z^2/2}\right]_{-\infty}^{\infty}
	+ \int_{-\infty}^{\infty} e^{-z^2/2}\,dz.
	\]
	El término evaluado se anula en ambos extremos porque
	$z\,e^{-z^2/2} \to 0$ cuando $z \to \pm\infty$. La integral
	restante es $\sqrt{2\pi}$ por la
	Proposición~\ref{prop:integral_gaussiana}. Por lo tanto:
	\[
	\mathrm{Var}(X) = \frac{\sigma^2}{\sqrt{2\pi}}
	\cdot \sqrt{2\pi} = \sigma^2. \qedhere
	\]
\end{proof}

\begin{observacion}
	El argumento de Maxwell y la integral gaussiana explican la
	\textbf{forma} de la fdp Normal --- por qué el exponente es
	cuadrático y de dónde viene $\sqrt{2\pi}$. Pero no responden
	la pregunta más importante: ¿por qué la Normal aparece en
	fenómenos tan distintos como alturas, errores de medición y
	precios financieros? La respuesta no está en la física de los
	gases sino en un resultado matemático profundo que estudiaremos
	a continuación: el \textbf{Teorema del Límite Central}, que
	establece que la suma de un gran número de variables aleatorias
	independientes converge a una distribución Normal,
	independientemente de la distribución original de cada una.
	Esa es la verdadera razón de la ubicuidad de la Normal.
\end{observacion}

\subsection{La distribuci\'{o}n Normal est\'{a}ndar}

La fdp Normal depende de dos parámetros: $\mu$ y $\sigma^2$.
Para calcular probabilidades habría que integrar numéricamente
para cada par $(\mu, \sigma^2)$ distinto. La solución es
trabajar con un caso de referencia y mostrar que toda Normal
general puede reducirse a él. Definimos ese caso de referencia
y demostramos la reducción inmediatamente después.

La \textbf{distribuci\'{o}n Normal est\'{a}ndar} es el caso
particular $\mu = 0$ y $\sigma^2 = 1$ de la
Definici\'{o}n~\ref{def:dist_normal}. Sustituyendo directamente:
\[
f(x)\big|_{\mu=0,\,\sigma^2=1}
= \frac{1}{\sqrt{2\pi}}\,e^{-z^2/2}
\eqqcolon \phi(z).
\]

\begin{definicion}{Distribuci\'{o}n Normal est\'{a}ndar}
	{normal_estandar}
	Una variable aleatoria $Z \sim N(0,1)$ tiene
	\textbf{distribuci\'{o}n Normal est\'{a}ndar}, con fdp
	$\phi(z) = \frac{1}{\sqrt{2\pi}}\,e^{-z^2/2}$ y FDA
	\[
	\Phi(z) = P(Z \leq z) =
	\int_{-\infty}^{z} \phi(t)\,dt,
	\qquad z \in \mathbb{R}.
	\]
	La función $\Phi$ no admite forma cerrada y debe tabularse
	o calcularse numéricamente.
\end{definicion}

\begin{proposicion}{Simetr\'{\i}a de la Normal est\'{a}ndar}
	{simetria_normal}
	Para todo $z \in \mathbb{R}$:
	\[
	\Phi(-z) = 1 - \Phi(z).
	\]
\end{proposicion}

\begin{proof}
	Como $\phi$ es simétrica respecto a cero --- $\phi(-z) =
	\phi(z)$ para todo $z$ --- se tiene:
	\[
	\Phi(-z) = \int_{-\infty}^{-z} \phi(t)\,dt
	= \int_{z}^{\infty} \phi(t)\,dt
	= 1 - \int_{-\infty}^{z} \phi(t)\,dt
	= 1 - \Phi(z). \qedhere
	\]
\end{proof}

\begin{observacion}
	Esta proposición es la herramienta práctica más usada al
	trabajar con la tabla de la Normal estándar: permite calcular
	probabilidades en la cola izquierda a partir de la cola
	derecha, y viceversa. Por ejemplo,
	$P(Z \leq -1.96) = 1 - P(Z \leq 1.96) = 1 - \Phi(1.96)$.
\end{observacion}

\begin{proposicion}{Estandarizaci\'{o}n}{estandarizacion}
	Si $X \sim N(\mu, \sigma^2)$, entonces
	\[
	Z = \frac{X - \mu}{\sigma} \sim N(0,1).
	\]
\end{proposicion}

\begin{proof}
	Calculamos la FDA de $Z$. Para $z \in \mathbb{R}$:
	\[
	P(Z \leq z) = P\!\left(\frac{X-\mu}{\sigma} \leq z\right)
	= P(X \leq \mu + \sigma z)
	= \int_{-\infty}^{\mu+\sigma z}
	\frac{1}{\sqrt{2\pi}\,\sigma}
	e^{-\frac{(x-\mu)^2}{2\sigma^2}}\,dx.
	\]
	Sustituimos $t = (x-\mu)/\sigma$, $dt = dx/\sigma$:
	\[
	= \int_{-\infty}^{z} \frac{1}{\sqrt{2\pi}}
	e^{-t^2/2}\,dt = \Phi(z).
	\]
	Como la FDA de $Z$ es $\Phi$, se concluye que
	$Z \sim N(0,1)$. \qedhere
\end{proof}

\begin{observacion}
	La estandarización tiene una interpretación directa: $Z$
	mide cuántas desviaciones estándar se aleja $X$ de su
	media. Es la transformación que hace que todas las Normales
	sean comparables entre sí y que una sola tabla --- la de
	$\Phi$ --- sea suficiente para calcular cualquier
	probabilidad Normal. Para calcular $P(X \leq x)$
	con $X \sim N(\mu,\sigma^2)$:
	\[
	P(X \leq x) = P\!\left(Z \leq \frac{x-\mu}{\sigma}\right)
	= \Phi\!\left(\frac{x-\mu}{\sigma}\right).
	\]
\end{observacion}

\begin{figure}[H]
	\centering
	\begin{tikzpicture}[scale=0.95]
		
		% ── Panel izquierdo: distintos mu, sigma=1 ─────────────
		\begin{scope}[shift={(0,0)}]
			\node[font=\small\bfseries, azul] at (3.0, 3.5)
			{Distinto $\mu$, mismo $\sigma = 1$};
			
			\draw[thick,->] (-0.5,0) -- (6.5,0) node[right] {$x$};
			\draw[thick,->] (0,-0.2) -- (0,3.3) node[above] {};
			
			% mu=1, sigma=1
			\draw[azul!50, thick, domain=0:6, samples=100]
			plot (\x, {3.5*(1/(sqrt(2*pi)))*exp(-(\x-1)^2/2)});
			\node[azul!50, font=\scriptsize] at (1.0, 2.0)
			{$\mu=1$};
			
			% mu=3, sigma=1
			\draw[azul, thick, domain=0:6, samples=100]
			plot (\x, {3.5*(1/(sqrt(2*pi)))*exp(-(\x-3)^2/2)});
			\node[azul, font=\scriptsize] at (3.0, 2.0)
			{$\mu=3$};
			
			% mu=5, sigma=1
			\draw[azul!80!black, thick, domain=0:6, samples=100]
			plot (\x, {3.5*(1/(sqrt(2*pi)))*exp(-(\x-5)^2/2)});
			\node[azul!80!black, font=\scriptsize] at (5.0, 2.0)
			{$\mu=5$};
		\end{scope}
		
		% ── Panel derecho: mismo mu=3, distintos sigma ──────────
		\begin{scope}[shift={(8.0,0)}]
			\node[font=\small\bfseries, azul] at (3.0, 3.5)
			{Mismo $\mu=3$, distinto $\sigma$};
			
			\draw[thick,->] (-0.5,0) -- (6.5,0) node[right] {$x$};
			\draw[thick,->] (0,-0.2) -- (0,3.3) node[above] {};
			
			% mu=3, sigma=0.5
			\draw[azul!80!black, very thick, domain=0:6, samples=100]
			plot (\x, {3.5*(1/(0.5*sqrt(2*pi)))*
				exp(-(\x-3)^2/(2*0.5^2))});
			\node[azul!80!black, font=\scriptsize] at (4.2, 2.8)
			{$\sigma=0.5$};
			
			% mu=3, sigma=1
			\draw[azul, thick, domain=0:6, samples=100]
			plot (\x, {3.5*(1/(sqrt(2*pi)))*exp(-(\x-3)^2/2)});
			\node[azul, font=\scriptsize] at (4.5, 1.6)
			{$\sigma=1$};
			
			% mu=3, sigma=1.5
			\draw[azul!50, thick, domain=0:6, samples=100]
			plot (\x, {3.5*(1/(1.5*sqrt(2*pi)))*
				exp(-(\x-3)^2/(2*1.5^2))});
			\node[azul!50, font=\scriptsize] at (5.0, 0.8)
			{$\sigma=1.5$};
		\end{scope}
		
	\end{tikzpicture}
	\caption{Izquierda: distintas medias $\mu$ con $\sigma=1$ ---
		$\mu$ desplaza la curva sin cambiar su forma. Derecha: misma
		media $\mu=3$ con distintas $\sigma$ --- $\sigma$ controla el
		ancho y la altura de la campana.}
	\label{fig:normal}
\end{figure}

En la práctica, el objetivo es calcular probabilidades del tipo
$P(a \leq X \leq b)$ para $X \sim N(\mu, \sigma^2)$. La
Proposici\'{o}n~\ref{prop:estandarizacion} reduce este problema
a uno sobre la Normal est\'{a}ndar:
\[
P(a \leq X \leq b)
= \Phi\!\left(\frac{b-\mu}{\sigma}\right)
- \Phi\!\left(\frac{a-\mu}{\sigma}\right).
\]
El problema entonces es evaluar $\Phi(z)$ --- y aquí aparece
la dificultad: la integral $\int_{-\infty}^z \phi(t)\,dt$ no
tiene forma cerrada, es decir, no puede expresarse en términos
de funciones elementales. Por eso se recurre a
\textbf{tablas numéricas} que precalculan $\Phi(z)$ para una
grilla de valores de $z$, o a software estadístico. El
procedimiento concreto para calcular cualquier probabilidad
Normal es:
\begin{enumerate}[label=\textbf{\arabic*.}]
	\item \textbf{Estandarizar} los límites: $z_a = (a-\mu)/\sigma$
	y $z_b = (b-\mu)/\sigma$.
	\item \textbf{Buscar} $\Phi(z_a)$ y $\Phi(z_b)$ en la tabla
	o software.
	\item \textbf{Restar:} $P(a \leq X \leq b) =
	\Phi(z_b) - \Phi(z_a)$.
\end{enumerate}
La propiedad de simetría $\Phi(-z) = 1 - \Phi(z)$
(Proposici\'{o}n~\ref{prop:simetria_normal}) es útil cuando
la tabla solo cubre valores positivos de $z$: basta con cambiar
el signo y complementar a $1$.

% ==========================================================
\chapter{Muestreo y distribuci\'{o}n de muestreo}

\section{Poblaci\'{o}n, muestra y par\'{a}metro}

En los capítulos anteriores trabajamos con variables aleatorias
cuya distribución era completamente conocida: sabíamos que
$X \sim N(\mu, \sigma^2)$ con valores específicos de $\mu$ y
$\sigma^2$, y calculábamos probabilidades a partir de eso. En
la práctica, sin embargo, ocurre algo distinto: a menudo
tenemos razones para creer que una variable sigue cierto
\emph{tipo} de distribución --- una Normal, una Poisson, una
Exponencial --- pero \textbf{no conocemos los parámetros
	específicos} que la determinan. Sabemos que las alturas de una
población se distribuyen normalmente, pero no sabemos cuáles son
$\mu$ y $\sigma^2$. Sabemos que las llegadas a un servidor
siguen un proceso de Poisson, pero no conocemos $\lambda$.

El problema central de la inferencia estadística es precisamente
este: dado que conocemos la \emph{familia} de distribuciones a
la que pertenece la variable de interés, ¿cómo estimamos el
parámetro desconocido que la identifica dentro de esa familia?
La herramienta es la \textbf{muestra} --- un conjunto de
observaciones generadas por la distribución real --- y el
objetivo es usar esas observaciones para aprender algo sobre el
parámetro desconocido. Este capítulo formaliza ese proceso.

\begin{definicion}{Familia indexada}{familia_indexada}
	Sea $\Theta$ un conjunto no vacío. Una \textbf{familia
		indexada} por $\Theta$ es una función
	\[
	\theta \mapsto P_\theta
	\]
	que asigna a cada elemento $\theta \in \Theta$ un objeto
	$P_\theta$ de alguna colección. El conjunto $\Theta$ se
	denomina \textbf{conjunto de índices} y cada $\theta$ es
	un \textbf{índice}.
\end{definicion}

\begin{observacion}
	La indexación no es más que una función con notación
	distinta. En lugar de escribir $f(\theta)$ para el
	objeto asociado a $\theta$, escribimos $P_\theta$ ---
	el índice aparece como subíndice en lugar de argumento.
	La razón es puramente notacional: cuando los objetos
	son distribuciones de probabilidad, $P_\theta$ es más
	natural que $P(\theta)$, que podría confundirse con
	la probabilidad de un evento $\theta$.
	
	En el caso de una familia paramétrica, los objetos
	indexados son distribuciones de probabilidad y el
	conjunto de índices $\Theta$ es el espacio paramétrico.
	Decir que $\mathcal{F} = \{P_\theta : \theta \in
	\Theta\}$ está indexada por $\Theta$ significa
	simplemente que existe una función biunívoca entre
	$\Theta$ y las distribuciones de $\mathcal{F}$ ---
	cada $\theta$ determina unívocamente una distribución
	y cada distribución corresponde a exactamente un
	$\theta$.
\end{observacion}

\begin{definicion}{Familia param\'{e}trica}{familia_parametrica}
	Sea $\Theta \subseteq \mathbb{R}^k$ un conjunto llamado
	\textbf{espacio param\'{e}trico}. Una \textbf{familia
		param\'{e}trica} es una colección de distribuciones de
	probabilidad indexada por $\Theta$:
	\[
	\mathcal{F} = \{P_\theta : \theta \in \Theta\},
	\]
	donde cada $\theta \in \Theta$ determina unívocamente una
	distribución $P_\theta$.
\end{definicion}

\begin{definicion}{Poblaci\'{o}n}{poblacion}
	Dada una familia paramétrica $\mathcal{F} = \{P_\theta :
	\theta \in \Theta\}$, la \textbf{población} es la
	distribución $P_{\theta_0} \in \mathcal{F}$ que gobierna
	la variable de interés, donde $\theta_0 \in \Theta$ es el
	\textbf{verdadero valor del parámetro} --- fijo pero
	desconocido.
\end{definicion}

\begin{definicion}{Par\'{a}metro}{parametro}
	El \textbf{parámetro} es el valor $\theta_0 \in \Theta$
	que indexa la población real $P_{\theta_0}$. Los parámetros
	más comunes son escalares como la media $\mu$ o la tasa
	$\lambda$, o vectores como $(\mu, \sigma^2)$ en el caso
	Normal.
\end{definicion}

\begin{observacion}
	La distinción entre familia y población es fundamental.
	La \emph{familia} $\mathcal{F}$ es el conocimiento previo
	que traemos al problema --- el tipo de distribución que
	creemos describe el fenómeno. La \emph{población}
	$P_{\theta_0}$ es la distribución específica dentro de
	esa familia que realmente opera en el mundo, pero cuyo
	parámetro $\theta_0$ desconocemos. El objetivo de la
	inferencia estadística es estimar $\theta_0$ a partir
	de datos observados.
\end{observacion}

\begin{definicion}{Muestra}{muestra}
	Una \textbf{muestra} de tamaño $n$ es un conjunto de
	observaciones $\{x_1, x_2, \ldots, x_n\}$ generadas por
	la población $P_{\theta_0}$.
\end{definicion}

En términos simples:

\begin{itemize}[leftmargin=2em]
	\item La \textbf{población} es el fenómeno completo que
	queremos estudiar --- la distribución real que gobierna
	la variable de interés. Existe, es fija, pero no la
	conocemos completamente.
	\item El \textbf{parámetro} es el número (o números) que
	describen esa población. Conocer el parámetro equivale
	a conocer la población: si supiéramos que las alturas
	siguen $N(170, 49)$, sabríamos todo lo que hay que saber
	sobre su distribución. El problema es que $\theta_0$
	es desconocido.
	\item La \textbf{muestra} es el subconjunto de
	observaciones que podemos obtener. Es nuestra única
	ventana hacia la población --- a partir de ella
	intentamos inferir $\theta_0$.
\end{itemize}

El objetivo de la inferencia estadística es usar la muestra
para aprender sobre el parámetro, sabiendo que nunca
tendremos acceso directo a la población completa.

\section{Muestra aleatoria}

Para que las conclusiones extraídas de una muestra sean
válidas, la selección de sus elementos no puede ser
arbitraria. La forma más rigurosa de garantizar esto es
modelar cada observación como una variable aleatoria
independiente con la misma distribución poblacional.

\begin{definicion}{Muestra aleatoria}{muestra_aleatoria}
	Las variables aleatorias $X_1, X_2, \ldots, X_n$
	constituyen una \textbf{muestra aleatoria} de tamaño $n$
	de la población $P_{\theta_0}$ si son independientes e
	idénticamente distribuidas (i.i.d.) con distribución
	$P_{\theta_0}$.
\end{definicion}

\begin{observacion}
	Esta definición es exactamente la misma que la de
	variables i.i.d.\ introducida en la
	Definición~\ref{def:iid}. La muestra aleatoria no es un
	concepto nuevo --- es el mismo objeto matemático aplicado
	ahora a un contexto de inferencia. Cada $X_i$ representa
	el resultado de observar un elemento de la población
	seleccionado al azar; la independencia garantiza que
	ninguna observación influye sobre las demás, y la identidad
	de distribución garantiza que todas provienen de la misma
	población $P_{\theta_0}$.
	
	Esta definición supone implícitamente que la población es
	suficientemente grande para que la extracción de cada
	elemento no altere la distribución de los restantes ---
	es decir, se asume muestreo con reposición o una población
	infinita.
\end{observacion}

\section{Estad\'{\i}stico y estimador}

A partir de una muestra aleatoria $X_1, \ldots, X_n$
podemos calcular distintas funciones de los datos. Cada
una de esas funciones es un \textbf{estadístico}.

\begin{definicion}{Estad\'{\i}stico}{estadistico}
	Un \textbf{estadístico} es cualquier función
	\[
	T = g(X_1, X_2, \ldots, X_n)
	\]
	de las observaciones de una muestra aleatoria que no
	depende de parámetros desconocidos. Un estadístico es
	en sí mismo una variable aleatoria --- su valor varía
	de muestra en muestra.
\end{definicion}

\begin{observacion}
	La condición de que el estadístico no dependa de
	parámetros desconocidos es esencial: si $T$ dependiera
	de $\theta_0$, no podríamos calcularlo sin conocer
	precisamente lo que queremos estimar. Un estadístico
	debe ser computable directamente a partir de los datos
	observados.
	
	Que un estadístico sea una variable aleatoria es
	fundamental: no tiene un valor fijo, sino una
	distribución propia --- la \textbf{distribución
		muestral} --- que depende de la distribución de la
	población $P_{\theta_0}$. Estudiar esa distribución
	es lo que permite hacer inferencia sobre $\theta_0$.
\end{observacion}

Cuando un estadístico se usa específicamente para aproximar
un parámetro desconocido, se denomina \textbf{estimador}.

\begin{definicion}{Estimador y estimaci\'{o}n puntual}
	{estimador}
	Un \textbf{estimador puntual} de $\theta_0$ es un
	estadístico $\hat{\theta} = g(X_1, \ldots, X_n)$ usado
	para aproximar $\theta_0$. Una vez observada la muestra
	$(x_1, \ldots, x_n)$, el valor numérico
	$\hat{\theta}(x_1, \ldots, x_n)$ se denomina
	\textbf{estimación puntual} de $\theta_0$.
\end{definicion}

\begin{observacion}
	La distinción entre estimador y estimación es análoga
	a la distinción entre variable aleatoria y realización:
	el estimador $\hat{\theta}$ es la \emph{regla},
	expresada antes de observar los datos; la estimación
	es el \emph{número} que resulta de aplicar esa regla
	a los datos observados. Por ejemplo, la media muestral
	\[
	\bar{X} = \frac{1}{n}\sum_{i=1}^n X_i
	\]
	es un estimador de $\mu$; el valor $\bar{x} = 28.75$
	calculado de una muestra particular es una estimación.
\end{observacion}

\section{Propiedades de los estimadores puntuales}
\label{sec:propiedades_estimadores}

Un parámetro $\theta_0$ puede tener muchos estimadores
distintos. ¿Cómo elegir entre ellos? La respuesta no es
arbitraria --- existen criterios matemáticos precisos que
permiten evaluar la calidad de un estimador. Estudiaremos
tres: \textbf{insesgadez}, \textbf{eficiencia} y
\textbf{consistencia}. Cada uno responde una pregunta
distinta sobre el comportamiento del estimador.

\subsection{Insesgadez}

La primera pregunta que podemos hacerle a un estimador es
la más básica: ¿apunta en la dirección correcta? Un
estimador que sistemáticamente sobreestima o subestima
$\theta_0$ tiene un error estructural que no desaparece
sin importar cuántos datos tengamos. La insesgadez
formaliza la ausencia de ese error sistemático.

\begin{definicion}{Estimador insesgado}{insesgadez}
	Un estimador $\hat{\theta}$ de $\theta_0$ es
	\textbf{insesgado} si
	\[
	E_{\theta_0}(\hat{\theta}) = \theta_0
	\qquad \text{para todo } \theta_0 \in \Theta.
	\]
	La cantidad
	\[
	B(\hat{\theta}) = E_{\theta_0}(\hat{\theta}) - \theta_0
	\]
	se denomina \textbf{sesgo} del estimador. Si
	$B(\hat{\theta}) = 0$, el estimador es insesgado.
\end{definicion}

\begin{observacion}
	La notación $E_{\theta_0}$ enfatiza que el valor esperado
	se calcula bajo la distribución $P_{\theta_0}$ ---
	asumiendo que el verdadero parámetro es $\theta_0$. La
	condición debe valer para \emph{todo} $\theta_0 \in
	\Theta$: un estimador insesgado apunta al parámetro
	correcto sin importar cuál sea ese parámetro.
	
	La insesgadez no garantiza que una estimación particular
	sea exacta --- solo garantiza que el estimador no tiene
	un error sistemático. En cualquier muestra concreta,
	$\hat{\theta}$ puede diferir de $\theta_0$; lo que la
	insesgadez asegura es que esas diferencias se cancelan
	en promedio.
\end{observacion}

\begin{observacion}
	El \textbf{principio de estimación insesgada} establece
	que, ante varios estimadores de $\theta_0$, se prefiere
	uno insesgado. Si todos son insesgados, se aplica el
	criterio de eficiencia que se describe a continuación.
\end{observacion}

\subsection{Eficiencia}

Dos estimadores pueden ser ambos insesgados --- ambos
apuntan en la dirección correcta en promedio --- y aun así
diferir en cuán dispersas son sus estimaciones alrededor
de $\theta_0$. El que produce estimaciones más concentradas
cerca del parámetro verdadero es preferible. La eficiencia
formaliza esta comparación.

\begin{definicion}{Eficiencia relativa}{eficiencia}
	Sean $\hat{\theta}_1$ y $\hat{\theta}_2$ dos estimadores
	insesgados de $\theta_0$. Se dice que $\hat{\theta}_1$
	es \textbf{más eficiente} que $\hat{\theta}_2$ si
	\[
	\mathrm{Var}_{\theta_0}(\hat{\theta}_1) 
	\mathrm{Var}_{\theta_0}(\hat{\theta}_2)
	\qquad \text{para todo } \theta_0 \in \Theta.
	\]
\end{definicion}

\begin{observacion}
	La raíz cuadrada de la varianza del estimador,
	\[
	\mathrm{SE}(\hat{\theta}) =
	\sqrt{\mathrm{Var}_{\theta_0}(\hat{\theta})},
	\]
	se denomina \textbf{error estándar} del estimador y es
	la medida práctica más usada de su precisión. Un
	estimador con menor error estándar produce estimaciones
	más estables y más cercanas a $\theta_0$ de muestra
	en muestra.
	
	Por ahora comparamos estimadores de a pares. En la
	Sección~\ref{sec:cramer_rao} veremos que el Teorema de
	Cramér-Rao establece una cota absoluta --- la menor
	varianza que puede tener cualquier estimador insesgado
	--- lo que permitirá hablar de eficiencia en un sentido
	más fuerte que la mera comparación entre dos
	estimadores.
\end{observacion}

La figura de los tiradores ilustra la diferencia entre
insesgadez y eficiencia. El centro de la diana representa
$\theta_0$; cada disparo representa una estimación obtenida
de una muestra distinta.

\begin{figure}[H]
	\centering
	\begin{tikzpicture}[scale=0.85]
		
		% ── Tirador A: insesgado, no eficiente ──────────────
		\begin{scope}[shift={(0,0)}]
			\foreach \r in {1.8, 1.2, 0.6}{
				\draw[gray, thin] (0,0) circle (\r);
			}
			\draw[gray, very thin, dashed] (-2,0)--(2,0);
			\draw[gray, very thin, dashed] (0,-2)--(0,2);
			\fill[rojo] (-1.2, 0.8) circle (2.5pt);
			\fill[rojo] ( 0.9,-1.3) circle (2.5pt);
			\fill[rojo] (-0.7,-1.0) circle (2.5pt);
			\fill[rojo] ( 1.4, 0.3) circle (2.5pt);
			\fill[rojo] (-0.3, 1.5) circle (2.5pt);
			\fill[rojo] ( 0.5,-0.9) circle (2.5pt);
			\node[font=\small\bfseries] at (0,-2.5) {A};
			\node[font=\scriptsize, gray, align=center] at (0,-3.1)
			{insesgado\\no eficiente};
		\end{scope}
		
		% ── Tirador B: sesgado, eficiente ───────────────────
		\begin{scope}[shift={(5.5,0)}]
			\foreach \r in {1.8, 1.2, 0.6}{
				\draw[gray, thin] (0,0) circle (\r);
			}
			\draw[gray, very thin, dashed] (-2,0)--(2,0);
			\draw[gray, very thin, dashed] (0,-2)--(0,2);
			\fill[azul] (-0.8, 1.2) circle (2.5pt);
			\fill[azul] (-1.0, 1.4) circle (2.5pt);
			\fill[azul] (-0.9, 1.1) circle (2.5pt);
			\fill[azul] (-0.7, 1.3) circle (2.5pt);
			\fill[azul] (-1.1, 1.0) circle (2.5pt);
			\fill[azul] (-0.8, 1.5) circle (2.5pt);
			\node[font=\small\bfseries] at (0,-2.5) {B};
			\node[font=\scriptsize, gray, align=center] at (0,-3.1)
			{sesgado\\eficiente};
		\end{scope}
		
		% ── Tirador C: insesgado, eficiente ─────────────────
		\begin{scope}[shift={(11,0)}]
			\foreach \r in {1.8, 1.2, 0.6}{
				\draw[gray, thin] (0,0) circle (\r);
			}
			\draw[gray, very thin, dashed] (-2,0)--(2,0);
			\draw[gray, very thin, dashed] (0,-2)--(0,2);
			\fill[verde!70!black] (-0.2, 0.3) circle (2.5pt);
			\fill[verde!70!black] ( 0.3,-0.2) circle (2.5pt);
			\fill[verde!70!black] (-0.3,-0.3) circle (2.5pt);
			\fill[verde!70!black] ( 0.2, 0.4) circle (2.5pt);
			\fill[verde!70!black] (-0.1,-0.3) circle (2.5pt);
			\fill[verde!70!black] ( 0.3, 0.1) circle (2.5pt);
			\node[font=\small\bfseries] at (0,-2.5) {C};
			\node[font=\scriptsize, gray, align=center] at (0,-3.1)
			{insesgado\\eficiente};
		\end{scope}
		
	\end{tikzpicture}
	\caption{Analogía entre estimadores y tiradores. El centro
		de la diana representa $\theta_0$. Cada punto es una
		estimación obtenida de una muestra distinta. El estimador~A
		es insesgado pero no eficiente: las estimaciones se centran
		en $\theta_0$ pero son muy dispersas. El estimador~B es
		eficiente pero sesgado: las estimaciones son muy
		concentradas pero apuntan sistemáticamente al lugar
		equivocado. El estimador~C es insesgado \textbf{y}
		eficiente: el ideal.}
	\label{fig:tiradores}
\end{figure}

\subsection{Consistencia}

Los criterios de insesgadez y eficiencia evalúan el
estimador para un tamaño de muestra $n$ fijo. La
consistencia evalúa su comportamiento cuando $n$ crece:
¿mejora el estimador a medida que tenemos más datos?

\begin{definicion}{Estimador consistente}{consistencia}
	Un estimador $\hat{\theta}_n$ de $\theta_0$ es
	\textbf{consistente} si converge en probabilidad a
	$\theta_0$ cuando $n \to \infty$: para todo
	$\varepsilon > 0$ y todo $\theta_0 \in \Theta$,
	\[
	\lim_{n \to \infty} P_{\theta_0}\!\left(
	|\hat{\theta}_n - \theta_0| > \varepsilon
	\right) = 0.
	\]
\end{definicion}

\begin{observacion}
	La consistencia es una propiedad asintótica: garantiza
	que con suficientes datos el estimador puede aproximar
	$\theta_0$ con tanta precisión como se desee. Es una
	condición mínima deseable --- un estimador que no
	mejora con más datos es de poco valor práctico.
	
	Un criterio suficiente para la consistencia, y el más
	utilizado en la práctica, es que el sesgo y la varianza
	del estimador converjan ambos a cero cuando
	$n \to \infty$:
	\[
	B(\hat{\theta}_n) \to 0
	\qquad \text{y} \qquad
	\mathrm{Var}_{\theta_0}(\hat{\theta}_n) \to 0
	\qquad \text{cuando } n \to \infty.
	\]
	Este criterio suficiente se justifica formalmente
	mediante la \textbf{desigualdad de Chebyshev}:
	\[
	P_{\theta_0}(|\hat{\theta}_n - \theta_0|
	> \varepsilon)
	\leq \frac{E_{\theta_0}[(\hat{\theta}_n -
		\theta_0)^2]}{\varepsilon^2}
	= \frac{\mathrm{Var}_{\theta_0}(\hat{\theta}_n)
		+ B(\hat{\theta}_n)^2}{\varepsilon^2}.
	\]
	Si sesgo y varianza convergen a cero, el lado derecho
	converge a cero para todo $\varepsilon > 0$, lo que
	establece la consistencia.
\end{observacion}

\section{Funci\'{o}n de verosimilitud e informaci\'{o}n de
	Fisher}

En la secci\'{o}n anterior definimos la eficiencia como la
comparaci\'{o}n de varianzas entre estimadores insesgados.
Pero esa definici\'{o}n es relativa --- solo dice que
$\hat{\theta}_1$ es mejor que $\hat{\theta}_2$, sin
responder si existe un l\'{i}mite absoluto a cu\'{a}n
preciso puede ser cualquier estimador insesgado. Las
preguntas que quedan abiertas son tres:

\begin{itemize}[leftmargin=2em]
	\item ¿Existe una varianza m\'{i}nima que ning\'{u}n
	estimador insesgado puede superar?
	\item Si existe, ¿de qu\'{e} depende ese l\'{i}mite?
	\item ¿Hay alg\'{u}n estimador que lo alcance?
\end{itemize}

La respuesta a las tres viene del Teorema de Cram\'{e}r-Rao,
que estudiaremos en la secci\'{o}n siguiente. Pero para
llegar ah\'{i} necesitamos primero un concepto nuevo:
cu\'{a}nta \textbf{informaci\'{o}n} contienen los datos
sobre el par\'{a}metro $\theta_0$. La intuici\'{o}n es
directa: si los datos son muy informativos sobre $\theta_0$,
es posible construir estimadores muy precisos. Si los datos
contienen poca informaci\'{o}n, ning\'{u}n estimador
podr\'{a} ser preciso --- hay un l\'{i}mite impuesto por
los propios datos. Formalizar esa intuici\'{o}n requiere
primero introducir la \textbf{funci\'{o}n de verosimilitud}.

\subsection{Funci\'{o}n de verosimilitud}

Dada una muestra observada $x_1, x_2, \ldots, x_n$,
distintos valores de $\theta$ hacen que esa muestra sea
m\'{a}s o menos probable. La pregunta natural es: ¿cu\'{a}n
probable era observar exactamente estos datos bajo cada
posible valor de $\theta$? La respuesta viene directamente
de la definici\'{o}n de muestra aleatoria
(Definici\'{o}n~\ref{def:muestra_aleatoria}): como las
observaciones son i.i.d., la probabilidad conjunta de
observar $x_1, x_2, \ldots, x_n$ bajo $P_\theta$ es
\[
P_\theta(X_1 = x_1, \ldots, X_n = x_n)
= \prod_{i=1}^n p_\theta(x_i).
\]
Esta cantidad, vista como funci\'{o}n de $\theta$ con
los datos fijos, es la \textbf{funci\'{o}n de
	verosimilitud} --- y es lo que define cu\'{a}n compatible
es cada valor de $\theta$ con los datos observados.

\begin{definicion}{Funci\'{o}n de verosimilitud}
	{verosimilitud}
	Sea $X_1, X_2, \ldots, X_n$ una muestra aleatoria de
	$P_\theta$ con funci\'{o}n de masa o densidad
	$p_\theta$. La \textbf{funci\'{o}n de verosimilitud}
	es la funci\'{o}n $L: \Theta \to \mathbb{R}$ definida
	por
	\[
	L(\theta) = L(\theta;\, x_1, \ldots, x_n)
	= \prod_{i=1}^{n} p_\theta(x_i),
	\qquad \theta \in \Theta.
	\]
\end{definicion}

\begin{observacion}
	Hay una distinci\'{o}n crucial entre $p_\theta(x)$ y
	$L(\theta)$:
	\begin{itemize}[leftmargin=2em]
		\item $p_\theta(x)$ es una funci\'{o}n de $x$ con
		$\theta$ fijo --- describe la distribuci\'{o}n de
		los datos dado el par\'{a}metro.
		\item $L(\theta)$ es una funci\'{o}n de $\theta$
		con los datos $x_1, \ldots, x_n$ fijos ---
		describe cu\'{a}n compatible es cada valor de
		$\theta$ con los datos observados.
	\end{itemize}
	La verosimilitud no es una funci\'{o}n de densidad
	sobre $\Theta$ --- no integra necesariamente a $1$.
	Es simplemente una medida relativa de compatibilidad
	entre par\'{a}metros y datos.
	
	El producto $\prod_{i=1}^n p_\theta(x_i)$ se justifica
	por la independencia de la muestra aleatoria: la
	probabilidad conjunta de observar $x_1, \ldots, x_n$
	es el producto de las probabilidades individuales.
\end{observacion}

\subsection{Log-verosimilitud y funci\'{o}n score}

Trabajar con el producto $\prod_{i=1}^n p_\theta(x_i)$
es algebraicamente inc\'{o}modo, especialmente al
diferenciar. La soluci\'{o}n est\'{a}ndar es tomar
logaritmo --- lo que convierte el producto en una suma
y preserva el orden de los m\'{a}ximos porque el logaritmo
es una funci\'{o}n estrictamente creciente.

\begin{definicion}{Log-verosimilitud y funci\'{o}n score}
	{score}
	La \textbf{log-verosimilitud} es
	\[
	\ell(\theta) = \ln L(\theta)
	= \sum_{i=1}^{n} \ln p_\theta(x_i).
	\]
	La \textbf{funci\'{o}n score} es la derivada de la
	log-verosimilitud respecto a $\theta$:
	\[
	s(\theta) = \frac{\partial}{\partial\theta}
	\ell(\theta)
	= \sum_{i=1}^{n} \frac{\partial}{\partial\theta}
	\ln p_\theta(x_i).
	\]
\end{definicion}

\begin{observacion}
	La funci\'{o}n score mide cu\'{a}n sensible es la
	log-verosimilitud al par\'{a}metro $\theta$. Si
	$s(\theta_0) = 0$, la log-verosimilitud alcanza un
	punto estacionario en $\theta_0$ --- intuitivamente,
	los datos son compatibles con $\theta_0$ y no
	``empujan'' hacia ning\'{u}n otro valor. Esta
	propiedad es la base del estimador de m\'{a}xima
	verosimilitud, que estudiaremos en cursos posteriores.
\end{observacion}

Antes de definir la informaci\'{o}n de Fisher necesitamos
un resultado fundamental sobre el score: que su valor
esperado bajo $P_{\theta_0}$ es cero.

\begin{proposicion}{Media cero del score}{score_media_cero}
	Bajo condiciones de regularidad sobre $p_\theta$,
	\[
	E_{\theta_0}[s(\theta_0)] = 0.
	\]
\end{proposicion}

\begin{proof}
	Como $p_\theta$ es una funci\'{o}n de densidad,
	integra a $1$ para todo $\theta$:
	\[
	\int_{-\infty}^{\infty} p_\theta(x)\,dx = 1.
	\]
	Diferenciando ambos lados respecto a $\theta$ e
	intercambiando derivada e integral (condici\'{o}n de
	regularidad):
	\[
	\int_{-\infty}^{\infty}
	\frac{\partial}{\partial\theta}
	p_\theta(x)\,dx = 0.
	\]
	Usando la identidad
	$\frac{\partial}{\partial\theta}\ln p_\theta(x)
	= \frac{1}{p_\theta(x)}
	\frac{\partial}{\partial\theta}p_\theta(x)$:
	\[
	\int_{-\infty}^{\infty}
	\frac{\partial}{\partial\theta}\ln p_\theta(x)
	\cdot p_\theta(x)\,dx = 0.
	\]
	El lado izquierdo es exactamente
	$E_{\theta_0}[s(\theta_0)]$ para una muestra de
	tama\~{n}o $n=1$. Para una muestra de tama\~{n}o
	$n$ general, por linealidad del valor esperado:
	\[
	E_{\theta_0}[s(\theta_0)]
	= \sum_{i=1}^n E_{\theta_0}\!\left[
	\frac{\partial}{\partial\theta}
	\ln p_\theta(X_i)\right] = 0. \qedhere
	\]
\end{proof}

\begin{observacion}
	Las \textbf{condiciones de regularidad} mencionadas en
	la proposici\'{o}n son condiciones t\'{e}cnicas sobre
	$p_\theta$ que permiten intercambiar el orden de
	derivaci\'{o}n e integraci\'{o}n. La m\'{a}s
	importante es que el soporte de $p_\theta$ --- el
	conjunto de $x$ donde $p_\theta(x) > 0$ --- no
	dependa de $\theta$. Todas las familias param\'{e}tricas
	cl\'{a}sicas (Normal, Poisson, Exponencial, Binomial)
	satisfacen estas condiciones.
\end{observacion}

\subsection{Informaci\'{o}n de Fisher}

Tenemos el score $s(\theta)$ --- una variable aleatoria
con media cero. Su varianza mide cu\'{a}nto fluctua el
score alrededor de cero, lo que a su vez mide cu\'{a}n
sensible es la distribuci\'{o}n al par\'{a}metro. Esa
varianza es la \textbf{informaci\'{o}n de Fisher}.

\begin{definicion}{Informaci\'{o}n de Fisher}{fisher_info}
	La \textbf{informaci\'{o}n de Fisher} sobre $\theta_0$
	contenida en una observaci\'{o}n $X$ es
	\[
	\mathcal{I}(\theta_0) =
	\mathrm{Var}_{\theta_0}(s(\theta_0))
	= E_{\theta_0}\!\left[
	\left(\frac{\partial}{\partial\theta}
	\ln p_{\theta_0}(X)\right)^2\right],
	\]
	donde la segunda igualdad usa que
	$E_{\theta_0}[s(\theta_0)] = 0$
	(Proposici\'{o}n~\ref{prop:score_media_cero}).
\end{definicion}

\begin{observacion}
	La informaci\'{o}n de Fisher tiene una forma alternativa
	\'{u}til bajo condiciones de regularidad:
	\[
	\mathcal{I}(\theta_0) = -E_{\theta_0}\!\left[
	\frac{\partial^2}{\partial\theta^2}
	\ln p_{\theta_0}(X)\right].
	\]
	Esta expresi\'{o}n --- la curvatura esperada negativa
	de la log-verosimilitud --- es frecuentemente m\'{a}s
	f\'{a}cil de calcular que la forma cuadr\'{a}tica.
	Interpretaci\'{o}n: si la log-verosimilitud tiene una
	curvatura muy pronunciada alrededor de $\theta_0$,
	los datos son muy informativos sobre $\theta_0$. Si
	la curvatura es suave, los datos dicen poco.
\end{observacion}

Esto lo podemos ilustrar mejor con un ejemplo:

\begin{ejemplo}{Informaci\'{o}n de Fisher para
		Bernoulli}{fisher_bernoulli}
	Sea $X \sim \mathrm{Bernoulli}(p)$ con $p \in (0,1)$.
	Calculamos $\mathcal{I}(p)$.
	
	La funci\'{o}n de masa es $p_p(x) = p^x(1-p)^{1-x}$
	para $x \in \{0,1\}$. La log-verosimilitud para una
	observaci\'{o}n es:
	\[
	\ln p_p(x) = x\ln p + (1-x)\ln(1-p).
	\]
	La funci\'{o}n score es:
	\[
	\frac{\partial}{\partial p}\ln p_p(x)
	= \frac{x}{p} - \frac{1-x}{1-p}
	= \frac{x - p}{p(1-p)}.
	\]
	La informaci\'{o}n de Fisher es la varianza del score:
	\[
	\mathcal{I}(p)
	= E_p\!\left[\left(\frac{X-p}{p(1-p)}\right)^2\right]
	= \frac{1}{p^2(1-p)^2}\,E_p[(X-p)^2]
	= \frac{\mathrm{Var}_p(X)}{p^2(1-p)^2}.
	\]
	Como $\mathrm{Var}_p(X) = p(1-p)$ para una Bernoulli
	(Proposici\'{o}n~\ref{prop:momentos_bern}):
	\[
	\mathcal{I}(p) = \frac{p(1-p)}{p^2(1-p)^2}
	= \frac{1}{p(1-p)}.
	\]
\end{ejemplo}

\begin{observacion}
	El resultado $\mathcal{I}(p) = 1/(p(1-p))$ tiene una
	interpretaci\'{o}n intuitiva clara. Cuando $p$ est\'{a}
	cerca de $0$ o de $1$, la funci\'{o}n $1/(p(1-p))$
	es muy grande --- los datos son muy informativos sobre
	$p$, porque casi siempre sale el mismo resultado y
	cualquier desviaci\'{o}n es muy reveladora. Cuando
	$p = 0.5$, la informaci\'{o}n es m\'{i}nima ---
	los resultados son m\'{a}xximamente impredecibles y
	cada observaci\'{o}n dice poco sobre el valor exacto
	de $p$.
	
	La cota de Cram\'{e}r-Rao para una muestra de
	tama\~{n}o $n$ es:
	\[
	\mathrm{Var}_p(\hat{p}) \geq
	\frac{1}{n\,\mathcal{I}(p)}
	= \frac{p(1-p)}{n}.
	\]
	Notemos que esta cota es exactamente la varianza de
	la media muestral $\bar{X}$ --- lo que anticipa que
	$\bar{X}$ es un estimador eficiente de $p$ en el
	sentido de Cram\'{e}r-Rao.
\end{observacion}

\begin{figure}[H]
	\centering
	\begin{tikzpicture}[scale=1.0]
		
		\draw[thick,->] (-0.2,0) -- (5.5,0)
		node[right] {$p$};
		\draw[thick,->] (0,-0.2) -- (0,5.0)
		node[above] {$\mathcal{I}(p)$};
		
		% Curva I(p) = 1/(p(1-p)) escalada
		% x en [0.1, 0.9] mapeado a [0.5, 4.5]
		% I(p) = 1/(p(1-p)), maximo en extremos
		% Escalamos: x_plot = p*5, y_plot = min(I(p)/2, 4.5)
		\draw[azul, very thick, domain=0.12:0.88, samples=200]
		plot ({5*\x}, {min(1/(5*\x*(1-\x)), 4.5)});
		
		% Mínimo en p=0.5: I(0.5) = 4, y_plot = 4/2 = 2
		\draw[rojo, dashed, thick] (2.5,0) --
		(2.5, {1/(5*0.5*0.5)});
		\draw[rojo, dashed, thick] (0, {1/(5*0.5*0.5)}) --
		(2.5, {1/(5*0.5*0.5)});
		\node[rojo, left, font=\small] at
		(0, {1/(5*0.5*0.5)}) {$4$};
		\node[rojo, below, font=\small] at (2.5,0) {$0.5$};
		
		% Etiquetas eje x
		\foreach \p/\lbl in {
			0.5/0.1, 1.0/0.2, 1.5/0.3, 2.0/0.4,
			3.0/0.6, 3.5/0.7, 4.0/0.8, 4.5/0.9}{
			\node[below, font=\scriptsize] at (\p, 0) {$\lbl$};
			\draw[thin] (\p, 0.05) -- (\p,-0.05);
		}
		
		% Anotaciones
		\node[azul, font=\small, align=center] at (1.2, 4.2)
		{alta informaci\'{o}n\\cerca de $0$};
		\draw[azul, ->, thin] (1.2, 3.8) -- (0.7, 3.0);
		
		\node[azul, font=\small, align=center] at (4.0, 4.2)
		{alta informaci\'{o}n\\cerca de $1$};
		\draw[azul, ->, thin] (4.0, 3.8) -- (4.4, 3.0);
		
		\node[rojo, font=\small, align=center] at (2.5, 1.2)
		{m\'{i}nimo en\\$p = 0.5$};
		
	\end{tikzpicture}
	\caption{Informaci\'{o}n de Fisher $\mathcal{I}(p) =
		1/(p(1-p))$ para $X \sim \mathrm{Bernoulli}(p)$. La
		informaci\'{o}n es m\'{a}xima cerca de los extremos $p=0$
		y $p=1$ --- donde los datos son m\'{a}s reveladores sobre
		$p$ --- y m\'{i}nima en $p=0.5$, donde la variable es
		m\'{a}xximamente impredecible.}
	\label{fig:fisher_bernoulli}
\end{figure}

\begin{observacion}
	Para una muestra aleatoria de tama\~{n}o $n$, la
	informaci\'{o}n de Fisher total es
	\[
	\mathcal{I}_n(\theta_0) = n\,\mathcal{I}(\theta_0).
	\]
	Es decir, la informaci\'{o}n se acumula linealmente
	con el tama\~{n}o muestral --- duplicar la muestra
	duplica la informaci\'{o}n disponible sobre $\theta_0$.
	Esta propiedad es consecuencia directa de la
	independencia de las observaciones.
\end{observacion}

\section{Teorema de Cram\'{e}r-Rao}
\label{sec:cramer_rao}

Tenemos ahora todas las herramientas necesarias para
responder las preguntas que dejamos abiertas al final de
la Secci\'{o}n~\ref{sec:propiedades_estimadores}: ¿existe
una cota absoluta para la varianza de cualquier estimador
insesgado? ¿De qu\'{e} depende? El Teorema de
Cram\'{e}r-Rao responde ambas con precisi\'{o}n.

Antes de enunciarlo, recordamos un resultado cl\'{a}sico
que usaremos en la demostraci\'{o}n.

\begin{observacion}
	La \textbf{desigualdad de Cauchy-Schwarz} establece
	que para dos variables aleatorias $U$ y $V$ con
	varianza finita:
	\[
	\mathrm{Cov}(U, V)^2 \leq
	\mathrm{Var}(U)\cdot\mathrm{Var}(V).
	\]
	Es la versi\'{o}n probabil\'{i}stica de un resultado
	general que vale en cualquier espacio con producto
	interno --- el mismo que se aplica en $\mathbb{R}^n$
	para vectores. Aqu\'{i} el producto interno entre
	dos variables aleatorias $U$ y $V$ es $E(UV)$, y la
	desigualdad dice que la covarianza no puede superar
	el producto de las desviaciones est\'{a}ndar.
\end{observacion}

\begin{teorema}{Teorema de Cram\'{e}r-Rao}{cramer_rao}
	Sea $X_1, \ldots, X_n$ una muestra aleatoria de
	$P_{\theta_0}$ y sea $\hat{\theta}$ un estimador
	insesgado de $\theta_0$. Bajo condiciones de
	regularidad,
	\[
	\mathrm{Var}_{\theta_0}(\hat{\theta})
	\geq \frac{1}{n\,\mathcal{I}(\theta_0)}.
	\]
	La cantidad $1/(n\,\mathcal{I}(\theta_0))$ se denomina
	\textbf{cota de Cram\'{e}r-Rao}.
\end{teorema}

\begin{proof}
	La demostraci\'{o}n tiene tres pasos.
	
	\medskip
	\textbf{Paso 1 --- Una identidad clave.}
	Mostraremos que
	\[
	\mathrm{Cov}_{\theta_0}(\hat{\theta},\,
	s(\theta_0)) = 1.
	\]
	Como $\hat{\theta}$ es insesgado,
	$E_{\theta_0}(\hat{\theta}) = \theta_0$ para todo
	$\theta_0 \in \Theta$. Diferenciando respecto a
	$\theta_0$ e intercambiando derivada e integral
	(condiciones de regularidad):
	\[
	\frac{\partial}{\partial\theta_0}
	\int \hat{\theta}(x_1,\ldots,x_n)\,
	\prod_{i=1}^n p_{\theta_0}(x_i)\,dx_i = 1.
	\]
	Aplicando la regla de la cadena al integrando y
	usando que $\frac{\partial}{\partial\theta_0}
	\ln L(\theta_0) = s(\theta_0)$:
	\[
	\int \hat{\theta}(x_1,\ldots,x_n)\cdot
	s(\theta_0)\cdot\prod_{i=1}^n
	p_{\theta_0}(x_i)\,dx_i = 1,
	\]
	es decir, $E_{\theta_0}[\hat{\theta}\cdot
	s(\theta_0)] = 1$. Como $E_{\theta_0}[s(\theta_0)]
	= 0$ por la
	Proposici\'{o}n~\ref{prop:score_media_cero}:
	\[
	\mathrm{Cov}_{\theta_0}(\hat{\theta},\,
	s(\theta_0))
	= E_{\theta_0}[\hat{\theta}\cdot s(\theta_0)]
	- E_{\theta_0}[\hat{\theta}]\cdot
	E_{\theta_0}[s(\theta_0)]
	= 1 - \theta_0 \cdot 0 = 1.
	\]
	
	\medskip
	\textbf{Paso 2 --- Desigualdad de Cauchy-Schwarz.}
	Aplicando la desigualdad de Cauchy-Schwarz a
	$\hat{\theta}$ y $s(\theta_0)$:
	\[
	\mathrm{Cov}_{\theta_0}(\hat{\theta},\,
	s(\theta_0))^2 \leq
	\mathrm{Var}_{\theta_0}(\hat{\theta})\cdot
	\mathrm{Var}_{\theta_0}(s(\theta_0)).
	\]
	
	\medskip
	\textbf{Paso 3 --- Combinar.}
	Sustituyendo el Paso 1 en el lado izquierdo y la
	definici\'{o}n de informaci\'{o}n de Fisher
	$\mathcal{I}_n(\theta_0) = n\,\mathcal{I}(\theta_0)
	= \mathrm{Var}_{\theta_0}(s(\theta_0))$ en el
	lado derecho:
	\[
	1 \leq \mathrm{Var}_{\theta_0}(\hat{\theta})
	\cdot n\,\mathcal{I}(\theta_0).
	\]
	Despejando:
	\[
	\mathrm{Var}_{\theta_0}(\hat{\theta}) \geq
	\frac{1}{n\,\mathcal{I}(\theta_0)}. \qedhere
	\]
\end{proof}

\begin{observacion}
	El teorema responde las tres preguntas abiertas:
	\begin{enumerate}[label=\textbf{\roman*)}]
		\item \textbf{¿Existe una cota absoluta?} Sí:
		$1/(n\,\mathcal{I}(\theta_0))$.
		\item \textbf{¿De qu\'{e} depende?} De la
		informaci\'{o}n de Fisher --- cuanto m\'{a}s
		informativa sea la muestra sobre $\theta_0$,
		menor es la cota y m\'{a}s preciso puede ser
		el estimador. Con m\'{a}s datos ($n$ mayor),
		la cota decrece como $1/n$.
		\item \textbf{¿Hay un estimador que la alcance?}
		Se dice que un estimador insesgado que alcanza
		la cota es \textbf{eficiente} en sentido
		absoluto. No siempre existe, pero cuando existe
		es el \'{o}ptimo indiscutible dentro de la
		clase de estimadores insesgados.
	\end{enumerate}
\end{observacion}

\begin{observacion}
	Cram\'{e}r-Rao justifica retroactivamente por qu\'{e}
	la varianza es el criterio correcto para comparar
	estimadores insesgados: no es una elecci\'{o}n
	arbitraria, sino la consecuencia de que existe una
	cota natural expresada en t\'{e}rminos de varianza.
	Comparar varianzas no es solo conveniente --- es lo
	\'{u}nico que tiene sentido dado este resultado.
\end{observacion}

\section{Por qu\'{e} estas propiedades y no otras}

Al definir insesgadez, eficiencia y consistencia en la
Secci\'{o}n~\ref{sec:propiedades_estimadores}, las
presentamos como criterios razonables para evaluar
estimadores. Pero ``razonable'' no es una justificaci\'{o}n
matem\'{a}tica. Ahora que tenemos el Teorema de
Cram\'{e}r-Rao, podemos responder con precisi\'{o}n por
qu\'{e} estas tres propiedades y no otras.

\subsection{Por qu\'{e} la insesgadez}

El Teorema de Cram\'{e}r-Rao establece una cota inferior
para la varianza de cualquier estimador
\textbf{insesgado}. Esto no es accidental --- la
demostraci\'{o}n us\'{o} de manera esencial que
$E_{\theta_0}(\hat{\theta}) = \theta_0$: fue precisamente
esa condici\'{o}n la que permiti\'{o} derivar la identidad
$\mathrm{Cov}_{\theta_0}(\hat{\theta}, s(\theta_0)) = 1$
en el Paso~1 de la demostraci\'{o}n.

Sin insesgadez, esa identidad no vale y la cota
desaparece. Esto tiene una consecuencia importante: si
permitimos estimadores sesgados, el problema de encontrar
el ``mejor'' estimador se vuelve mal planteado --- siempre
es posible reducir la varianza introduciendo sesgo, y no
hay un l\'{i}mite natural que diga cu\'{a}ndo parar.

La insesgadez no es una propiedad arbitraria. Es la
condici\'{o}n que hace que el problema de optimizaci\'{o}n
est\'{e} bien planteado: delimita la clase de estimadores
dentro de la cual existe un l\'{i}mite absoluto de
precisi\'{o}n, y por lo tanto dentro de la cual tiene
sentido hablar de un estimador \'{o}ptimo.

\subsection{Por qu\'{e} la eficiencia}

Dentro de la clase de estimadores insesgados,
Cram\'{e}r-Rao demuestra que la varianza satisface:
\[
\mathrm{Var}_{\theta_0}(\hat{\theta}) \geq
\frac{1}{n\,\mathcal{I}(\theta_0)}.
\]
Esto significa que la varianza tiene una cota inferior
\textbf{absoluta} --- no relativa a otro estimador, sino
fijada por la cantidad de informaci\'{o}n que los datos
contienen sobre $\theta_0$. Comparar varianzas entre
estimadores insesgados no es una elecci\'{o}n arbitraria:
es lo \'{u}nico que tiene sentido dado que existe ese
l\'{i}mite natural expresado en t\'{e}rminos de varianza.

M\'{a}s a\'{u}n, el teorema define impl\'{i}citamente
qu\'{e} significa ser \'{o}ptimo en sentido absoluto: un
estimador insesgado que alcanza la cota de Cram\'{e}r-Rao
es eficiente en el sentido m\'{a}s fuerte posible --- no
hay ning\'{u}n estimador insesgado que pueda superarlo,
independientemente de cu\'{a}n ingenioso sea su
dise\~{n}o.

\subsection{Por qu\'{e} la consistencia}

La cota de Cram\'{e}r-Rao es $1/(n\,\mathcal{I}(\theta_0))$
--- decrece como $1/n$ cuando el tama\~{n}o muestral
crece. Esto revela algo fundamental: la precisi\'{o}n
m\'{a}xima alcanzable mejora con m\'{a}s datos, y lo hace
a una tasa bien definida. Un buen estimador deber\'{i}a
aprovechar esa mejora --- su varianza deber\'{i}a decrecer
al menos tan r\'{a}pido como la cota.

La consistencia formaliza exactamente eso: exige que el
estimador converja a $\theta_0$ cuando $n \to \infty$.
Un estimador inconsistente no aprovecha la informaci\'{o}n
adicional que proveen m\'{a}s datos --- su error no
desaparece aunque la cota de Cram\'{e}r-Rao se acerque
a cero. Es, en ese sentido, un estimador que desperdicia
informaci\'{o}n.

La consistencia no es una condici\'{o}n arbitraria de
``mejora con m\'{a}s datos''. Es la condici\'{o}n m\'{i}nima
para que un estimador sea compatible con lo que
Cram\'{e}r-Rao garantiza que es posible lograr.

\subsection{Las tres propiedades como sistema coherente}

Las tres propiedades no son una lista arbitraria de
criterios deseables. Son un sistema coherente donde
cada una responde una pregunta precisa:

\begin{center}
	\begin{tabular}{lll}
		\toprule
		\textbf{Propiedad} & \textbf{Pregunta} &
		\textbf{Justificaci\'{o}n} \\
		\midrule
		Insesgadez &
		¿Apunta en la direcci\'{o}n correcta? &
		Delimita la clase donde existe \'{o}ptimo \\[4pt]
		Eficiencia &
		¿Aprovecha la informaci\'{o}n? &
		Cram\'{e}r-Rao fija el l\'{i}mite absoluto \\[4pt]
		Consistencia &
		¿Mejora con m\'{a}s datos? &
		La cota decrece como $1/n$ \\
		\bottomrule
	\end{tabular}
\end{center}

\medskip
Cram\'{e}r-Rao es el resultado que une las tres: muestra
que dentro de la clase de estimadores insesgados existe
un l\'{i}mite de precisi\'{o}n determinado por la
informaci\'{o}n de Fisher, que ese l\'{i}mite mejora con
m\'{a}s datos, y que un estimador \'{o}ptimo es aquel
que lo alcanza. Insesgadez, eficiencia y consistencia
son las tres caras de esa misma idea.

\section{La media muestral como estimador \'{o}ptimo}

A lo largo de este cap\'{i}tulo hemos construido un
marco te\'{o}rico completo: definimos qu\'{e} es un
estimador, establecimos las propiedades que debe cumplir
un buen estimador, y demostramos con Cram\'{e}r-Rao que
existe un l\'{i}mite absoluto de precisi\'{o}n determinado
por la informaci\'{o}n de Fisher. Ahora aplicamos todo ese
aparato al estimador m\'{a}s natural e importante de la
estad\'{i}stica: la \textbf{media muestral}
\[
\bar{X} = \frac{1}{n}\sum_{i=1}^n X_i.
\]
Demostraremos que $\bar{X}$ no es solo un estimador
conveniente --- es el estimador \'{o}ptimo dentro de la
clase insesgada, en el sentido de que alcanza la cota de
Cram\'{e}r-Rao.

\begin{proposicion}{Insesgadez de la media muestral}
	{xbar_insesgado}
	Sea $X_1, \ldots, X_n$ una muestra aleatoria de una
	poblaci\'{o}n con media $\mu$. Entonces
	\[
	E(\bar{X}) = \mu.
	\]
\end{proposicion}

\begin{proof}
	Por linealidad del valor esperado
	(Proposici\'{o}n~\ref{prop:ev_lineal}):
	\[
	E(\bar{X})
	= E\!\left(\frac{1}{n}\sum_{i=1}^n X_i\right)
	= \frac{1}{n}\sum_{i=1}^n E(X_i)
	= \frac{1}{n}\sum_{i=1}^n \mu
	= \frac{1}{n}\cdot n\mu
	= \mu. \qedhere
	\]
\end{proof}

\begin{observacion}
	La demostraci\'{o}n no asume ninguna distribuci\'{o}n
	espec\'{i}fica para la poblaci\'{o}n --- solo que
	$E(X_i) = \mu$. La insesgadez de $\bar{X}$ vale para
	cualquier familia param\'{e}trica con media finita.
\end{observacion}

\begin{proposicion}{Varianza de la media muestral}{var_xbar}
	Sea $X_1, \ldots, X_n$ una muestra aleatoria de una
	poblaci\'{o}n con varianza $\sigma^2$. Entonces
	\[
	\mathrm{Var}(\bar{X}) = \frac{\sigma^2}{n}.
	\]
\end{proposicion}

\begin{proof}
	Por la aditividad de la varianza bajo independencia
	(v\'{a}lida porque la muestra es i.i.d.):
	\[
	\mathrm{Var}(\bar{X})
	= \mathrm{Var}\!\left(\frac{1}{n}
	\sum_{i=1}^n X_i\right)
	= \frac{1}{n^2}\sum_{i=1}^n \mathrm{Var}(X_i)
	= \frac{1}{n^2}\cdot n\sigma^2
	= \frac{\sigma^2}{n}. \qedhere
	\]
\end{proof}

\begin{proposicion}{Eficiencia de $\bar{X}$ para la
		familia Normal}{xbar_eficiente}
	Sea $X_1, \ldots, X_n$ una muestra aleatoria de
	$N(\mu, \sigma^2)$ con $\sigma^2$ conocido. Entonces
	$\bar{X}$ alcanza la cota de Cram\'{e}r-Rao y es por
	lo tanto el estimador insesgado de $\mu$ con menor
	varianza posible.
\end{proposicion}

\begin{proof}
	Calculamos la informaci\'{o}n de Fisher para la familia
	$N(\mu, \sigma^2)$ con $\sigma^2$ conocido y $\theta
	= \mu$. La log-densidad es:
	\[
	\ln p_\mu(x) = -\frac{1}{2}\ln(2\pi\sigma^2)
	- \frac{(x-\mu)^2}{2\sigma^2}.
	\]
	La funci\'{o}n score es:
	\[
	\frac{\partial}{\partial\mu}\ln p_\mu(x)
	= \frac{x - \mu}{\sigma^2}.
	\]
	La informaci\'{o}n de Fisher para una observaci\'{o}n:
	\[
	\mathcal{I}(\mu)
	= E_\mu\!\left[\left(\frac{X-\mu}{\sigma^2}
	\right)^2\right]
	= \frac{1}{\sigma^4}\,\mathrm{Var}_\mu(X)
	= \frac{1}{\sigma^4}\cdot\sigma^2
	= \frac{1}{\sigma^2}.
	\]
	La cota de Cram\'{e}r-Rao para una muestra de
	tama\~{n}o $n$ es:
	\[
	\frac{1}{n\,\mathcal{I}(\mu)}
	= \frac{1}{n \cdot \frac{1}{\sigma^2}}
	= \frac{\sigma^2}{n}.
	\]
	Por la Proposici\'{o}n~\ref{prop:var_xbar},
	$\mathrm{Var}(\bar{X}) = \sigma^2/n$. Como la
	varianza de $\bar{X}$ coincide exactamente con la
	cota de Cram\'{e}r-Rao, $\bar{X}$ es el estimador
	insesgado de $\mu$ con menor varianza posible.
	\qedhere
\end{proof}

\begin{proposicion}{Consistencia de la media muestral}
	{xbar_consistente}
	Sea $X_1, \ldots, X_n$ una muestra aleatoria de una
	poblaci\'{o}n con media $\mu$ y varianza $\sigma^2 
	\infty$. Entonces $\bar{X}$ es consistente:
	\[
	\lim_{n\to\infty} P\!\left(|\bar{X} - \mu| >
	\varepsilon\right) = 0
	\qquad \text{para todo } \varepsilon > 0.
	\]
\end{proposicion}

\begin{proof}
	Aplicamos la desigualdad de Chebyshev al estimador
	$\bar{X}$. Como $E(\bar{X}) = \mu$ y
	$\mathrm{Var}(\bar{X}) = \sigma^2/n$:
	\[
	P\!\left(|\bar{X} - \mu| > \varepsilon\right)
	\leq \frac{\mathrm{Var}(\bar{X})}{\varepsilon^2}
	= \frac{\sigma^2}{n\varepsilon^2}.
	\]
	Cuando $n \to \infty$, el lado derecho converge a
	cero para todo $\varepsilon > 0$ fijo. Por lo tanto:
	\[
	\lim_{n\to\infty} P\!\left(|\bar{X} - \mu| >
	\varepsilon\right) = 0. \qedhere
	\]
\end{proof}

\begin{observacion}
	La consistencia de $\bar{X}$ es la \textbf{Ley D\'{e}bil
		de los Grandes N\'{u}meros} --- uno de los resultados
	m\'{a}s importantes de toda la probabilidad. Establece
	que el promedio de una muestra aleatoria converge en
	probabilidad a la media poblacional cuando el tama\~{n}o
	muestral crece. La demostraci\'{o}n via Chebyshev que
	presentamos aqu\'{i} es elemental pero completa ---
	no requiere ning\'{u}n resultado adicional m\'{a}s all\'{a}
	de lo que ya hemos desarrollado en este apunte.
\end{observacion}

Las cuatro proposiciones juntas establecen que $\bar{X}$
cumple todo lo que el marco te\'{o}rico de este cap\'{i}tulo
exige de un buen estimador --- y m\'{a}s: no solo es
insesgado, eficiente y consistente, sino que es
\textbf{\'{o}ptimo en sentido absoluto} para la familia
Normal, en el sentido de que ningún estimador insesgado
puede superarlo. Esto cierra el c\'{i}rculo: partimos de
propiedades definidas emp\'{i}ricamente, las justificamos
con Cram\'{e}r-Rao, y terminamos demostrando que el
estimador m\'{a}s natural de todos las satisface todas
--- y las satisface al m\'{a}ximo nivel posible.

\section{Distribuci\'{o}n muestral y Teorema del
	L\'{i}mite Central}

En la secci\'{o}n anterior demostramos que $\bar{X}$
tiene media $\mu$ y varianza $\sigma^2/n$. Pero conocer
los primeros dos momentos no es suficiente para calcular
probabilidades sobre $\bar{X}$ --- para eso necesitamos
conocer su \textbf{distribuci\'{o}n completa}. La
distribuci\'{o}n de probabilidad de un estad\'{i}stico
recibe un nombre propio, y determinarla para $\bar{X}$
es el objetivo de esta secci\'{o}n. La respuesta es el
resultado m\'{a}s importante de toda la probabilidad.

\begin{definicion}{Distribuci\'{o}n muestral}
	{dist_muestral}
	La \textbf{distribuci\'{o}n muestral} de un
	estad\'{i}stico $T = g(X_1, \ldots, X_n)$ es la
	distribuci\'{o}n de probabilidad de $T$ cuando
	$X_1, \ldots, X_n$ es una muestra aleatoria de
	$P_{\theta_0}$.
\end{definicion}

\begin{observacion}
	La distribuci\'{o}n muestral no es la distribuci\'{o}n
	de la poblaci\'{o}n --- es la distribuci\'{o}n del
	estad\'{i}stico calculado sobre la muestra. Para un
	mismo par\'{a}metro $\theta_0$, distintos estad\'{i}sticos
	tienen distintas distribuciones muestrales. La
	distribuci\'{o}n muestral de $\bar{X}$ responde la
	pregunta: si tomara muchas muestras de tama\~{n}o $n$
	y calculara $\bar{X}$ en cada una, ¿c\'{o}mo se
	distribuir\'{i}an esos valores?
\end{observacion}

Ya sabemos que $E(\bar{X}) = \mu$ y
$\mathrm{Var}(\bar{X}) = \sigma^2/n$ para cualquier
poblaci\'{o}n con media $\mu$ y varianza $\sigma^2$.
Pero ¿qu\'{e} \emph{forma} tiene la distribuci\'{o}n
muestral de $\bar{X}$? La respuesta depende de la
poblaci\'{o}n:

\begin{itemize}[leftmargin=2em]
	\item Si la poblaci\'{o}n es Normal, la distribuci\'{o}n
	de $\bar{X}$ es exactamente Normal para cualquier $n$.
	\item Si la poblaci\'{o}n \textbf{no} es Normal,
	el Teorema del L\'{i}mite Central garantiza que la
	distribuci\'{o}n de $\bar{X}$ se aproxima a una
	Normal cuando $n$ crece.
\end{itemize}

\begin{observacion}
	En el lenguaje formal de este cap\'{i}tulo, decir
	que \textbf{la poblaci\'{o}n es Normal} significa
	que la distribuci\'{o}n $P_{\theta_0}$ que gobierna
	la variable de inter\'{e}s pertenece a la familia
	param\'{e}trica $\mathcal{F} = \{N(\mu, \sigma^2) :
	\mu \in \mathbb{R},\, \sigma^2 > 0\}$, es decir,
	que cada observaci\'{o}n $X_i$ de la muestra
	aleatoria sigue una distribuci\'{o}n Normal. Cuando
	la poblaci\'{o}n \textbf{no} es Normal, $P_{\theta_0}$
	pertenece a otra familia --- Uniforme, Exponencial,
	Poisson, o cualquier otra --- y el TLC garantiza
	que $\bar{X}$ converge a la Normal de todas formas,
	independientemente de cu\'{a}l sea esa familia.
\end{observacion}

\begin{teorema}{Teorema del L\'{i}mite Central}{tlc}
	Sea $X_1, X_2, \ldots, X_n$ una muestra aleatoria
	con $E(X_i) = \mu$ y $\mathrm{Var}(X_i) = \sigma^2
	< \infty$. Se define
	\[
	Z_n \coloneqq \frac{\displaystyle\sum_{i=1}^n
		X_i - n\mu}{\sigma\sqrt{n}}.
	\]
	Entonces la funci\'{o}n de distribuci\'{o}n de $Z_n$
	converge hacia la funci\'{o}n de distribuci\'{o}n
	Normal est\'{a}ndar cuando $n \to \infty$:
	\[
	\lim_{n\to\infty} P(Z_n \leq z)
	= \Phi(z)
	= \int_{-\infty}^{z} \frac{1}{\sqrt{2\pi}}\,
	e^{-t^2/2}\,dt
	\qquad \text{para todo } z \in \mathbb{R}.
	\]
	Es equivalente escribirlo en funci\'{o}n de la media
	muestral $\bar{X}$:
	\[
	Z_n = \frac{\bar{X} - \mu}{\sigma/\sqrt{n}},
	\]
	puesto que ambas expresiones difieren solo en dividir
	numerador y denominador por $n$.
\end{teorema}

\begin{observacion}
	El teorema no dice nada sobre la distribuci\'{o}n de
	cada $X_i$ --- solo exige que existan su media y su
	varianza. La convergencia a la Normal ocurre
	independientemente de la forma de la distribuci\'{o}n
	poblacional. En la pr\'{a}ctica, la aproximaci\'{o}n
	$P(Z_n \leq z) \approx \Phi(z)$ es adecuada para
	$n \geq 30$, aunque este umbral depende de cu\'{a}n
	alejada est\'{e} la poblaci\'{o}n de la simetr\'{i}a.
\end{observacion}

\begin{observacion}
	La demostraci\'{o}n del TLC requiere el concepto de
	\textbf{funci\'{o}n caracter\'{i}stica}
	$\varphi_X(t) = E[e^{itX}]$ --- la transformada de
	Fourier de la distribuci\'{o}n de $X$ --- y el
	resultado de que la convergencia de funciones
	caracter\'{i}sticas implica convergencia en
	distribuci\'{o}n. Estas herramientas pertenecen al
	\'{a}mbito del an\'{a}lisis real y est\'{a}n fuera
	del alcance de este curso. Lo que s\'{i} podemos
	hacer es entender por qu\'{e} el resultado es
	verdadero intuitivamente, y verificarlo en casos
	particulares.
\end{observacion}

\begin{observacion}
	El TLC es la respuesta a la pregunta que dejamos
	abierta en el Tema~3 al estudiar la distribuci\'{o}n
	Normal: ¿por qu\'{e} la Normal aparece en fen\'{o}menos
	tan distintos como alturas, errores de medici\'{o}n
	y precios financieros? La respuesta es que en todos
	esos casos la variable observada es la suma --- o el
	promedio --- de muchos factores independientes
	peque\~{n}os. El TLC garantiza que esa suma converge
	a una Normal independientemente de la distribuci\'{o}n
	de los factores individuales. La Normal no aparece
	por coincidencia --- aparece porque es el atractor
	universal de las sumas de variables aleatorias
	independientes.
\end{observacion}

Lo anterior implica adem\'{a}s que la suma
$\sum_{i=1}^n X_i = n\bar{X}$ tambi\'{e}n se distribuye
aproximadamente Normal:
\[
\sum_{i=1}^n X_i \approx N(n\mu,\, n\sigma^2)
\qquad \text{para } n \geq 30.
\]

\begin{proposicion}{Caso exacto: poblaci\'{o}n Normal}
	{tlc_exacto}
	Si la poblaci\'{o}n es $N(\mu, \sigma^2)$, entonces
	para \textbf{cualquier} $n$:
	\[
	\bar{X} \sim N\!\left(\mu,\,\frac{\sigma^2}{n}
	\right)
	\qquad \text{y} \qquad
	\sum_{i=1}^n X_i \sim N(n\mu,\, n\sigma^2).
	\]
\end{proposicion}

\begin{proof}
	La suma de variables normales independientes es
	normal --- resultado que se demuestra usando funciones
	caracter\'{i}sticas o la propiedad de reproductividad
	de la distribuci\'{o}n Normal. Espec\'{i}ficamente,
	si $X_i \sim N(\mu, \sigma^2)$ i.i.d., entonces por
	la propiedad de reproductividad:
	\[
	\sum_{i=1}^n X_i \sim N(n\mu,\, n\sigma^2).
	\]
	Dividiendo por $n$ y usando las propiedades de la
	Normal bajo transformaciones afines
	(Proposici\'{o}n~\ref{prop:estandarizacion}):
	\[
	\bar{X} = \frac{1}{n}\sum_{i=1}^n X_i
	\sim N\!\left(\mu,\,\frac{\sigma^2}{n}\right).
	\qedhere
	\]
\end{proof}

\begin{observacion}
	La diferencia entre el caso exacto y el TLC es
	fundamental. Cuando la poblaci\'{o}n es Normal, la
	distribuci\'{o}n de $\bar{X}$ es exactamente Normal
	para cualquier $n$, incluso $n=1$. Cuando la
	poblaci\'{o}n no es Normal, la distribuci\'{o}n de
	$\bar{X}$ solo se \emph{aproxima} a la Normal, y
	la aproximaci\'{o}n mejora con $n$. En la pr\'{a}ctica,
	$n \geq 30$ suele ser suficiente para que la
	aproximaci\'{o}n sea adecuada, aunque esto depende
	de cu\'{a}n alejada est\'{e} la poblaci\'{o}n de
	la simetr\'{i}a.
\end{observacion}

\begin{figure}[H]
	\centering
	\begin{tikzpicture}[scale=0.95]
		
		% Mostramos la distribución de Xbar para distintos n
		% cuando la población es Uniforme(0,1): mu=0.5, sigma^2=1/12
		% Xbar ~ N(0.5, 1/(12n)) aproximadamente por TLC
		% sigma/sqrt(n) = 1/(sqrt(12)*sqrt(n)) = 0.2887/sqrt(n)
		
		\draw[thick,->] (-0.5,0) -- (10.5,0)
		node[right] {$\bar{x}$};
		\draw[thick,->] (0,-0.2) -- (0,5.5)
		node[above] {$f(\bar{x})$};
		
		% Escala: x en [0,1] mapeado a [0,10]
		% f escalada por 2 para visualización
		
		% n=1: Uniforme(0,1), densidad = 1, altura = 2
		\fill[azul!10] (0,0) rectangle (10,2);
		\draw[azul!40, very thick] (0,2) -- (10,2);
		\node[azul!60, font=\scriptsize] at (8.5, 2.3) {$n=1$};
		
		% n=5: Normal(0.5, 1/60), sigma=0.1291
		% sigma_plot = 0.1291*10 = 1.291
		% altura pico = 2/(1.291*sqrt(2pi)) * 2 = 0.617*2 = ...
		% usamos 0.3989/sigma_plot * escala
		% escala = 2*10 = 20 (para que área = 1 en plot)
		% altura = 0.3989/(0.1291) * 2 = 6.18 -- demasiado alto
		% Usemos escala = 1 y multiplicamos por 2
		\draw[azul!60, very thick, domain=0:10, samples=200]
		plot (\x, {2*0.3989/1.291*exp(-(\x-5)^2/(2*1.291^2))});
		\node[azul!70, font=\scriptsize] at (7.5, 0.9)
		{$n=5$};
		
		% n=10: sigma_plot = 0.2887/sqrt(10)*10 = 0.9129
		\draw[azul!80, very thick, domain=0:10, samples=200]
		plot (\x, {2*0.3989/0.9129*exp(-(\x-5)^2/(2*0.9129^2))});
		\node[azul!80, font=\scriptsize] at (6.8, 1.8)
		{$n=10$};
		
		% n=30: sigma_plot = 0.2887/sqrt(30)*10 = 0.5270
		\draw[azul, very thick, domain=0:10, samples=200]
		plot (\x, {2*0.3989/0.5270*exp(-(\x-5)^2/(2*0.5270^2))});
		\node[azul, font=\scriptsize] at (6.3, 3.2)
		{$n=30$};
		
		% n=100: sigma_plot = 0.2887/sqrt(100)*10 = 0.2887
		\draw[azul!80!black, very thick, domain=0:10, samples=200]
		plot (\x, {2*0.3989/0.2887*exp(-(\x-5)^2/(2*0.2887^2))});
		\node[azul!80!black, font=\scriptsize] at (6.0, 5.0)
		{$n=100$};
		
		% Línea en la media mu=0.5 -> x=5
		\draw[rojo, dashed, thick] (5,0) -- (5,5.3);
		\node[rojo, above, font=\small] at (5, 5.3) {$\mu = 0.5$};
		
		% Etiquetas eje x
		\foreach \x/\lbl in {0/0, 2.5/0.25, 5/0.5,
			7.5/0.75, 10/1}{
			\node[below, font=\scriptsize] at (\x,0) {$\lbl$};
			\draw[thin] (\x,0.05) -- (\x,-0.05);
		}
		
	\end{tikzpicture}
	\caption{Distribuci\'{o}n muestral de $\bar{X}$ para
		distintos tama\~{n}os de muestra $n$, cuando la poblaci\'{o}n
		es $U(0,1)$ --- una distribuci\'{o}n completamente no
		Normal (rect\'{a}ngulo). Para $n=1$, la distribuci\'{o}n
		de $\bar{X}$ es la misma Uniforme. A medida que $n$ crece,
		la distribuci\'{o}n se concentra alrededor de $\mu = 0.5$
		y adquiere la forma de campana gaussiana. Para $n \geq 30$
		la aproximaci\'{o}n Normal es adecuada.}
	\label{fig:tlc}
\end{figure}

El Teorema del L\'{i}mite Central es la respuesta
completa a la pregunta que dejamos abierta en el
Tema~3: ¿por qu\'{e} la distribuci\'{o}n Normal
aparece en fen\'{o}menos tan distintos como alturas,
errores de medici\'{o}n y precios financieros? La
respuesta no es que el mundo sea normal --- es que
el mundo est\'{a} lleno de fen\'{o}menos que son
sumas de muchos factores independientes peque\~{n}os,
y el TLC garantiza que esas sumas convergen a la
Normal sin importar la distribuci\'{o}n de cada
factor individual. La altura de una persona es la
suma de miles de factores gen\'{e}ticos y ambientales.
El error de medici\'{o}n es la suma de miles de
perturbaciones peque\~{n}as. El precio de una acci\'{o}n
es la acumulaci\'{o}n de miles de decisiones
independientes. En todos esos casos opera el TLC ---
y produce la campana gaussiana. En t\'{e}rminos
simples: \textbf{la distribuci\'{o}n Normal no es
	un supuesto sobre el mundo, es una consecuencia
	matem\'{a}tica inevitable de la independencia y
	la acumulaci\'{o}n.}

% ==========================================================
\chapter{Introducci\'{o}n a la inferencia estad\'{i}stica}

Todo el aparato que hemos construido --- variables
aleatorias, distribuciones, estimadores, Teorema del
L\'{i}mite Central --- ten\'{i}a un objetivo impl\'{i}cito:
aprender algo sobre una poblaci\'{o}n a partir de una
muestra. Hasta ahora ese aprendizaje tom\'{o} la forma
de una \textbf{estimaci\'{o}n puntual}: dado un par\'{a}metro
desconocido $\theta_0$, calculamos $\hat{\theta}$ como
nuestra mejor aproximaci\'{o}n.

Pero en la pr\'{a}ctica, la pregunta que enfrentan los
tomadores de decisiones rara vez es ``¿cu\'{a}nto vale
$\theta_0$?''. M\'{a}s frecuentemente es ``¿es $\theta_0$
mayor que cierto umbral?'', ``¿cambi\'{o} $\theta_0$
despu\'{e}s de una intervenci\'{o}n?'', o ``¿son
compatibles los datos con cierta afirmaci\'{o}n sobre
$\theta_0$?''. Esas son preguntas de \textbf{decisi\'{o}n},
no de estimaci\'{o}n.

La inferencia estad\'{i}stica provee dos herramientas
para responderlas:

\begin{center}
	\begin{tabular}{ll}
		\toprule
		\textbf{Herramienta} & \textbf{Pregunta que responde} \\
		\midrule
		Test de hip\'{o}tesis &
		¿Son los datos compatibles con cierta afirmaci\'{o}n
		sobre $\theta_0$? \\[4pt]
		Intervalo de confianza &
		¿En qu\'{e} rango de valores es plausible que
		est\'{e} $\theta_0$? \\
		\bottomrule
	\end{tabular}
\end{center}

\medskip
Ambas herramientas est\'{a}n \'{i}ntimamente relacionadas
--- un intervalo de confianza puede construirse a partir
de un test de hip\'{o}tesis y viceversa. Las estudiaremos
en orden.

\section{Test de hip\'{o}tesis}

Imaginemos que una empresa farmac\'{e}utica afirma que
su nuevo medicamento reduce la presi\'{o}n arterial media
a menos de 120 mmHg. Un regulador de salud quiere
verificar esa afirmaci\'{o}n. No puede medir a toda la
poblaci\'{o}n --- solo tiene acceso a una muestra. ¿C\'{o}mo
usa esa muestra para tomar una decisi\'{o}n?

La respuesta es el \textbf{test de hip\'{o}tesis}: un
procedimiento formal para decidir, sobre la base de
evidencia muestral, si una afirmaci\'{o}n sobre el
par\'{a}metro poblacional es compatible con los datos
o no.

\subsection{Hip\'{o}tesis nula e hip\'{o}tesis alternativa}

Todo test de hip\'{o}tesis parte de dos afirmaciones
contrapuestas sobre el par\'{a}metro $\theta_0$.

\begin{definicion}{Hip\'{o}tesis nula e hip\'{o}tesis
		alternativa}{hipotesis}
	La \textbf{hip\'{o}tesis nula} $H_0$ es la afirmaci\'{o}n
	que se supone cierta inicialmente --- la posici\'{o}n
	de referencia o \emph{statu quo}. La \textbf{hip\'{o}tesis
		alternativa} $H_1$ es la afirmaci\'{o}n contradictoria
	a $H_0$ que se desea detectar si los datos lo justifican.
	
	Si $\mu$ denota el par\'{a}metro de inter\'{e}s y
	$\mu_0$ es un valor espec\'{i}fico llamado
	\textbf{valor nulo}, la hip\'{o}tesis nula toma
	la forma
	\[
	H_0: \mu = \mu_0,
	\]
	y la hip\'{o}tesis alternativa tiene una de las
	tres formas:
	\begin{enumerate}[label=\textbf{\arabic*.}]
		\item $H_1: \mu \neq \mu_0$ \quad
		(prueba \textbf{de las dos colas})
		\item $H_1: \mu > \mu_0$ \quad
		(prueba \textbf{de cola superior})
		\item $H_1: \mu < \mu_0$ \quad
		(prueba \textbf{de cola inferior})
	\end{enumerate}
\end{definicion}

\begin{observacion}
	En el lenguaje formal del Cap\'{i}tulo~4, la
	hip\'{o}tesis nula $H_0: \mu = \mu_0$ afirma que
	el verdadero par\'{a}metro es $\theta_0 = \mu_0$.
	El test de hip\'{o}tesis pregunta: ¿son los datos
	compatibles con esa afirmaci\'{o}n, o constituyen
	evidencia suficiente para rechazarla?
\end{observacion}

\subsection{La l\'{o}gica del test: por qu\'{e} partimos
	de $H_0$}

La estructura del test de hip\'{o}tesis --- partir de
$H_0$ y buscar evidencia para rechazarla, nunca para
``probarla'' --- no es una convenci\'{o}n arbitraria.
Tiene ra\'{i}ces profundas en la filosof\'{i}a de la
ciencia, espec\'{i}ficamente en el \textbf{falsacionismo}
del fil\'{o}sofo austr\'{i}aco Karl Popper (1902--1994).

\subsubsection{El problema de la inducci\'{o}n}

Durante siglos, la ciencia opero bajo la l\'{o}gica
\textbf{inductiva}: observar muchos casos favorables
y concluir que una ley general es verdadera. Si veo
mil cisnes blancos, induzco que ``todos los cisnes son
blancos''. El problema --- se\~{n}alado ya por David
Hume en el siglo XVIII --- es que ninguna cantidad
finita de observaciones favorables puede \emph{probar}
una afirmaci\'{o}n universal. Siempre podr\'{i}a existir
un contraejemplo no observado a\'{u}n.

Y en efecto: cuando los europeos llegaron a Australia,
encontraron cisnes negros. Mil observaciones favorables
no bastaron.

\subsubsection{La soluci\'{o}n de Popper: falsabilidad}

Popper propuso invertir la l\'{o}gica. En lugar de
buscar confirmaciones, la ciencia deber\'{i}a buscar
\textbf{refutaciones}. Una afirmaci\'{o}n cient\'{i}fica
v\'{a}lida no es aquella que puede confirmarse, sino
aquella que puede \emph{falsarse} --- es decir, aquella
para la que existe al menos un experimento concebible
que podr\'{i}a demostrar que es falsa.

La asimetr\'{i}a es crucial:

\begin{itemize}[leftmargin=2em]
	\item \textbf{Confirmar} requiere infinitas
	observaciones favorables --- nunca se puede
	estar seguro.
	\item \textbf{Refutar} requiere una sola
	observaci\'{o}n contradictoria --- y eso s\'{i}
	es definitivo.
\end{itemize}

Por eso Popper argument\'{o} que la ciencia avanza
\textbf{descartando hip\'{o}tesis}, no confirm\'{a}ndolas.
Una teor\'{i}a que no puede refutarse en principio ---
que es compatible con cualquier observaci\'{o}n posible
--- no es cient\'{i}fica. Es infalsable, y por lo tanto
vac\'{i}a de contenido emp\'{i}rico.

\subsubsection{De Popper al test de hip\'{o}tesis}

El test de hip\'{o}tesis opera exactamente con esa
l\'{o}gica. La hip\'{o}tesis nula $H_0$ juega el papel
de la teor\'{i}a que queremos intentar refutar:

\begin{enumerate}[label=\textbf{\arabic*.}]
	\item \textbf{Asumimos $H_0$ como verdadera} ---
	el \emph{statu quo} que ponemos a prueba. No
	porque creamos que es verdadera, sino porque es
	la posici\'{o}n que la evidencia debe superar.
	
	\item \textbf{Calculamos qu\'{e} tan probable
		ser\'{i}a observar los datos que tenemos si
		$H_0$ fuera cierta}. Si esa probabilidad es
	muy peque\~{n}a, los datos son incompatibles
	con $H_0$.
	
	\item \textbf{Si la probabilidad es suficientemente
		peque\~{n}a, rechazamos $H_0$} --- al igual que
	Popper refuta una teor\'{i}a con un contraejemplo,
	rechazamos $H_0$ porque los datos son demasiado
	improbables bajo ella.
	
	\item \textbf{Si la probabilidad no es peque\~{n}a,
		no rechazamos $H_0$} --- pero esto \textbf{no
		significa que $H_0$ sea verdadera}. Solo significa
	que los datos no proporcionan evidencia suficiente
	para descartarla. Exactamente como en Popper: la
	ausencia de refutaci\'{o}n no es confirmaci\'{o}n.
\end{enumerate}

\begin{observacion}
	La distinci\'{o}n entre ``no rechazar $H_0$'' y
	``aceptar $H_0$'' es una de las m\'{a}s importantes
	--- y m\'{a}s frecuentemente violadas --- de toda
	la estad\'{i}stica aplicada. Un test que no rechaza
	$H_0$ no ha ``probado'' nada: solo ha constatado
	que la evidencia disponible es insuficiente para
	refutarla. Puede que $H_0$ sea falsa y la muestra
	simplemente demasiado peque\~{n}a para detectarlo.
	En la jerga estad\'{i}stica, esto se expresa
	diciendo que el test tiene \textbf{baja potencia}
	--- un concepto que estudiaremos m\'{a}s adelante.
	
	La confusi\'{o}n entre ambas es directamente
	popperiana: creer que la ausencia de refutaci\'{o}n
	implica confirmaci\'{o}n es el error inductivo que
	Popper pas\'{o} su vida criticando.
\end{observacion}

\begin{observacion}
	El falsacionismo tambi\'{e}n explica por qu\'{e}
	$H_0$ toma siempre la forma de una \textbf{igualdad}
	$\mu = \mu_0$ y no de una desigualdad. Una igualdad
	es m\'{a}s f\'{a}cil de refutar --- define un punto
	preciso en el espacio param\'{e}trico, y cualquier
	alejamiento suficiente de ese punto constituye
	evidencia en su contra. Una hip\'{o}tesis nula
	de la forma $\mu \geq \mu_0$ ser\'{i}a m\'{a}s
	dif\'{i}cil de tratar formalmente porque no define
	un \'{u}nico valor bajo el cual calcular
	probabilidades.
\end{observacion}

\subsection{Errores en el test de hip\'{o}tesis}

Al tomar una decisi\'{o}n basada en una muestra ---
y no en la poblaci\'{o}n completa --- siempre existe
la posibilidad de equivocarse. Hay exactamente dos
formas de cometer un error.

\begin{definicion}{Error de tipo I y error de tipo II}
	{errores}
	Un \textbf{error de tipo I} consiste en rechazar
	$H_0$ cuando $H_0$ es verdadera. Un \textbf{error
		de tipo II} consiste en no rechazar $H_0$ cuando
	$H_0$ es falsa.
	
	Sus probabilidades se denotan:
	\[
	P(\text{error tipo I}) = \alpha, \qquad
	P(\text{error tipo II}) = \beta.
	\]
	El par\'{a}metro $\alpha$ se denomina \textbf{nivel
		de significancia} del test.
\end{definicion}

El nivel de significancia $\alpha$ no es un
par\'{a}metro matem\'{a}tico neutro --- es una
\textbf{decisi\'{o}n humana} que refleja cu\'{a}n
conservador quiere ser el investigador. Fijar
$\alpha = 0{,}05$ significa: ``acepto un 5\% de
probabilidad de rechazar $H_0$ cuando es verdadera''.
Fijar $\alpha = 0{,}01$ es m\'{a}s exigente: solo
se tolera un 1\% de ese error. Cuanto m\'{a}s
peque\~{n}o sea $\alpha$, m\'{a}s dif\'{i}cil
ser\'{a} rechazar $H_0$ --- y por lo tanto
m\'{a}s s\'{o}lida ser\'{a} la evidencia requerida
para hacerlo.

\begin{table}[H]
	\centering
	\begin{tabular}{lccc}
		\toprule
		& & \multicolumn{2}{c}{\textbf{Realidad}} \\
		\cmidrule(lr){3-4}
		& & $H_0$ verdadera & $H_0$ falsa \\
		\midrule
		\multirow{2}{*}{\textbf{Decisi\'{o}n}} &
		No rechazar $H_0$ &
		\cellcolor{verde!20} Sin error &
		\cellcolor{rojo!20} Error tipo II ($\beta$) \\[6pt]
		& Rechazar $H_0$ &
		\cellcolor{rojo!20} Error tipo I ($\alpha$) &
		\cellcolor{verde!20} Sin error \\
		\bottomrule
	\end{tabular}
	\caption{Los cuatro posibles resultados de un test de
		hip\'{o}tesis. Solo dos de ellos constituyen errores.}
	\label{tab:errores}
\end{table}

\begin{observacion}
	Existe una \textbf{tensi\'{o}n fundamental} entre
	los dos tipos de error: reducir $\alpha$ ---
	ser m\'{a}s exigente para rechazar $H_0$ ---
	inevitablemente aumenta $\beta$, y viceversa.
	No es posible minimizar ambos simult\'{a}neamente
	con una muestra de tama\~{n}o fijo. La \'{u}nica
	forma de reducir ambos a la vez es \textbf{aumentar
		$n$} --- con m\'{a}s datos, la evidencia es m\'{a}s
	precisa y ambos errores pueden controlarse mejor.
	
	La convenci\'{o}n est\'{a}ndar en la pr\'{a}ctica
	es fijar $\alpha = 0{,}05$ o $\alpha = 0{,}01$.
	Esta convenci\'{o}n no tiene ning\'{u}n fundamento
	matem\'{a}tico especial --- es un est\'{a}ndar que
	se impuso en la pr\'{a}ctica cient\'{i}fica del
	siglo XX, en gran parte por influencia de Fisher.
	Muchos estad\'{i}sticos modernos critican el uso
	mec\'{a}nico de ese umbral, argumentando que la
	decisi\'{o}n deber\'{i}a depender del contexto:
	el costo relativo de los errores tipo~I y tipo~II
	var\'{i}a seg\'{u}n el problema. Esto refleja
	adem\'{a}s la asimetr\'{i}a popperiana: se es
	m\'{a}s cuidadoso con rechazar $H_0$
	err\'{o}neamente que con no detectar una
	alternativa verdadera.
\end{observacion}

% ==========================================================
\section{Pruebas de hip\'{o}tesis}

Con la arquitectura conceptual del test establecida
--- la l\'{o}gica de Popper, los errores de tipo~I
y~II, el nivel de significancia $\alpha$ --- el
siguiente paso es traducirla en un protocolo
operacional. La estad\'{i}stica aplicada exige que
el investigador fije todos sus criterios
\emph{antes} de observar los datos: de lo
contrario, siempre existe la tentaci\'{o}n de
ajustar las reglas para que los datos ``salgan
bien''. El procedimiento de cinco pasos convierte
la l\'{o}gica falsacionista en acciones concretas
y reproducibles.

% ----------------------------------------------------------
\subsection{El procedimiento de cinco pasos}

\begin{enumerate}[label=\textbf{Paso \arabic*.},
	leftmargin=*, itemsep=6pt]
	
	\item \textbf{Formular las hip\'{o}tesis nula
		y alternativa.} Se establece $H_0$ como la
	posici\'{o}n que se someter\'{a} a prueba ---
	siempre una igualdad $\mu = \mu_0$ --- y $H_1$
	como la alternativa que se aceptar\'{a} si
	$H_0$ es rechazada.
	
	La forma de $H_1$ no se determina
	matem\'{a}ticamente sino a partir del
	\textbf{contexto del problema}. La pregunta
	que gu\'{i}a la elecci\'{o}n es: \emph{¿qu\'{e}
		tipo de desviaci\'{o}n de $H_0$ ser\'{i}a
		relevante para el investigador?}
	\begin{itemize}[leftmargin=2em]
		\item Si el objetivo es verificar si un
		valor supuesto es correcto, sin direcci\'{o}n
		de sospecha prefijada $\longrightarrow$
		$H_1: \mu \neq \mu_0$, prueba de dos colas.
		\item Si hay razones previas para sospechar
		que el par\'{a}metro es \emph{mayor} que el
		supuesto $\longrightarrow$ $H_1: \mu > \mu_0$,
		prueba de cola superior.
		\item Si hay razones previas para sospechar
		que el par\'{a}metro es \emph{menor} que el
		supuesto $\longrightarrow$ $H_1: \mu < \mu_0$,
		prueba de cola inferior.
	\end{itemize}
	Esta elecci\'{o}n debe hacerse \textbf{antes
		de observar los datos} --- elegir $H_1$ despu\'{e}s
	de ver $\bar{x}$ invalida el nivel $\alpha$
	fijado en el Paso~2, pues equivale a ajustar
	las reglas del juego a la conveniencia.
	
	\item \textbf{Seleccionar el nivel de
		significancia $\alpha$.} Es la probabilidad
	m\'{a}xima de error tipo~I que se est\'{a}
	dispuesto a tolerar. Debe fijarse antes de
	observar los datos --- fijarla despu\'{e}s
	equivale a ajustar las reglas del juego a la
	conveniencia. Los valores convencionales son
	$\alpha = 0{,}05$ y $\alpha = 0{,}01$.
	
	\item \textbf{Identificar el estad\'{i}stico
		de prueba.} Se elige una funci\'{o}n de los
	datos cuya distribuci\'{o}n bajo $H_0$ es
	conocida. Es esa distribuci\'{o}n la que
	permite calcular probabilidades y construir
	la regi\'{o}n cr\'{i}tica. La elecci\'{o}n
	depende del par\'{a}metro bajo prueba y de
	los supuestos disponibles sobre la poblaci\'{o}n.
	
	\item \textbf{Formular la regla de decisi\'{o}n.}
	Se delimita la \textbf{regi\'{o}n cr\'{i}tica}
	--- el conjunto de valores del estad\'{i}stico
	que llevar\'{i}an a rechazar $H_0$ ---
	construida de modo que
	\[
	P\!\left(\text{estad\'{i}stico} \in
	\text{regi\'{o}n cr\'{i}tica}
	\mid H_0 \text{ verdadera}\right) = \alpha.
	\]
	Esta regla queda fijada antes de tomar la
	muestra.
	
	\item \textbf{Tomar la muestra y decidir.}
	Se calcula el valor observado del
	estad\'{i}stico y se verifica si cae dentro
	o fuera de la regi\'{o}n cr\'{i}tica. La
	decisi\'{o}n es binaria: \emph{rechazar
		$H_0$} o \emph{no rechazar $H_0$}.
	
\end{enumerate}

\begin{figure}[H]
	\centering
	\begin{tikzpicture}[
		node distance = 0.5cm and 0.15cm,
		box/.style = {
			rectangle, rounded corners = 3pt,
			draw = NavyBlue, fill = NavyBlue!10,
			text width = 1.95cm, align = center,
			minimum height = 1.4cm, font = \small
		},
		lbl/.style = {
			font = \footnotesize\itshape,
			text = gray, align = center
		},
		arr/.style = {->, thick, draw = NavyBlue!60}
		]
		\node[box] (p1) {Formular\\$H_0$ y $H_1$};
		\node[box, right = of p1] (p2) {Fijar\\$\alpha$};
		\node[box, right = of p2] (p3) {Elegir\\estad\'{\i}stico};
		\node[box, right = of p3] (p4) {Regi\'{o}n\\cr\'{\i}tica};
		\node[box, right = of p4] (p5) {Muestra\\y decisi\'{o}n};
		
		\draw[arr] (p1) -- (p2);
		\draw[arr] (p2) -- (p3);
		\draw[arr] (p3) -- (p4);
		\draw[arr] (p4) -- (p5);
		
		\node[lbl, below = 0.2cm of p1] {Paso 1};
		\node[lbl, below = 0.2cm of p2] {Paso 2};
		\node[lbl, below = 0.2cm of p3] {Paso 3};
		\node[lbl, below = 0.2cm of p4] {Paso 4};
		\node[lbl, below = 0.2cm of p5] {Paso 5};
		
	\end{tikzpicture}
	\caption{El procedimiento de cinco pasos para
		probar una hip\'{o}tesis.}
	\label{fig:cinco-pasos}
\end{figure}

% ----------------------------------------------------------
\subsection{Prueba $Z$: media con varianza conocida}

El procedimiento de cinco pasos es un esqueleto
vac\'{i}o hasta que se especifica un estad\'{i}stico
concreto en el Paso~3. La elecci\'{o}n del
estad\'{i}stico no es arbitraria --- depende de dos
factores: el \textbf{par\'{a}metro que se est\'{a}
	probando} y los \textbf{supuestos disponibles} sobre
la poblaci\'{o}n. El principio que gu\'{i}a la
elecci\'{o}n es siempre el mismo: se busca una
funci\'{o}n de los datos cuya distribuci\'{o}n bajo
$H_0$ sea completamente conocida, pues solo as\'{i}
es posible calcular probabilidades y construir la
regi\'{o}n cr\'{i}tica.

El primer caso --- y el m\'{a}s transparente
matem\'{a}ticamente --- es el de una poblaci\'{o}n
normal con varianza $\sigma^2$ conocida. Bajo
$H_0: \mu = \mu_0$, sabemos que $\bar{X}$ sigue
una distribuci\'{o}n normal con media $\mu_0$ y
varianza $\sigma^2/n$. Estandarizando $\bar{X}$
obtenemos un estad\'{i}stico con distribuci\'{o}n
$N(0,1)$ exacta bajo $H_0$ --- sin aproximaciones
ni supuestos adicionales. Ese estad\'{i}stico es
$Z_0$, y a la prueba que lo usa se la denomina
\textbf{prueba $Z$}.

\begin{proposicion}{Distribuci\'{o}n del estad\'{i}stico
		$Z_0$ bajo $H_0$}{dist-z0}
	Sea $X_1, X_2, \ldots, X_n$ una muestra
	aleatoria de una poblaci\'{o}n
	$N(\mu, \sigma^2)$ con $\sigma^2$ conocida.
	Bajo la hip\'{o}tesis nula $H_0: \mu = \mu_0$,
	el estad\'{i}stico
	\[
	Z_0 = \frac{\bar{X} - \mu_0}{\sigma /
		\sqrt{n}}
	\]
	sigue una distribuci\'{o}n normal est\'{a}ndar:
	\[
	Z_0 \sim N(0, 1).
	\]
\end{proposicion}

\begin{proof}
	Por la proposici\'{o}n~\ref{prop:tlc_exacto},
	la media muestral satisface
	\[
	\bar{X} \sim N\!\left(\mu,
	\frac{\sigma^2}{n}\right).
	\]
	Bajo $H_0: \mu = \mu_0$, esto se convierte en
	\[
	\bar{X} \sim N\!\left(\mu_0,
	\frac{\sigma^2}{n}\right).
	\]
	Aplicando la estandarizaci\'{o}n de la
	distribuci\'{o}n normal --- si
	$X \sim N(\mu, \sigma^2)$ entonces
	$(X - \mu)/\sigma \sim N(0,1)$ --- con
	$\mu = \mu_0$ y $\sigma^2$ reemplazada por
	$\sigma^2/n$:
	\[
	Z_0 = \frac{\bar{X} - \mu_0}{\sigma /
		\sqrt{n}} \sim N(0, 1). \qedhere
	\]
\end{proof}

Una vez que $Z_0$ est\'{a} calculado, la decisi\'{o}n
de rechazar o no $H_0$ depende de la forma de la
hip\'{o}tesis alternativa $H_1$. En todos los casos
la regi\'{o}n cr\'{i}tica se construye imponiendo
la misma condici\'{o}n: que su probabilidad bajo
$H_0$ sea exactamente $\alpha$. Como $Z_0 \sim
N(0,1)$ bajo $H_0$, esto se traduce en encontrar
el umbral $c$ tal que la probabilidad de que $Z_0$
caiga en la regi\'{o}n de rechazo sea $\alpha$.

\medskip
\textbf{Cola superior} ($H_1: \mu > \mu_0$).
Valores grandes de $Z_0$ son evidencia a favor
de $H_1$, por lo que la regi\'{o}n cr\'{i}tica
est\'{a} en la cola derecha. Imponemos:
\[
P(Z_0 > c \mid H_0) = \alpha
\implies 1 - \Phi(c) = \alpha
\implies \Phi(c) = 1 - \alpha
\implies c = z_{1-\alpha}.
\]
Se rechaza $H_0$ si $Z_0 > z_{1-\alpha}$.

\medskip
\textbf{Cola inferior} ($H_1: \mu < \mu_0$).
Valores peque\~{n}os de $Z_0$ son evidencia a
favor de $H_1$, por lo que la regi\'{o}n
cr\'{i}tica est\'{a} en la cola izquierda.
Por simetr\'{i}a de la $N(0,1)$:
\[
P(Z_0 < c \mid H_0) = \alpha
\implies \Phi(c) = \alpha
\implies c = -z_{1-\alpha}.
\]
Se rechaza $H_0$ si $Z_0 < -z_{1-\alpha}$.

\medskip
\textbf{Dos colas} ($H_1: \mu \neq \mu_0$).
Tanto valores muy grandes como muy peque\~{n}os
de $Z_0$ son evidencia a favor de $H_1$. El nivel
$\alpha$ se reparte en partes iguales entre ambas
colas --- $\alpha/2$ en cada una --- para mantener
la probabilidad total de rechazo bajo $H_0$
igual a $\alpha$:
\[
P(|Z_0| > c \mid H_0) = \alpha
\implies P(Z_0 > c) + P(Z_0 < -c) = \alpha
\implies 2(1 - \Phi(c)) = \alpha
\implies c = z_{1-\alpha/2}.
\]
Se rechaza $H_0$ si $|Z_0| > z_{1-\alpha/2}$.

\begin{table}[H]
	\centering
	\begin{tabular}{lll}
		\toprule
		\textbf{Hip\'{o}tesis alternativa $H_1$} &
		\textbf{Regi\'{o}n cr\'{i}tica} &
		\textbf{Tipo de prueba} \\
		\midrule
		$\mu > \mu_0$ &
		$Z_0 > z_{1-\alpha}$ &
		Cola superior \\[4pt]
		$\mu < \mu_0$ &
		$Z_0 < -z_{1-\alpha}$ &
		Cola inferior \\[4pt]
		$\mu \neq \mu_0$ &
		$|Z_0| > z_{1-\alpha/2}$ &
		Dos colas \\
		\bottomrule
	\end{tabular}
	\caption{Regiones de rechazo para la prueba $Z$
		seg\'{u}n la hip\'{o}tesis alternativa.}
	\label{tab:regiones-z}
\end{table}

\begin{definicion}{Valor cr\'{i}tico}{valor-critico}
	El \textbf{valor cr\'{i}tico} de un test es el
	umbral $c$ que delimita la regi\'{o}n cr\'{i}tica.
	Se obtiene despejando $c$ de la condici\'{o}n
	\[
	P(\text{rechazo} \mid H_0) = \alpha,
	\]
	que es exactamente lo que hicimos en las
	derivaciones anteriores. Su valor depende de
	$\alpha$ y del tipo de prueba:
	\begin{itemize}[leftmargin=2em]
		\item Cola superior: $c = z_{1-\alpha}$,
		pues $\Phi(z_{1-\alpha}) = 1-\alpha$.
		\item Cola inferior: $c = -z_{1-\alpha}$,
		por simetr\'{i}a de la $N(0,1)$.
		\item Dos colas: $c = z_{1-\alpha/2}$,
		pues el $\alpha$ se reparte en $\alpha/2$
		en cada cola.
	\end{itemize}
	Los valores cr\'{i}ticos m\'{a}s usados en la
	pr\'{a}ctica son:
	\[
	z_{0{,}95} = 1{,}645 \quad (\alpha = 0{,}10),
	\qquad
	z_{0{,}975} = 1{,}960 \quad (\alpha = 0{,}05),
	\]
	\[
	z_{0{,}99} = 2{,}326 \quad (\alpha = 0{,}01).
	\]
\end{definicion}

% --- figura cola superior ---
\begin{figure}[H]
	\centering
	\begin{tikzpicture}
		\begin{axis}[
			width = 10cm, height = 5cm,
			domain = -3.8:3.8,
			samples = 200,
			axis lines = left,
			xlabel = {$z$},
			ylabel = {},
			xtick = {1.645},
			xticklabels = {$z_{1-\alpha}$},
			ytick = \empty,
			ymin = 0, ymax = 0.45,
			axis line style = {draw = gray!60},
			tick style = {draw = gray!60},
			every axis x label/.style = {
				at = {(axis cs: 3.8, 0)},
				anchor = west, font = \small
			}
			]
			% curva normal
			\addplot[thick, NavyBlue, smooth]
			{exp(-x^2/2) / sqrt(2*pi)};
			
			% region critica cola superior
			\addplot[
			fill = red!25, opacity = 0.8,
			draw = none, domain = 1.645:3.8
			] {exp(-x^2/2) / sqrt(2*pi)} \closedcycle;
			
			% linea vertical
			\addplot[dashed, gray, thick]
			coordinates {(1.645, 0) (1.645, 0.105)};
			
			% etiqueta region critica
			\node[font=\small, text=red!70!black]
			at (axis cs: 2.5, 0.15) {$\alpha$};
			
		\end{axis}
	\end{tikzpicture}
	\caption{Regi\'{o}n cr\'{i}tica para la prueba de
		cola superior: se rechaza $H_0$ si
		$Z_0 > z_{1-\alpha}$.}
	\label{fig:cola-superior}
\end{figure}

% --- figura cola inferior ---
\begin{figure}[H]
	\centering
	\begin{tikzpicture}
		\begin{axis}[
			width = 10cm, height = 5cm,
			domain = -3.8:3.8,
			samples = 200,
			axis lines = left,
			xlabel = {$z$},
			ylabel = {},
			xtick = {-1.645},
			xticklabels = {$-z_{1-\alpha}$},
			ytick = \empty,
			ymin = 0, ymax = 0.45,
			axis line style = {draw = gray!60},
			tick style = {draw = gray!60},
			every axis x label/.style = {
				at = {(axis cs: 3.8, 0)},
				anchor = west, font = \small
			}
			]
			\addplot[thick, NavyBlue, smooth]
			{exp(-x^2/2) / sqrt(2*pi)};
			
			\addplot[
			fill = red!25, opacity = 0.8,
			draw = none, domain = -3.8:-1.645
			] {exp(-x^2/2) / sqrt(2*pi)} \closedcycle;
			
			\addplot[dashed, gray, thick]
			coordinates {(-1.645, 0) (-1.645, 0.105)};
			
			\node[font=\small, text=red!70!black]
			at (axis cs: -2.5, 0.15) {$\alpha$};
			
		\end{axis}
	\end{tikzpicture}
	\caption{Regi\'{o}n cr\'{i}tica para la prueba de
		cola inferior: se rechaza $H_0$ si
		$Z_0 < -z_{1-\alpha}$.}
	\label{fig:cola-inferior}
\end{figure}

% --- figura dos colas ---
\begin{figure}[H]
	\centering
	\begin{tikzpicture}
		\begin{axis}[
			width = 10cm, height = 5cm,
			domain = -3.8:3.8,
			samples = 200,
			axis lines = left,
			xlabel = {$z$},
			ylabel = {},
			xtick = {-1.96, 1.96},
			xticklabels = {$-z_{1-\alpha/2}$,
				$z_{1-\alpha/2}$},
			ytick = \empty,
			ymin = 0, ymax = 0.45,
			axis line style = {draw = gray!60},
			tick style = {draw = gray!60},
			every axis x label/.style = {
				at = {(axis cs: 3.8, 0)},
				anchor = west, font = \small
			}
			]
			\addplot[thick, NavyBlue, smooth]
			{exp(-x^2/2) / sqrt(2*pi)};
			
			\addplot[
			fill = red!25, opacity = 0.8,
			draw = none, domain = -3.8:-1.96
			] {exp(-x^2/2) / sqrt(2*pi)} \closedcycle;
			
			\addplot[
			fill = red!25, opacity = 0.8,
			draw = none, domain = 1.96:3.8
			] {exp(-x^2/2) / sqrt(2*pi)} \closedcycle;
			
			\addplot[dashed, gray, thick]
			coordinates {(-1.96, 0) (-1.96, 0.058)};
			\addplot[dashed, gray, thick]
			coordinates {(1.96, 0) (1.96, 0.058)};
			
			\node[font=\small, text=red!70!black]
			at (axis cs: -2.7, 0.12) {$\alpha/2$};
			\node[font=\small, text=red!70!black]
			at (axis cs: 2.7, 0.12) {$\alpha/2$};
			
		\end{axis}
	\end{tikzpicture}
	\caption{Regi\'{o}n cr\'{i}tica para la prueba de
		dos colas: se rechaza $H_0$ si
		$|Z_0| > z_{1-\alpha/2}$.}
	\label{fig:dos-colas}
\end{figure}

\begin{observacion}
	La notaci\'{o}n $z_{1-\alpha}$ denota el
	percentil $1 - \alpha$ de la distribuci\'{o}n
	normal est\'{a}ndar, es decir, el valor tal que
	\[
	P(Z \leq z_{1-\alpha}) = 1 - \alpha,
	\qquad Z \sim N(0,1).
	\]
	Por simetr\'{i}a de la distribuci\'{o}n normal
	est\'{a}ndar se tiene adem\'{a}s que
	$z_\alpha = -z_{1-\alpha}$, lo que permite
	expresar la regi\'{o}n cr\'{i}tica de cola
	inferior en t\'{e}rminos del mismo percentil.
	Los valores m\'{a}s usados en la pr\'{a}ctica
	son:
	\[
	z_{0{,}95} = 1{,}645, \qquad
	z_{0{,}975} = 1{,}960, \qquad
	z_{0{,}99} = 2{,}326.
	\]
\end{observacion}

% ----------------------------------------------------------
\subsection{Uso de la tabla normal est\'{a}ndar}

La funci\'{o}n $\Phi(z) = P(Z \leq z)$ no tiene
f\'{o}rmula cerrada --- su valor se obtiene
num\'{e}ricamente, en la pr\'{a}ctica mediante
tablas estad\'{i}sticas o software. La tabla
de la distribuci\'{o}n normal est\'{a}ndar
organiza los valores de $\Phi(z)$ para distintos
valores de $z$: las filas corresponden a la parte
entera y primer decimal de $z$, y las columnas
al segundo decimal.

La tabla se usa en dos direcciones:

\medskip
\textbf{Direcci\'{o}n directa.} Dado el valor
observado $z_0$, obtener $\Phi(z_0)$. Se busca
la fila correspondiente a la parte entera y
primer decimal de $z_0$, y la columna del segundo
decimal. Por ejemplo, para $z_0 = 1{,}96$: fila
$1{,}9$, columna $0{,}06$, valor
$\Phi(1{,}96) = 0{,}9750$. Esto se usa para
calcular el valor $p$ --- por ejemplo, para cola
superior: $p = 1 - \Phi(z_0) = 1 - 0{,}9750
= 0{,}0250$.

\medskip
\textbf{Direcci\'{o}n inversa.} Dada una
probabilidad $1 - \alpha$, obtener el percentil
$z_{1-\alpha}$. Se busca el valor $1 - \alpha$
dentro de la tabla y se lee el $z$ correspondiente.
Por ejemplo, para $\alpha = 0{,}05$ en cola
superior: se busca $0{,}9500$ en la tabla y se
lee $z_{0{,}95} = 1{,}645$. Esto se usa para
obtener los valores cr\'{i}ticos.

\begin{observacion}
	Por simetr\'{i}a de la $N(0,1)$, si
	necesitas $\Phi(-z)$ y la tabla solo
	cubre valores positivos de $z$, usa:
	\[
	\Phi(-z) = 1 - \Phi(z).
	\]
	Por ejemplo, $\Phi(-1{,}96) = 1 -
	\Phi(1{,}96) = 1 - 0{,}9750 = 0{,}0250$.
\end{observacion}

\begin{ejemplo}{Prueba $Z$ para la media: cl\'{i}nica
		Buena Salud}{clinica}
	
	Una cl\'{i}nica supone que el tiempo medio que los
	pacientes pasan en sala de espera es $\mu_0 = 9$
	minutos. Para verificar ese supuesto se toma una
	muestra aleatoria de $n = 120$ personas, obteniendo
	una media muestral de $\bar{x} = 10{,}2$ minutos.
	Los tiempos se distribuyen normalmente con varianza
	conocida $\sigma^2 = 16\ \text{min}^2$.
	
	La pregunta es: ¿es el supuesto de la cl\'{i}nica
	compatible con los datos observados, o la
	discrepancia entre $\bar{x} = 10{,}2$ y
	$\mu_0 = 9$ es demasiado grande para ser atribuida
	al azar del muestreo? Aplicamos el procedimiento
	de cinco pasos.
	
	\medskip
	\textbf{Paso 1. Hip\'{o}tesis.}
	
	La cl\'{i}nica quiere verificar si la media es
	exactamente $9$ minutos --- no tiene una direcci\'{o}n
	prefijada de desviaci\'{o}n. Por lo tanto la prueba
	es de dos colas:
	\[
	H_0: \mu = 9 \qquad \text{vs} \qquad
	H_1: \mu \neq 9.
	\]
	
	\medskip
	\textbf{Paso 2. Nivel de significancia.}
	
	Se fija $\alpha = 0{,}05$.
	
	\medskip
	\textbf{Paso 3. Estad\'{i}stico de prueba.}
	
	La poblaci\'{o}n es normal y $\sigma^2 = 16$ es
	conocida, por lo que corresponde usar la prueba
	$Z$. El estad\'{i}stico es:
	\[
	Z_0 = \frac{\bar{X} - \mu_0}{\sigma/\sqrt{n}}
	= \frac{\bar{X} - 9}{4/\sqrt{120}}.
	\]
	
	\medskip
	\textbf{Paso 4. Regi\'{o}n cr\'{i}tica.}
	
	La prueba es de dos colas, por lo que el valor
	cr\'{i}tico es $c = z_{1-\alpha/2}$
	(Definici\'{o}n~\ref{def:valor-critico}). Este
	valor no es una elecci\'{o}n arbitraria: es el
	\'{u}nico umbral que garantiza que la probabilidad
	de cometer un error tipo~I sea exactamente
	$\alpha$. En efecto, bajo $H_0$ se tiene
	$Z_0 \sim N(0,1)$, y para que
	$P(|Z_0| > c \mid H_0) = 0{,}05$ se necesita:
	\[
	2(1 - \Phi(c)) = 0{,}05
	\implies \Phi(c) = 0{,}975
	\implies c = z_{0{,}975} = 1{,}960.
	\]
	Cualquier otro valor de $c$ dar\'{i}a una
	probabilidad de error tipo~I distinta de $0{,}05$,
	violando el nivel de significancia fijado en el
	Paso~2. La regi\'{o}n cr\'{i}tica es por lo tanto:
	\[
	|Z_0| > 1{,}960,
	\]
	es decir, se rechaza $H_0$ si el estad\'{i}stico
	cae en alguna de las dos colas sombreadas en la
	Figura~\ref{fig:dos-colas}. En caso contrario
	--- si $-1{,}960 \leq Z_0 \leq 1{,}960$ --- no
	se rechaza $H_0$.
	
	\medskip
	\textbf{Paso 5. Muestra y decisi\'{o}n.}
	
	Con los datos observados $\bar{x} = 10{,}2$:
	\[
	z_0 = \frac{10{,}2 - 9}{4/\sqrt{120}}
	= \frac{1{,}2}{0{,}3651}
	= 3{,}286.
	\]
	Como $|z_0| = 3{,}286 > 1{,}960$, el valor
	observado cae dentro de la regi\'{o}n cr\'{i}tica.
	
	\medskip
	\textbf{Conclusi\'{o}n.} Se rechaza $H_0$ con
	nivel de significancia $\alpha = 0{,}05$. Los
	datos proporcionan evidencia suficiente para
	concluir que el tiempo medio en sala de espera
	no es $9$ minutos --- la discrepancia observada
	es demasiado grande para ser atribuida al azar
	del muestreo.
	
\end{ejemplo}

% ----------------------------------------------------------

\section{El valor $p$}

El procedimiento de regi\'{o}n cr\'{i}tica tiene
una limitaci\'{o}n pr\'{a}ctica: la decisi\'{o}n
depende del $\alpha$ que se fij\'{o} en el Paso~2.
Si dos investigadores analizan los mismos datos
con $\alpha = 0{,}05$ y $\alpha = 0{,}01$
respectivamente, podr\'{i}an llegar a conclusiones
distintas --- y ninguno de los dos sabr\'{i}a
cu\'{a}n ``cerca'' estuvo el estad\'{i}stico del
umbral de rechazo.

El \textbf{valor $p$} resuelve esto. En lugar de
preguntar ``¿cae $Z_0$ en la regi\'{o}n cr\'{i}tica
para este $\alpha$?'', pregunta algo m\'{a}s
informativo: ¿cu\'{a}l es el nivel de significancia
m\'{i}nimo para el cual estos datos llevar\'{i}an
a rechazar $H_0$? Equivalentemente, el valor $p$
mide cu\'{a}n incompatibles son los datos con
$H_0$: un valor $p$ peque\~{n}o indica que los
datos observados ser\'{i}an muy improbables si
$H_0$ fuera cierta.

\begin{definicion}{Valor $p$}{valor-p}
	El \textbf{valor $p$} es la probabilidad,
	bajo $H_0$, de obtener un valor del
	estad\'{i}stico de prueba tan extremo o
	m\'{a}s extremo que el valor observado $z_0$.
	Formalmente:
	\[
	p = P(Z_0 \in R_{z_0} \mid H_0),
	\]
	donde $R_{z_0}$ denota la regi\'{o}n del
	espacio muestral que es al menos tan extrema
	como $z_0$ en la direcci\'{o}n de $H_1$.
	Un valor $p$ peque\~{n}o indica que los datos
	observados ser\'{i}an muy improbables si $H_0$
	fuera cierta.
\end{definicion}

La intuici\'{o}n es simple. Imagina que lanzas
una moneda 100 veces y obtienes 95 caras. El
valor $p$ ser\'{i}a la probabilidad de obtener
95 o m\'{a}s caras \emph{si la moneda fuera
	justa}. Esa probabilidad es \'{i}nfima ---
lo que sugiere fuertemente que la moneda no
es justa. En general:
\begin{itemize}[leftmargin=2em]
	\item $p$ \textbf{grande} $\to$ lo que
	observaste no ser\'{i}a raro bajo $H_0$.
	No hay conflicto con $H_0$.
	\item $p$ \textbf{peque\~{n}o} $\to$ lo que
	observaste ser\'{i}a rar\'{i}simo bajo $H_0$.
	Los datos son incompatibles con $H_0$.
\end{itemize}

\begin{proposicion}{F\'{o}rmulas del valor $p$
		para la prueba $Z$}{formulas-p}
	Bajo $H_0$, con $Z_0 \sim N(0,1)$ y valor
	observado $z_0$, el valor $p$ se calcula como:
	\[
	\text{valor } p = \begin{cases}
		1 - \Phi(z_0)      & \text{cola superior}
		\; (H_1: \mu > \mu_0), \\[4pt]
		\Phi(z_0)          & \text{cola inferior}
		\; (H_1: \mu < \mu_0), \\[4pt]
		2(1-\Phi(|z_0|))   & \text{dos colas}
		\; (H_1: \mu \neq \mu_0).
	\end{cases}
	\]
\end{proposicion}

\begin{proof}
	En todos los casos aplicamos directamente la
	Definici\'{o}n~\ref{def:valor-p}: el valor $p$
	es la probabilidad, bajo $H_0$, de obtener un
	estad\'{i}stico tan extremo o m\'{a}s extremo
	que el observado $z_0$.
	
	\medskip
	\textbf{Cola superior.} Valores de $Z_0$
	mayores que $z_0$ son evidencia a\'{u}n m\'{a}s
	fuerte contra $H_0$ que el propio $z_0$. Por
	lo tanto:
	\[
	p = P(Z \geq z_0 \mid H_0)
	= 1 - P(Z < z_0 \mid H_0)
	= 1 - \Phi(z_0).
	\]
	
	\medskip
	\textbf{Cola inferior.} Valores de $Z_0$
	menores que $z_0$ son evidencia a\'{u}n m\'{a}s
	fuerte contra $H_0$. Por lo tanto:
	\[
	p = P(Z \leq z_0 \mid H_0)
	= \Phi(z_0).
	\]
	
	\medskip
	\textbf{Dos colas.} Tanto valores muy grandes
	como muy peque\~{n}os de $Z_0$ son evidencia
	contra $H_0$. La regi\'{o}n ``tan extrema o
	m\'{a}s'' que $z_0$ es $\{|Z| \geq |z_0|\}$.
	Por simetr\'{i}a de la $N(0,1)$:
	\[
	p = P(|Z| \geq |z_0| \mid H_0)
	= P(Z \geq |z_0|) + P(Z \leq -|z_0|)
	= 2\,P(Z \geq |z_0|)
	= 2(1 - \Phi(|z_0|)). \qedhere
	\]
\end{proof}

La regla de decisi\'{o}n es simple: se rechaza
$H_0$ si $p \leq \alpha$. Esto no es un nuevo
criterio --- es el mismo criterio de la
regi\'{o}n cr\'{i}tica, expresado de forma
distinta.

\begin{proposicion}{Equivalencia entre valor $p$
		y regi\'{o}n cr\'{i}tica}{equiv-p}
	Para la prueba $Z$, los siguientes criterios
	son equivalentes:
	\begin{enumerate}[label=\textit{\roman*.}]
		\item $Z_0$ cae en la regi\'{o}n cr\'{i}tica.
		\item $\text{valor } p \leq \alpha$.
	\end{enumerate}
\end{proposicion}

\begin{proof}
	Demostramos el caso de cola superior; los
	dem\'{a}s son an\'{a}logos. La regi\'{o}n
	cr\'{i}tica es $Z_0 > z_{1-\alpha}$, y el
	valor $p$ es $1 - \Phi(z_0)$. Como $\Phi$
	es estrictamente creciente:
	\[
	z_0 > z_{1-\alpha}
	\iff \Phi(z_0) > \Phi(z_{1-\alpha})
	= 1 - \alpha
	\iff 1 - \Phi(z_0) < \alpha
	\iff p < \alpha.
	\]
	An\'{a}logamente para cola inferior y dos
	colas. \qedhere
\end{proof}

\begin{corolario}{Rechazo para cualquier $\alpha
		\geq p$}{rechazo-alfa}
	Se rechaza $H_0$ para todo nivel de
	significancia $\alpha \geq p$, y no se
	rechaza para todo $\alpha < p$.
\end{corolario}

\begin{proof}
	Sea $\alpha^*$ un nivel de significancia
	arbitrario. Por la
	Proposici\'{o}n~\ref{prop:equiv-p}, rechazar
	$H_0$ con nivel $\alpha^*$ es equivalente
	a que $p \leq \alpha^*$. Por lo tanto:
	\[
	\alpha^* \geq p
	\iff p \leq \alpha^*
	\iff \text{se rechaza } H_0
	\text{ con nivel } \alpha^*.
	\]
	An\'{a}logamente, $\alpha^* < p \iff$ no
	se rechaza $H_0$ con nivel $\alpha^*$.
	\qedhere
\end{proof}

\begin{observacion}
	Esto hace al valor $p$ m\'{a}s informativo
	que la simple decisi\'{o}n binaria de la
	regi\'{o}n cr\'{i}tica: en lugar de decir
	``se rechaza con $\alpha = 0{,}05$'', el
	valor $p$ dice ``se rechazar\'{i}a para
	cualquier $\alpha \geq 0{,}001$'', lo que
	permite al lector juzgar la solidez de la
	evidencia con su propio criterio.
\end{observacion}

\begin{observacion}
	El valor $p$ \textbf{no} es la probabilidad
	de que $H_0$ sea verdadera, ni la probabilidad
	de haber cometido un error. Es la probabilidad
	de observar datos tan extremos como los
	obtenidos, \emph{asumiendo} que $H_0$ es
	verdadera. Esta distinci\'{o}n es una de las
	m\'{a}s frecuentemente mal interpretadas en
	la estad\'{i}stica aplicada.
\end{observacion}

Volviendo al ejemplo de la cl\'{i}nica, calculemos
el valor $p$ correspondiente al estad\'{i}stico
observado $z_0 = 3{,}286$.

\medskip
\textbf{Continuaci\'{o}n del ejemplo
	\ref{ej:clinica}.}

Como la prueba es de dos colas:
\[
p = 2(1 - \Phi(|z_0|))
= 2(1 - \Phi(3{,}286))
= 2(1 - 0{,}9995)
= 2 \times 0{,}0005
= 0{,}001.
\]
Como $p = 0{,}001 \leq 0{,}05 = \alpha$, se
rechaza $H_0$ --- la misma conclusi\'{o}n que
obtuvimos con la regi\'{o}n cr\'{i}tica, como
garantiza la Proposici\'{o}n~\ref{prop:equiv-p}.

Adem\'{a}s, el valor $p$ es m\'{a}s informativo
que la simple decisi\'{o}n de rechazo: indica
que $H_0$ ser\'{i}a rechazada para cualquier
nivel de significancia $\alpha \geq 0{,}001$
--- un resultado extremadamente s\'{o}lido.

% ==========================================================
\section{Prueba $t$: media con varianza desconocida}

La prueba $Z$ resuelve el problema de testear
$\mu$ cuando $\sigma^2$ es conocida. Pero en la
pr\'{a}ctica, $\sigma^2$ casi nunca se conoce
--- si se conociera, en muchos casos tambi\'{e}n
se conocer\'{i}a $\mu$. El escenario realista es
que ambos par\'{a}metros sean desconocidos y
deban inferirse a partir de los datos.

El paso natural es reemplazar $\sigma$ por un
estimador construido desde la muestra. Pero esa
sustituci\'{o}n tiene un costo: el estad\'{i}stico
resultante ya no sigue una $N(0,1)$ bajo $H_0$,
sino una distribuci\'{o}n con colas m\'{a}s
pesadas --- la $t$-Student. Para construir ese
estimador correctamente, primero necesitamos
definirlo con precisi\'{o}n.

% ----------------------------------------------------------
\subsection{La varianza muestral}

La varianza poblacional $\sigma^2$ mide la
dispersi\'{o}n de la poblaci\'{o}n respecto a
su media $\mu$. Cuando $\mu$ es desconocida,
no podemos calcular $\sigma^2$ directamente
--- pero s\'{i} podemos estimarla usando la
dispersi\'{o}n de la muestra respecto a
$\bar{X}$.

La candidata natural ser\'{i}a promediar las
desviaciones cuadr\'{a}ticas con $n$ en el
denominador. Pero eso produce un estimador
sesgado --- sistem\'{a}ticamente subestima
$\sigma^2$. La correcci\'{o}n es dividir por
$n-1$, y la raz\'{o}n de esa correcci\'{o}n
emerge directamente de la demostraci\'{o}n
del insesgamiento.

\begin{definicion}{Varianza muestral y
		desviaci\'{o}n est\'{a}ndar muestral}
	{varianza-muestral}
	Sea $X_1, X_2, \ldots, X_n$ una muestra
	aleatoria. La \textbf{varianza muestral}
	se define como:
	\[
	S^2 = \frac{1}{n-1}
	\sum_{i=1}^{n}(X_i - \bar{X})^2.
	\]
	La \textbf{desviaci\'{o}n est\'{a}ndar
		muestral} es $S = \sqrt{S^2}$.
\end{definicion}

\begin{proposicion}{Insesgamiento de $S^2$}
	{insesgamiento-s2}
	Sea $X_1, X_2, \ldots, X_n$ una muestra
	aleatoria de una poblaci\'{o}n con media
	$\mu$ y varianza $\sigma^2$. Entonces:
	\[
	E(S^2) = \sigma^2.
	\]
	Es decir, $S^2$ es un estimador insesgado
	de $\sigma^2$.
\end{proposicion}

\begin{proof}
	Partimos expandiendo la suma definitoria
	de $S^2$. Notemos que:
	\[
	\sum_{i=1}^{n}(X_i - \bar{X})^2
	= \sum_{i=1}^{n} X_i^2
	- n\bar{X}^2.
	\]
	Tomando esperanza a ambos lados:
	\[
	E\!\left[\sum_{i=1}^{n}
	(X_i - \bar{X})^2\right]
	= \sum_{i=1}^{n} E(X_i^2)
	- n\,E(\bar{X}^2).
	\]
	Usando que $E(X^2) = \text{Var}(X) +
	[E(X)]^2$ para cualquier variable
	aleatoria:
	\[
	E(X_i^2) = \sigma^2 + \mu^2,
	\qquad
	E(\bar{X}^2) = \text{Var}(\bar{X})
	+ [E(\bar{X})]^2
	= \frac{\sigma^2}{n} + \mu^2.
	\]
	Sustituyendo:
	\[
	E\!\left[\sum_{i=1}^{n}
	(X_i - \bar{X})^2\right]
	= n(\sigma^2 + \mu^2)
	- n\!\left(\frac{\sigma^2}{n}
	+ \mu^2\right)
	= n\sigma^2 + n\mu^2
	- \sigma^2 - n\mu^2
	= (n-1)\sigma^2.
	\]
	Por lo tanto:
	\[
	E(S^2) = \frac{1}{n-1}\,
	E\!\left[\sum_{i=1}^{n}
	(X_i - \bar{X})^2\right]
	= \frac{(n-1)\sigma^2}{n-1}
	= \sigma^2. \qedhere
	\]
\end{proof}

\begin{observacion}
	La demostraci\'{o}n revela por qu\'{e}
	el denominador es $n-1$ y no $n$: es
	exactamente el factor necesario para
	cancelar el $(n-1)$ que aparece al
	calcular $E\!\left[\sum(X_i -
	\bar{X})^2\right]$. Si se dividiera
	por $n$, el estimador ser\'{i}a
	$\frac{n-1}{n}\sigma^2 < \sigma^2$
	--- sistem\'{a}ticamente sesgado
	hacia abajo. La raz\'{o}n intuitiva
	es que $\bar{X}$ ya usa los datos
	para aproximar $\mu$, consumiendo
	un grado de libertad --- de ah\'{i}
	el $n-1$.
\end{observacion}

% ----------------------------------------------------------
\subsection{La distribuci\'{o}n $t$-Student}

Ahora que tenemos $S^2$ bien definido, podemos
construir el estad\'{i}stico an\'{a}logo a $Z_0$
pero con $S$ en lugar de $\sigma$. La pregunta
inmediata es: ¿qu\'{e} distribuci\'{o}n sigue
ese estad\'{i}stico bajo $H_0$?

Ya no puede ser $N(0,1)$, porque $S$ es una
variable aleatoria --- introduce variabilidad
adicional respecto a $Z_0$. Esa variabilidad
extra produce una distribuci\'{o}n con
\textbf{colas m\'{a}s pesadas} que la normal:
asigna m\'{a}s probabilidad a valores extremos.

La consecuencia pr\'{a}ctica es concreta: los
valores cr\'{i}ticos de esa distribuci\'{o}n
son m\'{a}s grandes que los de la $N(0,1)$
para el mismo $\alpha$. Por ejemplo, para
$\alpha = 0{,}05$ en una prueba de dos colas:
\begin{center}
	\begin{tabular}{lc}
		\toprule
		\textbf{Distribuci\'{o}n} &
		\textbf{Valor cr\'{i}tico} \\
		\midrule
		$N(0,1)$            & $1{,}960$ \\
		$t$ con 30 g.l.     & $2{,}042$ \\
		$t$ con 10 g.l.     & $2{,}228$ \\
		$t$ con 5 g.l.      & $2{,}571$ \\
		\bottomrule
	\end{tabular}
\end{center}
Es decir, con muestras peque\~{n}as se necesita
evidencia m\'{a}s extrema para rechazar $H_0$
--- lo que tiene sentido, porque la estimaci\'{o}n
de $\sigma^2$ es menos confiable. A medida que
$n \to \infty$, $S \to \sigma$ y la distribuci\'{o}n
converge a $N(0,1)$.

Esa distribuci\'{o}n es la \textbf{$t$-Student},
descubierta por William Sealy Gosset en 1908,
quien la public\'{o} bajo el seud\'{o}nimo
``Student'' porque su empleador --- la cervecer\'{i}a
Guinness --- no permit\'{i}a publicaciones a
sus empleados.

\begin{definicion}{Distribuci\'{o}n
		$t$-Student}{t-student}
	Se dice que una variable aleatoria $T$
	sigue una \textbf{distribuci\'{o}n
		$t$-Student con $\nu$ grados de libertad},
	y se escribe $T \sim t_{\nu}$, si su
	funci\'{o}n de densidad es:
	\[
	f(t) = \frac{\Gamma\!\left(
		\frac{\nu+1}{2}\right)}
	{\sqrt{\nu\pi}\;
		\Gamma\!\left(\frac{\nu}{2}\right)}
	\left(1 + \frac{t^2}{\nu}
	\right)^{-(\nu+1)/2},
	\qquad t \in \mathbb{R},
	\]
	donde $\Gamma$ es la funci\'{o}n Gamma.
	La distribuci\'{o}n $t_\nu$ es sim\'{e}trica
	respecto a cero, tiene colas m\'{a}s pesadas
	que la $N(0,1)$, y converge a $N(0,1)$
	cuando $\nu \to \infty$.
\end{definicion}

\begin{observacion}
	La notaci\'{o}n $t_{n-1,\,\alpha}$ denota
	el valor tal que
	\[
	P(T > t_{n-1,\,\alpha}) = \alpha,
	\qquad T \sim t_{n-1}.
	\]
	Es decir, $t_{n-1,\,\alpha}$ es el percentil
	$1-\alpha$ de la distribuci\'{o}n $t$ con
	$n-1$ grados de libertad. Por simetr\'{i}a,
	$-t_{n-1,\,\alpha}$ es el percentil $\alpha$.
\end{observacion}

\begin{figure}[H]
	\centering
	\begin{tikzpicture}
		\begin{axis}[
			width = 10cm, height = 6cm,
			domain = -4:4,
			samples = 200,
			axis lines = left,
			xlabel = {$t$},
			ylabel = {},
			ytick = \empty,
			ymin = 0, ymax = 0.45,
			axis line style = {draw = gray!60},
			tick style = {draw = gray!60},
			legend pos = north east,
			legend style = {font = \small},
			every axis x label/.style = {
				at = {(axis cs: 4, 0)},
				anchor = west, font = \small
			}
			]
			% N(0,1)
			\addplot[thick, NavyBlue, smooth]
			{exp(-x^2/2) / sqrt(2*pi)};
			\addlegendentry{$N(0,1)$}
			
% t con 10 g.l.
\addplot[thick, rojo, smooth, dashed]
{0.3891 * (1 + x^2/10)^(-5.5)};
\addlegendentry{$t_{10}$}

% t con 3 g.l.
\addplot[thick, verde!70!black, smooth,
dotted]
{0.3676 * (1 + x^2/3)^(-2)};
\addlegendentry{$t_{3}$}
			
		\end{axis}
	\end{tikzpicture}
	\caption{Comparaci\'{o}n entre la $N(0,1)$ y
		la distribuci\'{o}n $t$-Student con distinto
		n\'{u}mero de grados de libertad. A menor
		$\nu$, las colas son m\'{a}s pesadas y el
		pico central m\'{a}s bajo.}
	\label{fig:t-student}
\end{figure}

% ----------------------------------------------------------
\subsection{El estad\'{i}stico $T_0$ y las
	regiones de rechazo}

Con la distribuci\'{o}n $t$-Student establecida,
podemos construir el estad\'{i}stico de prueba
para el caso de varianza desconocida. La idea
es la misma que en la prueba $Z$: estandarizar
$\bar{X}$ bajo $H_0$, pero ahora usando $S$
en lugar de $\sigma$.

\begin{definicion}{Estad\'{i}stico $T_0$}
	{estadistico-t0}
	Sea $X_1, X_2, \ldots, X_n$ una muestra
	aleatoria de una poblaci\'{o}n
	$N(\mu, \sigma^2)$ con $\sigma^2$
	desconocida. Bajo $H_0: \mu = \mu_0$,
	el \textbf{estad\'{i}stico de prueba}
	es:
	\[
	T_0 = \frac{\bar{X} - \mu_0}
	{S / \sqrt{n}}.
	\]
\end{definicion}

\begin{proposicion}{Distribuci\'{o}n de $T_0$
		bajo $H_0$}{dist-t0}
	Bajo $H_0: \mu = \mu_0$, el
	estad\'{i}stico $T_0$ sigue una
	distribuci\'{o}n $t$-Student con
	$n-1$ grados de libertad:
	\[
	T_0 \sim t_{n-1}.
	\]
\end{proposicion}

\begin{proof}
	Reescribimos $T_0$ multiplicando y
	dividiendo por $\sigma$:
	\[
	T_0 = \frac{\bar{X} - \mu_0}
	{S/\sqrt{n}}
	= \frac{(\bar{X} - \mu_0)/
		(\sigma/\sqrt{n})}
	{S/\sigma}
	= \frac{Z_0}{S/\sigma},
	\]
	donde $Z_0 = \frac{\bar{X} - \mu_0}
	{\sigma/\sqrt{n}} \sim N(0,1)$ bajo
	$H_0$ por la
	Proposici\'{o}n~\ref{prop:dist-z0}.
	
	\medskip
	Reescribimos el denominador en
	t\'{e}rminos de la varianza muestral:
	\[
	\frac{S}{\sigma}
	= \sqrt{\frac{S^2}{\sigma^2}}
	= \sqrt{\frac{(n-1)S^2/\sigma^2}
		{n-1}}.
	\]
	
	\begin{observacion}
		Se puede demostrar que, cuando
		la poblaci\'{o}n es normal,
		\[
		\frac{(n-1)S^2}{\sigma^2}
		\sim \chi^2_{n-1},
		\]
		donde $\chi^2_{n-1}$ denota la
		\textbf{distribuci\'{o}n
			chi-cuadrado con $n-1$ grados
			de libertad} --- una distribuci\'{o}n
		continua de soporte positivo que
		surge naturalmente al sumar
		cuadrados de variables normales
		est\'{a}ndar independientes. El
		desarrollo formal de la
		distribuci\'{o}n $\chi^2$ est\'{a}
		fuera del alcance de este curso.
	\end{observacion}
	
	Por lo tanto:
	\[
	\frac{S}{\sigma}
	= \sqrt{\frac{\chi^2_{n-1}}{n-1}}.
	\]
	Sustituyendo en $T_0$:
	\[
	T_0 = \frac{Z_0}{\sqrt{\chi^2_{n-1}
			/(n-1)}}.
	\]
	
	\begin{observacion}
		Se puede demostrar adem\'{a}s que,
		cuando la poblaci\'{o}n es normal,
		$\bar{X}$ y $S^2$ son
		\textbf{independientes} --- resultado
		no trivial cuya demostraci\'{o}n
		tambi\'{e}n est\'{a} fuera del
		alcance del curso.
	\end{observacion}
	
	Con la independencia de $Z_0$ y
	$\chi^2_{n-1}$ garantizada, el cociente
	anterior es exactamente la
	\textbf{definici\'{o}n} de una variable
	$t$-Student: si $Z \sim N(0,1)$ y
	$V \sim \chi^2_\nu$ son independientes,
	entonces
	\[
	T = \frac{Z}{\sqrt{V/\nu}}
	\sim t_\nu.
	\]
	Aplicando con $\nu = n-1$:
	\[
	T_0 = \frac{Z_0}{\sqrt{\chi^2_{n-1}
			/(n-1)}} \sim t_{n-1}. \qedhere
	\]
\end{proof}

Con $T_0 \sim t_{n-1}$ bajo $H_0$, las
regiones de rechazo se construyen de forma
an\'{a}loga a la prueba $Z$, reemplazando
los percentiles de la $N(0,1)$ por los
de la $t_{n-1}$.

\begin{table}[H]
	\centering
	\begin{tabular}{lll}
		\toprule
		\textbf{Hip\'{o}tesis alternativa
			$H_1$} &
		\textbf{Regi\'{o}n cr\'{i}tica} &
		\textbf{Tipo de prueba} \\
		\midrule
		$\mu > \mu_0$ &
		$T_0 > t_{n-1,\,\alpha}$ &
		Cola superior \\[4pt]
		$\mu < \mu_0$ &
		$T_0 < -t_{n-1,\,\alpha}$ &
		Cola inferior \\[4pt]
		$\mu \neq \mu_0$ &
		$|T_0| > t_{n-1,\,\alpha/2}$ &
		Dos colas \\
		\bottomrule
	\end{tabular}
	\caption{Regiones de rechazo para la
		prueba $t$ seg\'{u}n la hip\'{o}tesis
		alternativa.}
	\label{tab:regiones-t}
\end{table}

\begin{ejemplo}{Prueba $t$ para la media:
		latas de pi\~{n}a}{latas}
	
	Una persona sospecha que las latas de pi\~{n}a
	de 16 onzas se est\'{a}n llenando en exceso.
	El departamento de control de calidad toma una
	muestra aleatoria de $n = 15$ latas, obteniendo
	una media muestral de $\bar{x} = 16{,}3$ onzas
	y una desviaci\'{o}n est\'{a}ndar muestral de
	$s = 0{,}6$ onzas. Los pesos se distribuyen
	normalmente y $\sigma^2$ es desconocida.
	
	N\'{o}tese que la sospecha tiene una
	\textbf{direcci\'{o}n prefijada} --- se cree
	que las latas se est\'{a}n llenando \emph{en
		exceso}, no simplemente que el peso es distinto
	de 16. Eso determina directamente la forma de
	$H_1$.
	
	\medskip
	\textbf{Paso 1. Hip\'{o}tesis.}
	
	La sospecha es que el llenado es excesivo,
	es decir, que $\mu > 16$. Por lo tanto la
	prueba es de cola superior:
	\[
	H_0: \mu = 16 \qquad \text{vs} \qquad
	H_1: \mu > 16.
	\]
	
	\medskip
	\textbf{Paso 2. Nivel de significancia.}
	
	El problema no especifica $\alpha$, por lo
	que adoptamos el valor convencional
	$\alpha = 0{,}05$.
	
	\medskip
	\textbf{Paso 3. Estad\'{i}stico de prueba.}
	
	La poblaci\'{o}n es normal y $\sigma^2$ es
	desconocida, por lo que corresponde usar la
	prueba $t$. El estad\'{i}stico es:
	\[
	T_0 = \frac{\bar{X} - \mu_0}
	{S/\sqrt{n}}
	= \frac{\bar{X} - 16}
	{0{,}6/\sqrt{15}}.
	\]
	
	\medskip
	\textbf{Paso 4. Regi\'{o}n cr\'{i}tica.}
	
	Como la prueba es de cola superior, el valor
	cr\'{i}tico es $t_{n-1,\,\alpha} =
	t_{14,\,0{,}05}$. Este es el \'{u}nico umbral
	que garantiza que $P(T_0 > c \mid H_0) =
	0{,}05$ cuando $T_0 \sim t_{14}$:
	\[
	P(T_0 > c \mid H_0) = 0{,}05
	\implies c = t_{14,\,0{,}05} = 1{,}761.
	\]
	Se rechaza $H_0$ si $T_0 > 1{,}761$.
	
	\medskip
	\textbf{Paso 5. Muestra y decisi\'{o}n.}
	
	Con los datos observados $\bar{x} = 16{,}3$,
	$s = 0{,}6$, $n = 15$:
	\[
	t_0 = \frac{16{,}3 - 16}
	{0{,}6/\sqrt{15}}
	= \frac{0{,}3}{0{,}1549}
	= 1{,}936.
	\]
	Como $t_0 = 1{,}936 > 1{,}761$, el valor
	observado cae dentro de la regi\'{o}n
	cr\'{i}tica.
	
	\medskip
	\textbf{Conclusi\'{o}n.} Se rechaza $H_0$
	con nivel de significancia $\alpha = 0{,}05$.
	Los datos proporcionan evidencia suficiente
	para concluir que las latas se est\'{a}n
	llenando en exceso --- la discrepancia
	observada es demasiado grande para ser
	atribuida al azar del muestreo.
	
\end{ejemplo}

% ==========================================================
\section{Estimaci\'{o}n por intervalos de
	confianza}

El test de hip\'{o}tesis responde una pregunta
binaria: ¿es $\mu = \mu_0$ o no? Pero esa
pregunta asume que ya tenemos un valor $\mu_0$
espec\'{i}fico para testear. En muchos casos
pr\'{a}cticos no tenemos ese valor --- simplemente
queremos saber \textbf{en qu\'{e} rango} est\'{a}
$\mu$.

La \textbf{estimaci\'{o}n por intervalos de
	confianza} responde esa pregunta. En lugar de
decir ``$\mu$ vale exactamente $\hat{\mu}$''
--- lo cual es casi con certeza falso ---
construimos un intervalo $[l, u]$ que contiene
a $\mu$ con una probabilidad prefijada
$1 - \alpha$. Esa probabilidad se denomina
\textbf{coeficiente de confianza}.

La idea clave es que el intervalo es aleatorio
--- depende de los datos --- pero el par\'{a}metro
$\mu$ es fijo. Lo que tiene probabilidad
$1 - \alpha$ no es que ``$\mu$ est\'{e} en el
intervalo'', sino que ``el procedimiento de
construcci\'{o}n produzca un intervalo que
contenga a $\mu$''.

% ----------------------------------------------------------
\subsection{Intervalo de confianza con varianza
	conocida}

Cuando $\sigma^2$ es conocida, sabemos que
$\bar{X} \sim N(\mu, \sigma^2/n)$, luego el
estad\'{i}stico
\[
Z = \frac{\bar{X} - \mu}{\sigma/\sqrt{n}}
\sim N(0,1).
\]
Este estad\'{i}stico se denomina
\textbf{pivote}: es una funci\'{o}n de los
datos y del par\'{a}metro $\mu$ cuya
distribuci\'{o}n es completamente conocida
e independiente de $\mu$. Eso nos permite
construir el intervalo imponiendo una
condici\'{o}n de probabilidad:
\[
P\!\left(-z_{1-\alpha/2} \leq
\frac{\bar{X} - \mu}{\sigma/\sqrt{n}}
\leq z_{1-\alpha/2}\right) = 1 - \alpha.
\]
Despejando $\mu$:
\[
P\!\left(\bar{X} - z_{1-\alpha/2}
\frac{\sigma}{\sqrt{n}} \leq \mu \leq
\bar{X} + z_{1-\alpha/2}
\frac{\sigma}{\sqrt{n}}\right)
= 1 - \alpha.
\]

\begin{definicion}{Intervalo de confianza
		para $\mu$ con varianza conocida}
	{ic-varianza-conocida}
	Sea $\bar{x}$ la media de una muestra
	aleatoria de tama\~{n}o $n$ de una
	poblaci\'{o}n $N(\mu, \sigma^2)$ con
	$\sigma^2$ conocida. Un intervalo de
	confianza del $100(1-\alpha)\%$ para
	$\mu$ es:
	\[
	\bar{x} - z_{1-\alpha/2}
	\frac{\sigma}{\sqrt{n}}
	\leq \mu \leq
	\bar{x} + z_{1-\alpha/2}
	\frac{\sigma}{\sqrt{n}}.
	\]
	Se define:
	\begin{itemize}[leftmargin=2em]
		\item $l = \bar{x} - z_{1-\alpha/2}
		\dfrac{\sigma}{\sqrt{n}}$:
		\textbf{l\'{i}mite inferior} de
		confianza.
		\item $u = \bar{x} + z_{1-\alpha/2}
		\dfrac{\sigma}{\sqrt{n}}$:
		\textbf{l\'{i}mite superior} de
		confianza.
		\item $u - l = 2z_{1-\alpha/2}
		\dfrac{\sigma}{\sqrt{n}}$:
		\textbf{longitud} del intervalo.
		\item $E = z_{1-\alpha/2}
		\dfrac{\sigma}{\sqrt{n}}$:
		\textbf{error de estimaci\'{o}n}
		o precisi\'{o}n del estimador.
	\end{itemize}
\end{definicion}

\begin{observacion}
	Existe una tensi\'{o}n entre confianza
	e informatividad: a mayor $1-\alpha$,
	el intervalo es m\'{a}s ancho y por
	lo tanto menos informativo sobre
	d\'{o}nde est\'{a} $\mu$. Un intervalo
	del $99\%$ de confianza siempre
	ser\'{a} m\'{a}s ancho que uno del
	$95\%$ con los mismos datos. La
	\'{u}nica forma de aumentar la
	confianza sin sacrificar informatividad
	es aumentar $n$.
\end{observacion}

\begin{ejemplo}{Intervalo de confianza para
		el gasto en sitios web}{negocios-web}
	
	En una investigaci\'{o}n sobre negocios
	peque\~{n}os con sitio web se encontr\'{o}
	que el gasto medio anual es de
	\$11.500. Se dispone de una muestra de
	$n = 60$ negocios con desviaci\'{o}n
	est\'{a}ndar conocida $\sigma = \$4.000$.
	Construya un intervalo de confianza del
	$95\%$ para la media del gasto anual.
	
	\medskip
	Con $\alpha = 0{,}05$ el valor cr\'{i}tico
	es $z_{0{,}975} = 1{,}960$. El error de
	estimaci\'{o}n es:
	\[
	E = 1{,}960 \cdot
	\frac{4.000}{\sqrt{60}}
	= 1{,}960 \cdot 516{,}4
	= 1.012{,}1.
	\]
	El intervalo de confianza del $95\%$ es:
	\[
	11.500 - 1.012{,}1
	\leq \mu \leq
	11.500 + 1.012{,}1,
	\]
	\[
	10.487{,}9 \leq \mu \leq 12.512{,}1.
	\]
	Con un $95\%$ de confianza, el gasto medio
	anual de los negocios peque\~{n}os con
	sitio web se encuentra entre \$10.487{,}9
	y \$12.512{,}1.
	
\end{ejemplo}

% ----------------------------------------------------------
\subsection{Tama\~{n}o muestral \'{o}ptimo}

En la pr\'{a}ctica, el investigador suele
fijar de antemano el error de estimaci\'{o}n
m\'{a}ximo tolerable $E$ y el coeficiente
de confianza $1-\alpha$, y desea saber
cu\'{a}ntas observaciones necesita. Despejando
$n$ de la expresi\'{o}n del error:
\[
E = z_{1-\alpha/2}
\frac{\sigma}{\sqrt{n}}
\implies \sqrt{n} =
z_{1-\alpha/2} \frac{\sigma}{E}
\implies n =
\left(\frac{z_{1-\alpha/2}\,
	\sigma}{E}\right)^2.
\]

\begin{definicion}{Tama\~{n}o muestral
		\'{o}ptimo}{tamano-muestral}
	Para estimar $\mu$ con un error de
	estimaci\'{o}n m\'{a}ximo $E$ y un
	coeficiente de confianza $1-\alpha$,
	el tama\~{n}o muestral requerido es:
	\[
	n = \left(\frac{z_{1-\alpha/2}\,
		\sigma}{E}\right)^2.
	\]
	Si $n$ no es entero, se redondea
	\textbf{hacia arriba} para garantizar
	que el error no supere $E$.
\end{definicion}

\begin{ejemplo}{Tama\~{n}o muestral para
		estimar salarios}{consejos}
	
	Un estudiante desea estimar el salario
	medio mensual de los miembros de consejos
	ciudadanos de grandes ciudades, con un
	error inferior a $100$ u.m.\ y un nivel
	de confianza del $95\%$. Un informe del
	Departamento del Trabajo reporta
	$\sigma = 1.000$ u.m. ¿Cu\'{a}ntas
	observaciones se necesitan?
	
	\medskip
	Con $\alpha = 0{,}05$, $z_{0{,}975} =
	1{,}960$, $E = 100$ y $\sigma = 1.000$:
	\[
	n = \left(\frac{1{,}960 \times 1.000}
	{100}\right)^2
	= (19{,}6)^2
	= 384{,}16.
	\]
	Redondeando hacia arriba: $n = 385$.
	Se requieren al menos $385$ observaciones
	para garantizar un error inferior a
	$100$ u.m.\ con un $95\%$ de confianza.
	
\end{ejemplo}

% ----------------------------------------------------------
\subsection{Intervalo de confianza con
	varianza desconocida}

Cuando $\sigma^2$ es desconocida, la
reemplazamos por $S$ --- exactamente como
en la prueba $t$. El pivote ya no es
$Z \sim N(0,1)$ sino:
\[
T = \frac{\bar{X} - \mu}{S/\sqrt{n}}
\sim t_{n-1},
\]
por la Proposici\'{o}n~\ref{prop:dist-t0}.
Imponiendo la condici\'{o}n de probabilidad:
\[
P\!\left(-t_{n-1,\,\alpha/2} \leq
\frac{\bar{X} - \mu}{S/\sqrt{n}}
\leq t_{n-1,\,\alpha/2}\right)
= 1 - \alpha,
\]
y despejando $\mu$:
\[
P\!\left(\bar{X} - t_{n-1,\,\alpha/2}
\frac{S}{\sqrt{n}} \leq \mu \leq
\bar{X} + t_{n-1,\,\alpha/2}
\frac{S}{\sqrt{n}}\right)
= 1 - \alpha.
\]

\begin{definicion}{Intervalo de confianza
		para $\mu$ con varianza desconocida}
	{ic-varianza-desconocida}
	Sea $\bar{x}$ y $s$ la media y
	desviaci\'{o}n est\'{a}ndar muestral
	de una muestra aleatoria de tama\~{n}o
	$n$ de una poblaci\'{o}n normal con
	$\sigma^2$ desconocida. Un intervalo
	de confianza del $100(1-\alpha)\%$
	para $\mu$ es:
	\[
	\bar{x} - t_{n-1,\,\alpha/2}
	\frac{s}{\sqrt{n}}
	\leq \mu \leq
	\bar{x} + t_{n-1,\,\alpha/2}
	\frac{s}{\sqrt{n}}.
	\]
\end{definicion}

\begin{ejemplo}{Intervalo de confianza para
		tweets diarios}{tweets}
	
	El n\'{u}mero de tweets diarios de j\'{o}venes
	entre 16 y 24 a\~{n}os sigue una distribuci\'{o}n
	normal con media y varianza desconocidas.
	Se toma una muestra de $n = 21$ usuarios,
	obteniendo $\bar{x} = 12$ tweets y
	$s = 6$ tweets. Construya un intervalo
	de confianza del $90\%$ para la media
	poblacional.
	
	\medskip
	Con $\alpha = 0{,}10$, los grados de
	libertad son $n - 1 = 20$, y el valor
	cr\'{i}tico es $t_{20,\,0{,}05} = 1{,}725$.
	El error de estimaci\'{o}n es:
	\[
	E = 1{,}725 \cdot
	\frac{6}{\sqrt{21}}
	= 1{,}725 \cdot 1{,}309
	= 2{,}258.
	\]
	El intervalo de confianza del $90\%$ es:
	\[
	12 - 2{,}258
	\leq \mu \leq
	12 + 2{,}258,
	\]
	\[
	9{,}742 \leq \mu \leq 14{,}258.
	\]
	Con un $90\%$ de confianza, el n\'{u}mero
	medio de tweets diarios se encuentra entre
	$9{,}742$ y $14{,}258$.
	
\end{ejemplo}

\end{document}
